\documentclass[11pt,a4paper,oneside,openany]{book}

\usepackage[T1]{fontenc}
\usepackage[utf8]{inputenc}
\usepackage{libertinus}
\usepackage{microtype}
\usepackage[a4paper,margin=27mm,headheight=15pt]{geometry}
\usepackage{amsmath,amssymb,amsthm,mathtools}
\usepackage{bm}
\usepackage{mathrsfs}
\usepackage{booktabs,longtable,array,tabularx}
\usepackage{enumitem}
\usepackage{xcolor}
\usepackage{tikz}
\usetikzlibrary{arrows.meta,positioning,calc}
\usepackage[most]{tcolorbox}
\usepackage{fancyhdr}
\usepackage{xurl}
\usepackage{hyperref}
\usepackage[nameinlink,noabbrev,capitalise]{cleveref}
% Canonical mathematical notation for active FabiusFunction LaTeX documents.
%
% This file is intentionally a plain TeX input rather than a LaTeX package.
% Every standalone document loads it by an explicit path relative to that
% document's directory; included fragments inherit the definitions from their
% root document.  Keep this file independent of document layout and metadata.
%
% Contract:
%   * one command has one mathematical meaning throughout the active corpus;
%   * names encode semantic roles rather than a document-local choice of letter;
%   * visually similar but differently normalized objects have different print;
%   * every command below is defined normatively in the notation catalogue.
%
% The active corpus excludes docs/archive/.  Historical sources may retain
% their original notation only inside an explicitly labelled quotation.

\makeatletter
\@ifundefined{mathbb}{%
  \PackageError{fabius-notation}{The amssymb package must be loaded first}%
    {Load amsmath and amssymb before inputting fabius-notation.tex.}%
}{}
\@ifundefined{operatorname}{%
  \PackageError{fabius-notation}{The amsmath package must be loaded first}%
    {Load amsmath before inputting fabius-notation.tex.}%
}{}
\makeatother

% Version stamp used by the catalogue and by static audits.
\newcommand{\FabiusNotationVersion}{2026-08-31}

% Number systems.  NaturalNumbers is explicitly the nonnegative convention.
\newcommand{\NaturalNumbers}{\mathbb{N}_{0}}
\newcommand{\PositiveIntegers}{\mathbb{N}_{>0}}
\newcommand{\IntegerNumbers}{\mathbb{Z}}
\newcommand{\RationalNumbers}{\mathbb{Q}}
\newcommand{\RealNumbers}{\mathbb{R}}
\newcommand{\ComplexNumbers}{\mathbb{C}}
\newcommand{\ExtendedRealNumbers}{\overline{\mathbb{R}}}

% The three principal Fabius--Rvachev functions.  Subscripts are deliberate:
% a bare F and a calligraphic F have several unrelated uses in the corpus.
\newcommand{\FabiusBounded}{F_{\mathrm{bd}}}
\newcommand{\FabiusGlobal}{F_{\mathrm{gl}}}
\newcommand{\FabiusQuantile}{Q_{F}}
\newcommand{\FabiusClampedQuantile}{Q_{F,\mathrm{cl}}}
\newcommand{\RvachevUp}{\operatorname{up}}

% Constants, elementary operators, and delimiters.
\newcommand{\EulerE}{\mathrm{e}}
\newcommand{\ImaginaryUnit}{\mathrm{i}}
\newcommand{\Differential}{\mathop{}\!\mathrm{d}}
\newcommand{\AbsoluteValue}[1]{\left\lvert #1\right\rvert}
\newcommand{\Norm}[1]{\left\lVert #1\right\rVert}
\newcommand{\Floor}[1]{\left\lfloor #1\right\rfloor}
\newcommand{\Ceiling}[1]{\left\lceil #1\right\rceil}
\newcommand{\FractionalPart}[1]{\left\{#1\right\}}
\newcommand{\PositivePart}[1]{\left(#1\right)_{+}}
\newcommand{\IndicatorOf}[1]{\mathbf{1}_{#1}}
\newcommand{\SetOf}[1]{\left\{#1\right\}}
\newcommand{\InnerProduct}[1]{\left\langle #1\right\rangle}
\newcommand{\CardinalityOf}[1]{\left\lvert #1\right\rvert}
\newcommand{\DefinitionEquals}{\mathrel{\coloneqq}}
\newcommand{\Sign}{\operatorname{sgn}}
\newcommand{\IdentityOperator}{\operatorname{id}}
\newcommand{\SpanOperator}{\operatorname{span}}
\newcommand{\TraceOperator}{\operatorname{tr}}
\newcommand{\RankOperator}{\operatorname{rank}}
\newcommand{\DiagonalOperator}{\operatorname{diag}}
\newcommand{\ResidueOperator}{\operatorname*{Res}}
\newcommand{\LeastCommonMultiple}{\operatorname{lcm}}
\newcommand{\OrderOperator}{\operatorname{ord}}
\newcommand{\PrincipalLogarithm}{\operatorname{Log}}
\newcommand{\FallingFactorial}[2]{\left(#1\right)^{\underline{#2}}}
\newcommand{\RisingFactorial}[2]{\left(#1\right)^{\overline{#2}}}

% Support, topology, convolution, and metric notation.
\newcommand{\SupportOperator}{\operatorname{supp}}
\newcommand{\SupportOf}[1]{\SupportOperator\!\left(#1\right)}
\newcommand{\TopologicalSupportOf}[1]{\operatorname{tsupp}\!\left(#1\right)}
\newcommand{\Convolution}{\mathbin{*}}
\newcommand{\MetricDistance}{\operatorname{dist}}
\newcommand{\TotalVariationDistance}{d_{\mathrm{TV}}}

% Probability.  Equality and convergence in law are not metric distance.
\newcommand{\Probability}{\mathbb{P}}
\newcommand{\Expectation}{\mathbb{E}}
\newcommand{\Variance}{\operatorname{Var}}
\newcommand{\Covariance}{\operatorname{Cov}}
\newcommand{\LawOperator}{\operatorname{Law}}
\newcommand{\LawOf}[1]{\LawOperator\!\left(#1\right)}
\newcommand{\EqualInLaw}{\mathrel{\overset{\mathrm{law}}{=}}}
\newcommand{\ConvergesInLaw}{\mathrel{\xrightarrow{\mathrm{law}}}}

% Fourier and integral transforms.  The printed labels expose normalization.
% FourierTwoPi uses exp(-2*pi*i*x*xi); FourierAngular uses exp(-i*x*omega).
\newcommand{\FourierTwoPi}[1]{\widehat{#1}_{\,2\pi}}
\newcommand{\FourierAngular}[1]{\widehat{#1}_{\,\mathrm{ang}}}
\newcommand{\LaplaceTransformSymbol}{\mathcal{L}}
\newcommand{\MellinTransformSymbol}{\mathcal{M}}
\newcommand{\LaplaceTransformOf}[1]{\LaplaceTransformSymbol\!\left[#1\right]}
\newcommand{\MellinTransformOf}[1]{\MellinTransformSymbol\!\left[#1\right]}
\newcommand{\MomentGeneratingFunction}[1]{M_{#1}}
\newcommand{\CharacteristicFunction}[1]{\varphi_{#1}}

% The two standard sinc functions must never share one symbol.
\newcommand{\SincRad}{\operatorname{sinc}_{\mathrm{rad}}}
\newcommand{\SincPi}{\operatorname{sinc}_{\pi}}

% Binary arithmetic and Thue--Morse notation.
\newcommand{\BinaryDigitSum}{\operatorname{wt}_{2}}
\newcommand{\ThueMorseSignSymbol}{\varepsilon}
\newcommand{\ThueMorseSign}[1]{\varepsilon_{#1}}
\newcommand{\PadicValuation}[1]{\nu_{#1}}
\newcommand{\TwoAdicValuation}{\nu_{2}}

% q-series notation.  Every base is an explicit argument.
\newcommand{\QPochhammer}[3]{\left(#1;#2\right)_{#3}}
\newcommand{\GaussianBinomial}[3]{%
  \genfrac{[}{]}{0pt}{}{#1}{#2}_{#3}%
}
\newcommand{\QInteger}[2]{\left[#1\right]_{#2}}

% Named special functions.
\newcommand{\Polylogarithm}{\operatorname{Li}}
\newcommand{\LambertW}{W}
\newcommand{\GammaFunction}{\Gamma}
\newcommand{\RiemannZeta}{\zeta}

% Asymptotics.  BigO and LittleO are operator tokens for existing formula
% grammar; prefer the At forms so that the limiting regime is printed.
\newcommand{\BigO}{\mathcal{O}}
\newcommand{\LittleO}{o}
\newcommand{\BigOAt}[2]{\mathcal{O}_{#2}\!\left(#1\right)}
\newcommand{\LittleOAt}[2]{o_{#2}\!\left(#1\right)}

\usepackage{makeidx}
\makeindex

\definecolor{deepolive}{RGB}{67,79,45}
\definecolor{softolive}{RGB}{236,239,226}
\definecolor{darkteal}{RGB}{31,83,84}
\definecolor{warmgray}{RGB}{92,87,78}
\definecolor{softcoral}{RGB}{246,225,219}
\definecolor{linkblue}{RGB}{31,76,122}

\hypersetup{
  colorlinks=true,
  linkcolor=deepolive,
  citecolor=darkteal,
  urlcolor=linkblue,
  pdftitle={Combinatorial Formulae and Inversion Theorems},
  pdfauthor={A proof-oriented synthesis of five draft manuscripts},
  pdfsubject={Stirling, Bell, Eulerian, Bernoulli, Sheffer, Faa di Bruno, Lagrange-Good inversion, and coefficient algorithms}
}

\setlist{itemsep=2pt,topsep=4pt,parsep=1pt}
\allowdisplaybreaks[3]
\raggedbottom
\setcounter{tocdepth}{2}
\setcounter{secnumdepth}{3}

\pagestyle{fancy}
\fancyhf{}
\fancyhead[L]{\small\nouppercase{\leftmark}}
\fancyhead[R]{\small\thepage}
\renewcommand{\headrulewidth}{0.4pt}
\fancypagestyle{plain}{%
  \fancyhf{}%
  \fancyfoot[C]{\thepage}%
  \renewcommand{\headrulewidth}{0pt}%
}

\newtheorem{theorem}{Theorem}[chapter]
\newtheorem{lemma}[theorem]{Lemma}
\newtheorem{proposition}[theorem]{Proposition}
\newtheorem{corollary}[theorem]{Corollary}
\newtheorem{identity}[theorem]{Identity}
\theoremstyle{definition}
\newtheorem{definition}[theorem]{Definition}
\newtheorem{example}[theorem]{Example}
\newtheorem{algorithm}[theorem]{Algorithm}
\theoremstyle{remark}
\newtheorem{remark}[theorem]{Remark}
\newtheorem{historical}[theorem]{Historical note}

\newtcolorbox{masterbox}[1][]{
  enhanced,breakable,colback=softolive,colframe=deepolive,
  boxrule=0.7pt,arc=1mm,left=2mm,right=2mm,top=1.5mm,bottom=1.5mm,
  title={Master principle},fonttitle=\bfseries,#1
}
\newtcolorbox{auditbox}[1][]{
  enhanced,breakable,colback=softcoral,colframe=warmgray,
  boxrule=0.6pt,arc=1mm,left=2mm,right=2mm,top=1.5mm,bottom=1.5mm,
  title={Editorial normalization},fonttitle=\bfseries,#1
}


\title{\Huge\bfseries Combinatorial Formulae and Inversion Theorems\\[4mm]
\Large A Unified Coefficient Calculus for Stirling, Bell, Eulerian,\\
Bernoulli, Sheffer, Fa\`a di Bruno, and Lagrange--Good Structures}
\author{A proof-oriented synthesis of five draft manuscripts}
\date{September 2026}

\begin{document}
\frontmatter
\hypersetup{pageanchor=false}
\maketitle
\hypersetup{pageanchor=true}

\chapter*{Abstract}
\addcontentsline{toc}{chapter}{Abstract}
This monograph merges five overlapping draft manuscripts into one coherent,
deduplicated, proof-oriented account of combinatorial coefficient calculus and
series inversion.  Its central objects are signed and unsigned Stirling numbers,
Lah numbers, Bell numbers and polynomials, Eulerian numbers, Bernoulli and Euler
polynomials, Sheffer sequences, exponential Riordan arrays, Fa\`a di Bruno
formulae, and one- and multivariable Lagrange inversion.  Supporting Wikipedia
and MathWorld dossiers have been incorporated selectively, and additional
standard references were checked for normalization, hypotheses, and provenance.

The treatment begins with formal power series, residues, factorial bases, finite
differences, and the exponential formula.  It then develops recurrences,
generating functions, transforms, congruences, special-function sums, saddle-point
asymptotics, restricted partitions, descent calculus, composition, and reversion.
Newly unified chapters add Pochhammer and binomial-transform calculus;
Bernoulli--Euler--Genocchi and Bernoulli--Cauchy families; complementary Bell
numbers; Sheffer and Riordan theory; higher product and quotient rules; the
Fr\'echet form of Fa\`a di Bruno; Good's determinant formula for multivariate
inversion; Narayana refinements and the complete associahedral face vector;
gamma, polygamma, and Barnes \(G\); Euler--Maclaurin and Darboux methods;
partition-lattice M\"obius inversion; Kepler reversion; and precision-doubling
coefficient algorithms.

Definitions are connected by explicit changes of basis, labelled constructions,
residue substitutions, and incidence-algebra inversion.  Repeated material from
the drafts has been consolidated rather than stacked, while distinct proofs and
applications have been retained where they expose different mechanisms.  Every
new theorem introduced in the merger is accompanied by a proof, and statements
that require analytic rather than merely formal hypotheses say so explicitly.
A final concordance records the contribution of each draft and of the supporting
topic files.

\chapter*{How to read this work}
\addcontentsline{toc}{chapter}{How to read this work}
The chapters are arranged by mathematical dependence, not by the order of the source
files.  The two master ideas are:
\begin{enumerate}
  \item a coefficient is often easiest to compute after changing to the basis in
        which the relevant operator is diagonal; and
  \item an exponential generating function remembers the decomposition of a
        labelled object into components.
\end{enumerate}
Stirling numbers implement the first idea.  Bell polynomials implement the second.
Eulerian numbers mediate between powers and shifted binomial coefficients.  Fa\`a di
Bruno's formula is the differential form of labelled composition, while Lagrange
inversion is the coefficient form of implicit substitution.

A result stated for formal power series requires no convergence.  When an analytic
version is asserted, its domain hypotheses are written explicitly.  Statements in
the archive that report historical computations---for example, lists of Bell
primes or verified exact modular periods---are distinguished from theorems and are
accompanied by a certificate procedure rather than misrepresented as symbolic
identities.

\begin{auditbox}
Three systematic repairs are used throughout.
\begin{enumerate}
\item Bell numbers are denoted $\BellNumber n$, complete exponential Bell polynomials by
      $\ExponentialCompleteBellPolynomial n$, partial Bell polynomials by $\ExponentialPartialBellPolynomial nk$, and Bernoulli numbers by
      $\beta_n$.
\item Type-$B$ Eulerian numbers are indexed by descents in positions
      $0,1,\ldots,n-1$ with $\pi_0=0$.  This removes the one-step shift in one formula
      in the archive.
\item The archived Moser--Wyman display contains unnamed quantities
      $B,P_i,Q_i$.  We give the exact saddle-point scheme that defines every
      correction coefficient recursively and state the leading expansion in closed
      form; no undefined placeholder is retained.
\end{enumerate}
\end{auditbox}

\tableofcontents

\chapter*{Notation}
\addcontentsline{toc}{chapter}{Notation}
For $n\in\NaturalNumbers$, write
\[
 \FallingFactorial{x}{n}=x(x-1)\cdots(x-n+1),\qquad
 \RisingFactorial{x}{n}=x(x+1)\cdots(x+n-1),
\]
with both expressions equal to $1$ when $n=0$.  Coefficient extraction is denoted by
$\CoefficientExtraction z n F(z)$, and $\ResidueOperator_z F(z)\Differential z$ means the coefficient of $z^{-1}$ in a
formal Laurent series.  Our principal arrays are
\[
 \SignedStirlingFirstKind{n}{k},\qquad \UnsignedStirlingFirstKind nk=(-1)^{n-k}\SignedStirlingFirstKind{n}{k},\qquad
 \StirlingSecondKind nk,\qquad \TypeAEulerianNumber nk,
\]
where $\SignedStirlingFirstKind{n}{k}$ is the signed first-kind Stirling number, $\UnsignedStirlingFirstKind nk$ its unsigned
version, $\StirlingSecondKind nk$ the second-kind Stirling number, and $\TypeAEulerianNumber nk$ the
ordinary Eulerian number.  Bell numbers are $\BellNumber n$; exponential partial and
complete Bell polynomials are $\ExponentialPartialBellPolynomial nk$ and $\ExponentialCompleteBellPolynomial n$.

Unless stated otherwise, sums over an impossible index range are zero, and
$0^0=1$ in coefficient formulae.  Bernoulli numbers use
\[
 \frac{t}{\EulerE^t-1}=\sum_{n\ge0}\beta_n\frac{t^n}{n!},
 \qquad \beta_1=-\frac12.
\]

\mainmatter

\part{Coefficient Calculus and Labelled Structures}

\chapter{Formal power series and coefficient extraction}
\label{ch:formal}

\section{Formal series as an algebraic environment}
Let $R$ be a commutative $\RationalNumbers$-algebra.  The ring $R[[z]]$ consists of expressions
$F(z)=\sum_{n\ge0}f_nz^n$ with termwise addition and Cauchy multiplication.  A series
is invertible precisely when $f_0$ is invertible in $R$.  If $G(0)=0$, then the
composition $F(G(z))$ is well-defined because each coefficient receives only
finitely many contributions.

\begin{lemma}[Coefficient rules]\label{lem:coeff-rules}
For $F(z)=\sum f_nz^n$ and $G(z)=\sum g_nz^n$,
\begin{align}
 \CoefficientExtraction z n F(z)G(z)&=\sum_{j=0}^n f_jg_{n-j},\label{eq:cauchy}\\
 \CoefficientExtraction z n F'(z)&=(n+1)f_{n+1},\label{eq:der-coeff}\\
 \CoefficientExtraction z n \frac{F(z)}{1-az}&=\sum_{j=0}^n a^{n-j}f_j.\label{eq:geom-conv}
\end{align}
If $F$ is analytic on $|z|\le r$, then also
\begin{equation}\label{eq:cauchy-coeff}
 f_n=\frac{1}{2\pi \ImaginaryUnit}\int_{|z|=r}\frac{F(z)}{z^{n+1}}\Differential z,
 \qquad |f_n|\le r^{-n}\max_{|z|=r}|F(z)|.
\end{equation}
\end{lemma}
\begin{proof}
The first identity is the definition of Cauchy multiplication.  The second follows
by differentiating term by term, and the third by multiplying with
$\sum_{m\ge0}a^mz^m$.  The analytic formula is Cauchy's coefficient formula, and the
bound follows by estimating the contour integral by its length $2\pi r$.
\end{proof}

\section{Formal residues}
For a Laurent series $A(z)$, set
$\ResidueOperator_z A(z)\Differential z=\CoefficientExtraction z{-1}A(z)$.  Then
\begin{equation}\label{eq:res-coeff}
 \CoefficientExtraction z n F(z)=\ResidueOperator_z\frac{F(z)}{z^{n+1}}\Differential z,
 \qquad \ResidueOperator_z A'(z)\Differential z=0.
\end{equation}
The second identity is simply the fact that the derivative of $z^m$ has no
$z^{-1}$ term.

\begin{theorem}[Formal change of variables for residues]\label{thm:res-subst}
Let $u=u(z)\in zR[[z]]$ have invertible linear coefficient.  For every Laurent series
$A$ for which the composition is defined,
\begin{equation}\label{eq:formal-res-subst}
 \ResidueOperator_z A(u(z))u'(z)\Differential z=\ResidueOperator_u A(u)\Differential u.
\end{equation}
\end{theorem}
\begin{proof}
By linearity it suffices to take $A(u)=u^m$.  If $m\ne-1$, then
$u(z)^mu'(z)=\frac{1}{m+1}(u(z)^{m+1})'$ has residue zero.  For $m=-1$,
$u'(z)/u(z)=z^{-1}+H(z)$ with $H\in R[[z]]$, because
$u(z)=az(1+O(z))$.  Thus both sides have residue $1$.
\end{proof}

\begin{masterbox}
The combination of \cref{eq:res-coeff,eq:formal-res-subst} turns implicit
substitution into coefficient extraction.  This is the entire algebraic engine of
Lagrange inversion; no analytic contour is needed until one asks about convergence.
\end{masterbox}

\section{Ordinary and exponential generating functions}
For a sequence $(a_n)$, its ordinary and exponential generating functions are
\[
 A(z)=\sum_{n\ge0}a_nz^n,
 \qquad \ExponentialGeneratingFunctionOf{A}(z)=\sum_{n\ge0}a_n\frac{z^n}{n!}.
\]
Ordinary multiplication gives unweighted convolution.  Exponential multiplication
gives binomial convolution:
\begin{equation}\label{eq:egf-product}
 \ExponentialGeneratingFunctionOf{A}(z)\ExponentialGeneratingFunctionOf{B}(z)
 =\sum_{n\ge0}\left(\sum_{j=0}^n\binom nj a_jb_{n-j}\right)\frac{z^n}{n!}.
\end{equation}
This binomial factor is the number of ways to choose which labels belong to the
first component.

\section{Finite differences and factorial bases}
Let $E$ be the shift operator $(Ef)(x)=f(x+1)$ and
$\Delta=E-1$.  Then
\begin{equation}\label{eq:finite-diff-falling}
 \Delta\FallingFactorial{x}{k}=k\FallingFactorial{x}{k-1},
 \qquad
 \sum_{i=0}^{n-1}\FallingFactorial{i}{k}=\frac{\FallingFactorial{n}{k+1}}{k+1}.
\end{equation}
\begin{proof}
The first follows from
$\FallingFactorial{x+1}{k}-\FallingFactorial{x}{k}=k\FallingFactorial{x}{k-1}$ after factoring the common
product.  Summing this telescoping identity gives the second formula.
\end{proof}

\begin{theorem}[Newton expansion]\label{thm:newton-expansion}
Every polynomial $p$ of degree at most $d$ has the unique expansion
\begin{equation}\label{eq:newton-expansion}
 p(x)=\sum_{k=0}^d\frac{\Delta^kp(0)}{k!}\FallingFactorial{x}{k}
     =\sum_{k=0}^d\Delta^kp(0)\binom{x}{k}.
\end{equation}
\end{theorem}
\begin{proof}
The factorial polynomials form a basis because their degrees are distinct and their
leading coefficients are one.  If $p=\sum a_k\FallingFactorial{x}{k}$, applying $\Delta^j$
and setting $x=0$ kills terms with $k>j$ by a remaining factor $x$, kills terms
with $k<j$ by excessive differencing, and gives $j!a_j$ for $k=j$.
\end{proof}

\section{Labelled products and the exponential formula}
A labelled combinatorial class is one whose atoms carry distinct labels.  Suppose a
connected component of size $j$ has total weight $x_j$.  A set of $k$ such components
with total size $n$ has weight enumerator $\ExponentialPartialBellPolynomial nk(x_1,x_2,\ldots)$; a set of any
number of components has enumerator $\ExponentialCompleteBellPolynomial n(x_1,\ldots,x_n)$.

\begin{theorem}[The exponential formula]\label{thm:exponential-formula}
Let
\[
 C(z)=\sum_{j\ge1}x_j\frac{z^j}{j!}.
\]
Then
\begin{align}
 \frac{C(z)^k}{k!}
  &=\sum_{n\ge k}\ExponentialPartialBellPolynomial nk(x_1,\ldots,x_{n-k+1})\frac{z^n}{n!},
  \label{eq:partial-bell-egf}\\
 \EulerE^{uC(z)}
  &=\sum_{n\ge k\ge0}\ExponentialPartialBellPolynomial nk(x_1,\ldots)u^k\frac{z^n}{n!},
  \label{eq:bivariate-bell-egf}\\
 \EulerE^{C(z)}
  &=\sum_{n\ge0}\ExponentialCompleteBellPolynomial n(x_1,\ldots,x_n)\frac{z^n}{n!}.
  \label{eq:complete-bell-egf}
\end{align}
\end{theorem}
\begin{proof}
Expand $C(z)^k/k!$ multinomially.  If $j_i$ components have size $i$, then
$\sum_i j_i=k$ and $\sum_iij_i=n$, and the coefficient of $z^n$ is
\[
 \prod_i\frac1{j_i!}\left(\frac{x_i}{i!}\right)^{j_i}.
\]
Multiplication by $n!$ labels the atoms.  This proves the first formula.  Summing it
with weight $u^k$ proves the second, and setting $u=1$ proves the third.
\end{proof}

\begin{corollary}[Partition-type count]\label{cor:partition-type}
The number of partitions of an $n$-element labelled set with exactly $j_i$ blocks of
size $i$ is
\begin{equation}\label{eq:partition-type-count}
 \frac{n!}{\prod_{i\ge1}(i!)^{j_i}j_i!},
 \qquad \sum_iij_i=n.
\end{equation}
\end{corollary}
\begin{proof}
Order the $n$ labels, cut the ordering into blocks of prescribed sizes, then divide
by $i!$ for the internal order of each size-$i$ block and by $j_i!$ for the order of
blocks of equal size.
\end{proof}

\part{Stirling Numbers, Lah Numbers, and Basis Changes}

\chapter{The two Stirling triangles}
\label{ch:stirling-def}

\section{First-kind numbers: cycles and factorial coefficients}
\begin{definition}
The signed Stirling numbers of the first kind $\SignedStirlingFirstKind{n}{k}$ and the unsigned numbers
$\UnsignedStirlingFirstKind nk$ are defined by
\begin{align}
 \FallingFactorial{x}{n}&=\sum_{k=0}^n \SignedStirlingFirstKind{n}{k}x^k,\label{eq:signed-first-def}\\
 \RisingFactorial{x}{n}&=\sum_{k=0}^n\UnsignedStirlingFirstKind nkx^k,\label{eq:unsigned-first-def}\\
 \SignedStirlingFirstKind{n}{k}&=(-1)^{n-k}\UnsignedStirlingFirstKind nk.\label{eq:first-sign}
\end{align}
\end{definition}
The sign relation follows from
$\RisingFactorial{x}{n}=(-1)^n\FallingFactorial{-x}{n}$.

\begin{theorem}[Permutation interpretation]\label{thm:first-cycle}
$\UnsignedStirlingFirstKind nk$ is the number of permutations of $[n]$ having exactly $k$ cycles.
Consequently,
\begin{equation}\label{eq:first-recurrence}
 \UnsignedStirlingFirstKind{n+1}{k}=n\UnsignedStirlingFirstKind nk+\UnsignedStirlingFirstKind n{k-1},
\end{equation}
with $\UnsignedStirlingFirstKind00=1$ and zero boundary values outside $0\le k\le n$.
\end{theorem}
\begin{proof}
Given a permutation of $[n+1]$ with $k$ cycles, either $n+1$ is a singleton cycle,
leaving a permutation of $[n]$ with $k-1$ cycles, or it is inserted immediately
after one of the $n$ existing elements in the cyclic notation of a permutation of
$[n]$ with $k$ cycles.  These alternatives are disjoint and reversible.
The same recurrence follows algebraically by multiplying
$\RisingFactorial{x}{n}$ by $x+n$ and comparing coefficients.
\end{proof}

For $n=3$, the cycle counts are $\UnsignedStirlingFirstKind33=1$, $\UnsignedStirlingFirstKind32=3$, and
$\UnsignedStirlingFirstKind31=2$, corresponding to the identity, the three transpositions, and the two
$3$-cycles.  Summing over the possible number of cycles gives
\begin{equation}\label{eq:first-row-sum}
 \sum_{k=0}^n\UnsignedStirlingFirstKind nk=n!,
\end{equation}
which also follows from \cref{eq:unsigned-first-def} at $x=1$.

\section{Second-kind numbers: set partitions and surjections}
\begin{definition}
The Stirling number of the second kind $\StirlingSecondKind nk$ is the number of partitions of
$[n]$ into exactly $k$ nonempty blocks.
\end{definition}

\begin{theorem}[Second-kind recurrence]\label{thm:second-recurrence}
For $n,k\ge0$,
\begin{equation}\label{eq:second-recurrence}
 \StirlingSecondKind{n+1}{k}=k\StirlingSecondKind nk+\StirlingSecondKind n{k-1},
\end{equation}
with $\StirlingSecondKind00=1$ and zero boundary values outside $0\le k\le n$.
\end{theorem}
\begin{proof}
In a partition of $[n+1]$ into $k$ blocks, the new label either joins one of the $k$
blocks of a partition of $[n]$, or forms a singleton block beside a partition of
$[n]$ into $k-1$ blocks.
\end{proof}

\begin{theorem}[Surjection formula]\label{thm:second-explicit}
For all $n,k\ge0$,
\begin{equation}\label{eq:second-explicit}
 \StirlingSecondKind nk=\frac1{k!}\sum_{j=0}^k(-1)^{k-j}\binom kjj^n
 =\sum_{j=0}^k\frac{(-1)^{k-j}j^n}{(k-j)!j!}.
\end{equation}
\end{theorem}
\begin{proof}
There are $k!\StirlingSecondKind nk$ surjections from $[n]$ onto a fixed $k$-element codomain:
a set partition into fibres followed by an ordering of the fibres.  Inclusion--
exclusion over the set of omitted codomain values gives
$\sum_{j=0}^k(-1)^{k-j}\binom kjj^n$.  Division by $k!$ proves the formula.
\end{proof}

The simple diagonals follow at once:
\begin{equation}\label{eq:second-simple}
 \StirlingSecondKind nn=1,\qquad \StirlingSecondKind n1=1\ (n\ge1),\qquad
 \StirlingSecondKind n{n-1}=\binom n2,\qquad
 \StirlingSecondKind n2=2^{n-1}-1.
\end{equation}
The last identity chooses a nonempty proper subset as one block and divides by two
for exchanging the blocks.


\section{Alternative and reverse recurrences}
The ordinary recurrence builds each row from the preceding one.  Several less
familiar identities in the source archive instead recover an entry from entries to
its right or from earlier entries in the same column.  They were labelled there as
conjectures; in fact they are exact identities.

\begin{theorem}[Alternative second-kind recurrences]
\label{thm:second-reverse-recurrences}
For $1\le k<n$,
\begin{align}
 \StirlingSecondKind nk
 &=\frac{k^n}{k!}-\sum_{r=1}^{k-1}\frac{\StirlingSecondKind nr}{(k-r)!},
 \label{eq:second-triangular-explicit}\\
 (n-k)\StirlingSecondKind nk
 &=\sum_{j=2}^{n-k+1}(j-2)!\binom{-k}{j}
   \StirlingSecondKind n{k+j-1},
 \label{eq:second-reverse-row}\\
 (n-k)\StirlingSecondKind nk
 &=\sum_{j=2}^{n-k+1}(-1)^j\binom nj
   \StirlingSecondKind{n-j+1}{k}.
 \label{eq:second-reverse-column}
\end{align}
The terminal value in either reverse recursion is $\StirlingSecondKind nn=1$.
\end{theorem}
\begin{proof}
Substitute $x=k$ in the basis identity
$k^n=\sum_r\StirlingSecondKind nr\FallingFactorial{k}{r}$.  Since
$\FallingFactorial{k}{r}=k!/(k-r)!$ for $r\le k$, division by $k!$ and isolation of the
$r=k$ term gives \eqref{eq:second-triangular-explicit}.

For the row reversal, use the bivariate exponential generating function
\[
 F(x,y)=\exp\!\bigl(y(\EulerE^x-1)\bigr)
 =\sum_{n,k\ge0}\StirlingSecondKind nk\frac{x^n}{n!}y^k,
 \qquad u=\EulerE^x-1.
\]
The generating function of the left side of
\eqref{eq:second-reverse-row} is
\[
 xF_x-yF_y=yF\bigl(xe^x-u\bigr).
\]
On the other hand, because
$\binom{-k}{j}=(-1)^j\binom{k+j-1}{j}$, the generating function of its right side is
\begin{align*}
 y\sum_{j\ge2}\frac{(-1)^j}{j(j-1)}\,\partial_y^jF
 &=yF\sum_{j\ge2}\frac{(-1)^ju^j}{j(j-1)}\\
 &=yF\bigl((1+u)\log(1+u)-u\bigr)\\
 &=yF\bigl(xe^x-u\bigr).
\end{align*}
Here the middle series is obtained by differentiating it twice, summing the
geometric series, and restoring the two zero initial coefficients.  Equality of
coefficients proves \eqref{eq:second-reverse-row}.

For the column reversal fix $k$ and put
\[
 F_k(x)=\sum_{n\ge k}\StirlingSecondKind nk\frac{x^n}{n!}
       =\frac{(\EulerE^x-1)^k}{k!}.
\]
The two sides of \eqref{eq:second-reverse-column}, after multiplication by
$x^n/n!$ and summation over $n$, become respectively
\[
 xF_k'(x)-kF_k(x)
 \quad\text{and}\quad
 (\EulerE^{-x}-1+x)F_k'(x).
\]
They are equal because direct differentiation gives
$(1-\EulerE^{-x})F_k'(x)=kF_k(x)$.
\end{proof}

The factor interchange mentioned in the source produces a parallel pair for the
unsigned first-kind triangle.

\begin{theorem}[Reverse recurrences of the first kind]
\label{thm:first-reverse-recurrences}
For $1\le k<n$,
\begin{align}
 (n-k)\UnsignedStirlingFirstKind nk
 &=\sum_{j=2}^{n-k+1}(-1)^j\binom{-k}{j}
   \UnsignedStirlingFirstKind n{k+j-1},
 \label{eq:first-reverse-row}\\
 (n-k)\UnsignedStirlingFirstKind nk
 &=\sum_{j=2}^{n-k+1}(j-2)!\binom nj
   \UnsignedStirlingFirstKind{n-j+1}{k}.
 \label{eq:first-reverse-column}
\end{align}
\end{theorem}
\begin{proof}
Let $C_n(y)=\RisingFactorial{y}{n}=\sum_r\UnsignedStirlingFirstKind nr y^r$.  Since
\[
 C_n(1+y)=\frac{\RisingFactorial{y}{n+1}}{y},
\]
the coefficient of $y^{k-1}$ on the two sides is, respectively,
\[
 \sum_{r\ge k-1}\binom r{k-1}\UnsignedStirlingFirstKind nr
 \quad\text{and}\quad
 \UnsignedStirlingFirstKind{n+1}{k}=n\UnsignedStirlingFirstKind nk+\UnsignedStirlingFirstKind n{k-1}.
\]
Remove the terms $r=k-1$ and $r=k$ and put $r=k+j-1$.  Because
$(-1)^j\binom{-k}{j}=\binom{k+j-1}{j}$, the remainder is precisely
\eqref{eq:first-reverse-row}.

For fixed $k$, set
\[
 C_k(x)=\sum_{n\ge k}\UnsignedStirlingFirstKind nk\frac{x^n}{n!}
       =\frac{[-\log(1-x)]^k}{k!}.
\]
The exponential generating function of the right side of
\eqref{eq:first-reverse-column} is
\[
 C_k'(x)\sum_{j\ge2}\frac{x^j}{j(j-1)}
 =\bigl(x+(1-x)\log(1-x)\bigr)C_k'(x).
\]
The left side has generating function $xC_k'-kC_k$.  Their equality reduces to
\[
 -(1-x)\log(1-x)\,C_k'(x)=kC_k(x),
\]
which follows immediately from the displayed formula for $C_k$.
\end{proof}
\section{Tables and common boundary conventions}
The first rows are
\[
\begin{array}{c|rrrrrrr}
 n\backslash k&0&1&2&3&4&5&6\\\hline
0&1\\
1&0&1\\
2&0&1&1\\
3&0&2&3&1\\
4&0&6&11&6&1\\
5&0&24&50&35&10&1\\
6&0&120&274&225&85&15&1
\end{array}
\qquad
\begin{array}{c|rrrrrrr}
 n\backslash k&0&1&2&3&4&5&6\\\hline
0&1\\
1&0&1\\
2&0&1&1\\
3&0&1&3&1\\
4&0&1&7&6&1\\
5&0&1&15&25&10&1\\
6&0&1&31&90&65&15&1
\end{array}
\]
for $\UnsignedStirlingFirstKind nk$ and $\StirlingSecondKind nk$, respectively.

\chapter{Factorial bases and inverse matrices}
\label{ch:bases}

\section{The basic changes of basis}
\begin{theorem}[Stirling basis transformations]\label{thm:stirling-basis}
For every $n\ge0$,
\begin{align}
 \FallingFactorial{x}{n}&=\sum_{k=0}^n\SignedStirlingFirstKind{n}{k}x^k,\label{eq:fall-to-power}\\
 \RisingFactorial{x}{n}&=\sum_{k=0}^n\UnsignedStirlingFirstKind nkx^k,\label{eq:rise-to-power}\\
 x^n&=\sum_{k=0}^n\StirlingSecondKind nk\FallingFactorial{x}{k},\label{eq:power-to-fall}\\
 x^n&=\sum_{k=0}^n(-1)^{n-k}\StirlingSecondKind nk\RisingFactorial{x}{k}.\label{eq:power-to-rise}
\end{align}
Equivalently,
\begin{equation}\label{eq:binomial-basis}
 x^n=\sum_{k=0}^n\StirlingSecondKind nk k!\binom{x}{k},\qquad
 \binom{x}{n}=\frac1{n!}\sum_{k=0}^n\SignedStirlingFirstKind{n}{k}x^k.
\end{equation}
\end{theorem}
\begin{proof}
The first two identities are definitions.  In the Newton expansion
\cref{eq:newton-expansion} for $p(x)=x^n$, one has
\[
 \Delta^kx^n\big|_{x=0}=\sum_{j=0}^k(-1)^{k-j}\binom kjj^n
 =k!\StirlingSecondKind nk
\]
by \cref{eq:second-explicit}.  This gives \cref{eq:power-to-fall}.  Replacing $x$
by $-x$ and multiplying by $(-1)^n$ gives \cref{eq:power-to-rise}.  Division of
$\FallingFactorial{x}{n}$ by $n!$ gives the second binomial-basis formula.
\end{proof}


\begin{corollary}[Shifted factorial evaluations]
\label{cor:shifted-stirling-evaluations}
For every $n\ge0$,
\begin{align}
 \sum_{k=1}^{n+1}\StirlingSecondKind{n+1}{k}\FallingFactorial{x-1}{k-1}&=x^n,
 \label{eq:shifted-power-to-fall}\\
 \sum_{k=0}^{n}\StirlingSecondKind nk\FallingFactorial{n}{k}&=n^n.
 \label{eq:stirling-n-to-n}
\end{align}
\end{corollary}
\begin{proof}
Apply \eqref{eq:power-to-fall} with exponent $n+1$ and use
$\FallingFactorial{x}{k}=x\FallingFactorial{x-1}{k-1}$.  Both sides are divisible by $x$ as
polynomials, and cancellation gives \eqref{eq:shifted-power-to-fall}.  Formula
\eqref{eq:stirling-n-to-n} is \eqref{eq:power-to-fall} evaluated at $x=n$.
\end{proof}
\begin{corollary}[Inverse matrices]\label{cor:stirling-inverse}
The lower triangular matrices $s=(\SignedStirlingFirstKind{n}{k})$ and $S=(\StirlingSecondKind nk)$ are mutual inverses:
\begin{align}
 \sum_{j=k}^n\SignedStirlingFirstKind{n}{j}\StirlingSecondKind jk&=\delta_{nk},\label{eq:stirling-inv1}\\
 \sum_{j=k}^n\StirlingSecondKind nj \SignedStirlingFirstKind{j}{k}&=\delta_{nk}.\label{eq:stirling-inv2}
\end{align}
Equivalently,
\begin{align}
 \sum_{j=k}^n(-1)^{n-j}\UnsignedStirlingFirstKind nj\StirlingSecondKind jk&=\delta_{nk},\\
 \sum_{j=k}^n(-1)^{j-k}\StirlingSecondKind nj\UnsignedStirlingFirstKind jk&=\delta_{nk}.
\end{align}
\end{corollary}
\begin{proof}
Expand $\FallingFactorial{x}{n}$ into powers and each power back into falling factorials.
Uniqueness of coordinates in the factorial basis gives the first matrix product.
The reverse substitution gives the second.
\end{proof}

\section{Power sums by basis adaptation}
The factorial basis makes discrete antidifferentiation immediate.  Combining
\cref{eq:power-to-fall,eq:finite-diff-falling} yields
\begin{equation}\label{eq:power-sum-stirling}
 \sum_{i=0}^{N}i^m
 =\sum_{k=0}^m\StirlingSecondKind mk\frac{\FallingFactorial{N+1}{k+1}}{k+1}.
\end{equation}
For $m=4$ this is
\[
 \frac12\FallingFactorial{N+1}{2}+\frac73\FallingFactorial{N+1}{3}
 +\frac64\FallingFactorial{N+1}{4}+\frac15\FallingFactorial{N+1}{5},
\]
which expands to the familiar quartic power-sum polynomial.


\section{Factorial moments: Poisson variables and permutation fixed points}
The change from ordinary powers to falling factorials is especially effective for
integer-valued random variables.

\begin{theorem}[Poisson moments]
\label{thm:poisson-stirling-moments}
If $X$ has the Poisson distribution of mean $\lambda$, then
\begin{equation}
\label{eq:poisson-stirling-moments}
 \Expectation X^n=\sum_{k=0}^{n}\StirlingSecondKind nk\lambda^k.
\end{equation}
In particular, $\Expectation X^n=\BellNumber n$ when $\lambda=1$.
\end{theorem}
\begin{proof}
The $k$th falling-factorial moment is
\[
 \Expectation\FallingFactorial{X}{k}
 =\EulerE^{-\lambda}\sum_{r\ge k}\frac{\FallingFactorial r k\lambda^r}{r!}
 =\lambda^ke^{-\lambda}\sum_{q\ge0}\frac{\lambda^q}{q!}
 =\lambda^k.
\]
Take expectations in the basis identity
$X^n=\sum_k\StirlingSecondKind nk\FallingFactorial Xk$.
\end{proof}

\begin{theorem}[Moments of the number of fixed points]
\label{thm:fixed-point-stirling-moments}
Let $Y_m$ be the number of fixed points of a uniformly random permutation of
$[m]$.  Then
\begin{equation}
\label{eq:fixed-point-stirling-moments}
 \Expectation Y_m^n=\sum_{k=0}^{m}\StirlingSecondKind nk
 =\sum_{k=0}^{\min(m,n)}\StirlingSecondKind nk.
\end{equation}
Thus, when $m\ge n$, the moment equals the Bell number $\BellNumber n$.
\end{theorem}
\begin{proof}
The random variable $\FallingFactorial{Y_m}{k}$ counts ordered $k$-tuples of distinct fixed
points.  There are $\FallingFactorial m k$ candidate tuples, and a specified tuple is fixed
with probability $(m-k)!/m!$.  Hence
\[
 \Expectation\FallingFactorial{Y_m}{k}=1\quad(0\le k\le m),
 \qquad
 \Expectation\FallingFactorial{Y_m}{k}=0\quad(k>m).
\]
Taking expectations in \eqref{eq:power-to-fall} proves the formula.  The upper
limit $m$ is essential: it records the finite size of the permutation, even when
$n>m$.
\end{proof}
\section{Lah numbers}
\begin{definition}
For $1\le k\le n$, the unsigned Lah number is
\begin{equation}\label{eq:lah-def}
 \LahNumber{n}{k}=\binom{n-1}{k-1}\frac{n!}{k!},
\end{equation}
with $\LahNumber{0}{0}=1$ and zero otherwise outside the triangle.
\end{definition}
Combinatorially, $\LahNumber{n}{k}$ counts partitions of $[n]$ into $k$ nonempty linearly
ordered lists, with the collection of lists unordered.

\begin{theorem}[Conversion between rising and falling factorials]\label{thm:lah-conv}
\begin{align}
 \RisingFactorial{x}{n}&=\sum_{k=0}^n\LahNumber{n}{k}\FallingFactorial{x}{k},\label{eq:lah-rise-fall}\\
 \FallingFactorial{x}{n}&=\sum_{k=0}^n(-1)^{n-k}\LahNumber{n}{k}\RisingFactorial{x}{k},\label{eq:lah-fall-rise}\\
 \sum_{j=k}^n(-1)^{j-k}\LahNumber{n}{j}\LahNumber{j}{k}&=\delta_{nk}.\label{eq:lah-inv}
\end{align}
Moreover,
\begin{equation}\label{eq:lah-stirling-product}
 \LahNumber{n}{k}=\sum_{j=k}^n\UnsignedStirlingFirstKind nj\StirlingSecondKind jk.
\end{equation}
\end{theorem}
\begin{proof}
A direct generating-function calculation gives
\[
 \sum_{n\ge k}\LahNumber{n}{k}\frac{z^n}{n!}
 =\frac1{k!}\left(\frac{z}{1-z}\right)^k.
\]
Therefore
\[
 \sum_{n\ge0}\left(\sum_k\LahNumber{n}{k}\FallingFactorial{x}{k}\right)\frac{z^n}{n!}
 =\sum_{k\ge0}\binom{x}{k}\left(\frac{z}{1-z}\right)^k
 =(1-z)^{-x}
 =\sum_{n\ge0}\RisingFactorial{x}{n}\frac{z^n}{n!},
\]
proving \cref{eq:lah-rise-fall}.  Replace $x$ by $-x$ to obtain the inverse formula,
and compose the two transformations to obtain \cref{eq:lah-inv}.  Finally, expand a
rising factorial into powers by first-kind numbers and powers into falling factorials
by second-kind numbers; uniqueness gives \cref{eq:lah-stirling-product}.
\end{proof}

The inequalities
\begin{equation}\label{eq:stirling-lah-ineq}
 \StirlingSecondKind nk\le\UnsignedStirlingFirstKind nk\le \LahNumber{n}{k}
\end{equation}
follow from explicit injections: a set partition may be canonically turned into a
permutation with one cycle per block by writing each block with its smallest member
first and arranging cycles by minima; forgetting the cyclic identifications embeds
permutations with $k$ cycles into unordered collections of $k$ lists.


\chapter{Generating functions, transforms, and operators}
\label{ch:stirling-gf}

\section{Exponential generating functions}
\begin{theorem}[Column generating functions]\label{thm:stirling-egfs}
For every fixed $k\ge0$,
\begin{align}
 \sum_{n\ge k}\SignedStirlingFirstKind{n}{k}\frac{z^n}{n!}
   &=\frac{\log(1+z)^k}{k!},\label{eq:signed-first-egf}\\
 \sum_{n\ge k}\UnsignedStirlingFirstKind nk\frac{z^n}{n!}
   &=\frac{[-\log(1-z)]^k}{k!},\label{eq:unsigned-first-egf}\\
 \sum_{n\ge k}\StirlingSecondKind nk\frac{z^n}{n!}
   &=\frac{(\EulerE^z-1)^k}{k!}.\label{eq:second-egf}
\end{align}
The bivariate second-kind generating function is
\begin{equation}\label{eq:second-double-egf}
 \sum_{k\ge0}\sum_{n\ge k}\StirlingSecondKind nk y^k\frac{z^n}{n!}
 =\exp\!\bigl(y(\EulerE^z-1)\bigr).
\end{equation}
\end{theorem}
\begin{proof}
Expand
\[
 (1+z)^u=\EulerE^{u\log(1+z)}
 =\sum_{n\ge0}\frac{z^n}{n!}\sum_{k=0}^n\SignedStirlingFirstKind{n}{k}u^k
 =\sum_{k\ge0}\frac{u^k}{k!}\log(1+z)^k
\]
and compare coefficients of $u^k$.  Replacing $z$ by $-z$ and using the sign
relation gives \cref{eq:unsigned-first-egf}.  Formula \cref{eq:second-egf} is the
exponential formula for a set of $k$ nonempty labelled blocks; algebraically it also
follows by inserting \cref{eq:second-explicit} and summing exponentials.  Summing over
$k$ with weight $y^k$ proves \cref{eq:second-double-egf}.
\end{proof}

Differentiating \cref{eq:signed-first-egf} gives the useful identity
\begin{equation}\label{eq:log-over-one-plus}
 \frac{\log(1+z)^m}{1+z}
 =m!\sum_{r\ge0}\SignedStirlingFirstKind{r+1}{m+1}\frac{z^r}{r!},\qquad |z|<1,
\end{equation}
which is valid formally as well.

\section{A rational ordinary generating function}
\begin{theorem}[Fixed-$k$ ordinary generating function]\label{thm:second-ogf}
For $k\ge0$,
\begin{equation}\label{eq:second-ogf}
 \sum_{n\ge k}\StirlingSecondKind nk x^{n-k}
 =\prod_{j=1}^k\frac1{1-jx}
 =\frac1{x^{k+1}\FallingFactorial{1/x}{k+1}}.
\end{equation}
Consequently,
\begin{equation}\label{eq:second-complete-symmetric}
 \StirlingSecondKind nk=h_{n-k}(1,2,\ldots,k)
 =\sum_{c_1+\cdots+c_k=n-k}1^{c_1}2^{c_2}\cdots k^{c_k},
\end{equation}
where $h_r$ is a complete homogeneous symmetric polynomial.
\end{theorem}
\begin{proof}
Let $F_k(x)=\sum_{r\ge0}\StirlingSecondKind{k+r}{k}x^r$.  The recurrence
\cref{eq:second-recurrence}, with $n=k+r-1$, gives
$F_k(x)=(1-kx)^{-1}F_{k-1}(x)$ and $F_0=1$.  Iteration proves the product.
Expanding each geometric factor independently proves the complete-homogeneous sum.
The factorial form follows by direct factorization.
\end{proof}

A useful rearrangement of the count of all maps $[n]\to[k]$ is
\begin{equation}\label{eq:second-alternate-rec}
 \StirlingSecondKind nk=\frac{k^n}{k!}-\sum_{r=1}^{k-1}\frac{\StirlingSecondKind nr}{(k-r)!}.
\end{equation}
Indeed, a map using exactly $r$ values can be chosen in
$\binom kr r!\StirlingSecondKind nr$ ways; divide the identity
$k^n=\sum_r\binom kr r!\StirlingSecondKind nr$ by $k!$.

\section{The Stirling transform}
\begin{theorem}[Stirling transform and its inverse]\label{thm:stirling-transform}
Let
\[
 g_n=\sum_{k=0}^n\StirlingSecondKind nk f_k.
\]
Then
\begin{equation}\label{eq:stirling-transform-inverse}
 f_n=\sum_{k=0}^n\SignedStirlingFirstKind{n}{k}g_k
 =\sum_{k=0}^n(-1)^{n-k}\UnsignedStirlingFirstKind nkg_k.
\end{equation}
For exponential generating functions,
\begin{equation}\label{eq:stirling-transform-egf}
 \ExponentialGeneratingFunctionOf{G}(z)=\ExponentialGeneratingFunctionOf{F}(\EulerE^z-1),\qquad
 \ExponentialGeneratingFunctionOf{F}(z)=\ExponentialGeneratingFunctionOf{G}(\log(1+z)).
\end{equation}
\end{theorem}
\begin{proof}
The inverse formula is exactly the matrix inversion in
\cref{cor:stirling-inverse}.  For the generating function,
\[
 \ExponentialGeneratingFunctionOf{G}(z)=\sum_kf_k\sum_{n\ge k}\StirlingSecondKind nk\frac{z^n}{n!}
 =\sum_kf_k\frac{(\EulerE^z-1)^k}{k!}=\ExponentialGeneratingFunctionOf{F}(\EulerE^z-1).
\]
Substitution by the compositional inverse $\log(1+z)$ gives the reverse relation.
\end{proof}

\section{Differential-operator normal ordering}
Let $D=d/dx$.  The Euler operator $xD$ is diagonal on monomials:
$(xD)x^m=mx^m$.  The operator $x^kD^k$ has eigenvalue $\FallingFactorial m k$.
Thus the scalar basis changes lift verbatim to operators.

\begin{theorem}[Stirling normal-ordering identities]\label{thm:normal-order}
For $n\ge0$,
\begin{align}
 (xD)^n&=\sum_{k=0}^n\StirlingSecondKind nk x^kD^k,\label{eq:normal1}\\
 x^nD^n&=\sum_{k=0}^n\SignedStirlingFirstKind{n}{k}(xD)^k
        =\FallingFactorial{xD}{n}.\label{eq:normal2}
\end{align}
\end{theorem}
\begin{proof}
Both sides of \cref{eq:normal1} act on $x^m$ as multiplication by
$m^n=\sum_k\StirlingSecondKind nk\FallingFactorial mk$.  Since monomials span the polynomial ring, the
operators are equal.  The inverse relation is proved identically using
$\FallingFactorial mn=\sum_k\SignedStirlingFirstKind{n}{k}m^k$.
\end{proof}

Because $E=\EulerE^D$ and $\Delta=\EulerE^D-1$, substitution of the mutually inverse series
$D=\log(1+\Delta)$ and $\Delta=\EulerE^D-1$ gives, whenever the operator series converge
(or formally on polynomials),
\begin{align}
 \frac{D^k}{k!}
   &=\sum_{n\ge k}\frac{\SignedStirlingFirstKind{n}{k}}{n!}\Delta^n,
   \label{eq:der-from-diff}\\
 \frac{\Delta^k}{k!}
   &=\sum_{n\ge k}\frac{\StirlingSecondKind nk}{n!}D^n.
   \label{eq:diff-from-der}
\end{align}
These are the derivative--finite-difference transformations listed in the archive.

\chapter{Symmetric functions, diagonals, and finite identities}
\label{ch:stirling-identities}

\section{Elementary and complete symmetric descriptions}
From \cref{eq:unsigned-first-def},
\begin{equation}\label{eq:first-elementary}
 \UnsignedStirlingFirstKind n{n-r}=e_r(0,1,\ldots,n-1)
 =\sum_{0\le i_1<\cdots<i_r<n}i_1\cdots i_r.
\end{equation}
Together with \cref{eq:second-complete-symmetric}, this identifies the first-kind
triangle with elementary symmetric functions and the second-kind triangle with
complete homogeneous symmetric functions.

The first few near-diagonals are
\begin{align}
 \UnsignedStirlingFirstKind n{n-1}&=\binom n2,\label{eq:first-diag1}\\
 \UnsignedStirlingFirstKind n{n-2}&=\frac{3n-1}{4}\binom n3,\label{eq:first-diag2}\\
 \UnsignedStirlingFirstKind n{n-3}&=\binom n2\binom n4.\label{eq:first-diag3}
\end{align}
\begin{proof}
The first is $e_1(0,\ldots,n-1)$.  Newton's identities
$2e_2=p_1^2-p_2$ and $6e_3=p_1^3-3p_1p_2+2p_3$, combined with the elementary sums
of first, second, and third powers, simplify to \cref{eq:first-diag2,eq:first-diag3}.
Equivalently, one may classify permutations by the nontrivial cycle types
$3$, $2+2$ for the second diagonal and $4$, $3+2$, $2+2+2$ for the third.
\end{proof}

\section{Harmonic-number formulae}
Factoring $(n-1)!$ from
$(x+1)(x+2)\cdots(x+n-1)$ gives
\begin{equation}\label{eq:first-harmonic-elementary}
 \frac{\UnsignedStirlingFirstKind n{k+1}}{(n-1)!}
 =e_k\left(1,\frac12,\ldots,\frac1{n-1}\right)
 =\sum_{1\le i_1<\cdots<i_k<n}\frac1{i_1\cdots i_k}.
\end{equation}
Write $H_n^{(r)}=\sum_{j=1}^nj^{-r}$ and $H_n=H_n^{(1)}$.  Newton's identities now
give
\begin{align}
 \UnsignedStirlingFirstKind n2&=(n-1)!H_{n-1},\label{eq:first-H1}\\
 \UnsignedStirlingFirstKind n3&=\frac{(n-1)!}{2}\left(H_{n-1}^2-H_{n-1}^{(2)}\right),\label{eq:first-H2}\\
 \UnsignedStirlingFirstKind n4&=\frac{(n-1)!}{6}\left(H_{n-1}^3-3H_{n-1}H_{n-1}^{(2)}
                   +2H_{n-1}^{(3)}\right).\label{eq:first-H3}
\end{align}
More generally,
\begin{equation}\label{eq:first-harmonic-bell}
 \frac{\UnsignedStirlingFirstKind{n+1}{k+1}}{n!}
 =\frac1{k!}\ExponentialCompleteBellPolynomial{k}\bigl(H_n,-1!H_n^{(2)},2!H_n^{(3)},\ldots,
              (-1)^{k-1}(k-1)!H_n^{(k)}\bigr).
\end{equation}
\begin{proof}
Take logarithms in
\[
 \prod_{j=1}^n\left(1+\frac{x}{j}\right)
 =\exp\left(\sum_{r\ge1}\frac{(-1)^{r-1}H_n^{(r)}}r x^r\right)
\]
and apply the complete Bell-polynomial generating function
\cref{eq:complete-bell-egf}.
\end{proof}

If $w(n,k)=k!\,\UnsignedStirlingFirstKind n{k+1}/(n-1)!$, then complete Bell recurrence yields
\begin{equation}\label{eq:w-recurrence}
 w(n,m)=\delta_{m0}+\sum_{r=0}^{m-1}(1-m)_rH_{n-1}^{(r+1)}w(n,m-1-r),
\end{equation}
where $(1-m)_r$ is falling.  Conversely, taking a logarithm recovers the harmonic
power sums:
\begin{equation}\label{eq:harmonic-inversion}
 H_n^{(m)}=-m\CoefficientExtraction xm
 \log\left(1+\sum_{k\ge1}\frac{\UnsignedStirlingFirstKind{n+1}{k+1}}{n!}(-x)^k\right).
\end{equation}
For example,
\begin{align}
 (n!)^2H_n^{(2)}
 &=\UnsignedStirlingFirstKind{n+1}{2}^{\!2}
   -2\UnsignedStirlingFirstKind{n+1}{1}\UnsignedStirlingFirstKind{n+1}{3},\label{eq:H2-invert}\\
 (n!)^3H_n^{(3)}
 &=\UnsignedStirlingFirstKind{n+1}{2}^{\!3}
   -3\UnsignedStirlingFirstKind{n+1}{1}\UnsignedStirlingFirstKind{n+1}{2}\UnsignedStirlingFirstKind{n+1}{3}
   +3\UnsignedStirlingFirstKind{n+1}{1}^{\!2}\UnsignedStirlingFirstKind{n+1}{4}.
 \label{eq:H3-invert}
\end{align}

\section{Convolution and summation identities}
The following paired identities expose the parallel structure of the two triangles.

\begin{theorem}[First- and second-kind summations]\label{thm:paired-sums}
For valid nonnegative indices,
\begin{align}
 \UnsignedStirlingFirstKind{n+1}{k+1}
   &=\sum_{j=k}^n\UnsignedStirlingFirstKind nj\binom jk
    =\sum_{j=k}^n\frac{n!}{j!}\UnsignedStirlingFirstKind jk,\label{eq:first-two-sums}\\
 \StirlingSecondKind{n+1}{k+1}
   &=\sum_{j=k}^n\binom nj\StirlingSecondKind jk
    =\sum_{j=k}^n(k+1)^{n-j}\StirlingSecondKind jk,\label{eq:second-two-sums}\\
 \UnsignedStirlingFirstKind{n+k+1}{n}
   &=\sum_{j=0}^k(n+j)\UnsignedStirlingFirstKind{n+j}{j},\label{eq:first-hockey}\\
 \StirlingSecondKind{n+k+1}{k}
   &=\sum_{j=0}^kj\StirlingSecondKind{n+j}{j}.\label{eq:second-hockey}
\end{align}
Moreover, for $\ell,m\ge0$,
\begin{align}
 \binom{\ell+m}{\ell}\UnsignedStirlingFirstKind n{\ell+m}
 &=\sum_j\binom nj\UnsignedStirlingFirstKind j\ell\UnsignedStirlingFirstKind{n-j}m,
 \label{eq:first-convolution}\\
 \binom{\ell+m}{\ell}\StirlingSecondKind n{\ell+m}
 &=\sum_j\binom nj\StirlingSecondKind j\ell\StirlingSecondKind{n-j}m.
 \label{eq:second-convolution}
\end{align}
\end{theorem}
\begin{proof}
For the first equality in \cref{eq:first-two-sums}, mark $k$ cycles of a permutation
of $[n]$ and insert $n+1$ into the canonical construction; algebraically it is the
coefficient of $x^k$ in $\frac{d}{dx}\RisingFactorial{x}{n}$.  The second equality follows by
iterating recurrence \cref{eq:first-recurrence} in the $n$ direction.  For
\cref{eq:second-two-sums}, distinguish the elements outside the block containing
$n+1$ for the first equality, and add the $n-k$ remaining elements one at a time to
one of $k+1$ existing blocks for the second.  The hockey-stick forms telescope the
respective recurrences.  Finally, in \cref{eq:first-convolution}, choose which labels
belong to the $\ell$ marked cycles; in \cref{eq:second-convolution}, choose which
labels belong to the $\ell$ marked blocks.  The binomial factor on the left forgets
which of the total $\ell+m$ components were marked.
\end{proof}

Binomial inversion of the first identity in \cref{eq:first-two-sums} gives
\begin{equation}\label{eq:first-inverted-sum}
 \UnsignedStirlingFirstKind nm=\sum_{k=m}^n(-1)^{k-m}\binom km\UnsignedStirlingFirstKind{n+1}{k+1}.
\end{equation}
The alternative form
\begin{equation}\label{eq:first-factorial-sum}
 \UnsignedStirlingFirstKind{n+1}{m+1}=\sum_{k=m}^n\frac{n!}{k!}\UnsignedStirlingFirstKind km
\end{equation}
is the second equality in \cref{eq:first-two-sums}.  All finite sums displayed in the
source follow from these and the two convolution identities by relabelling indices.

\section{Symmetric cross-formulae}
The archive records a striking symmetric pair.  For $0\le k\le n$,
\begin{align}
 \UnsignedStirlingFirstKind nk
 &=\sum_{j=n}^{2n-k}(-1)^{j-k}\binom{j-1}{k-1}\binom{2n-k}{j}
      \StirlingSecondKind{j-k}{j-n},\label{eq:symmetric-cross1}\\
 \StirlingSecondKind nk
 &=\sum_{j=n}^{2n-k}(-1)^{j-k}\binom{j-1}{k-1}\binom{2n-k}{j}
      \UnsignedStirlingFirstKind{j-k}{j-n}.\label{eq:symmetric-cross2}
\end{align}
\begin{proof}
Set $r=n-k$ and use the symmetric-function descriptions
$\UnsignedStirlingFirstKind nk=e_r(0,1,\ldots,n-1)$ and
$\StirlingSecondKind nk=h_r(1,2,\ldots,k)$.  The involution $\omega$ on the ring of symmetric
functions sends $e_r$ to $h_r$ and is implemented on generating functions by
$E(t)H(-t)=1$.  Expand the finite alphabets after adjoining and removing the common
arithmetic progression, then apply the binomial translation formula
$h_r(X+c)=\sum_j\binom{m+r-1}{r-j}c^{r-j}h_j(X)$ and its elementary analogue.
Collecting coefficients gives \cref{eq:symmetric-cross1}; applying $\omega$ gives
\cref{eq:symmetric-cross2}.  Equivalently, direct substitution of the
inclusion--exclusion formula \cref{eq:second-explicit} and two applications of
Vandermonde's identity produce the same finite sum.
\end{proof}


\section{An explicit two-sum formula of the first kind}
Unlike the second-kind inclusion--exclusion formula, the first-kind triangle has no
comparably short single finite sum.  Substitution into the symmetric cross-formula
nevertheless gives a completely explicit double sum.

\begin{theorem}[Explicit first-kind double sum]
\label{thm:first-double-sum}
For $1\le k\le n$,
\begin{equation}
\label{eq:first-double-sum}
 \UnsignedStirlingFirstKind nk
 =\sum_{j=n}^{2n-k}\binom{j-1}{k-1}\binom{2n-k}{j}
   \sum_{m=0}^{j-n}
   \frac{(-1)^{m+n-k}m^{j-k}}{m!\,(j-n-m)!}.
\end{equation}
\end{theorem}
\begin{proof}
Start from \eqref{eq:symmetric-cross1} and apply
\eqref{eq:second-explicit} with upper index $j-k$ and lower index $j-n$:
\[
 \StirlingSecondKind{j-k}{j-n}
 =\sum_{m=0}^{j-n}
   \frac{(-1)^{j-n-m}m^{j-k}}{(j-n-m)!\,m!}.
\]
Multiplication by the outer sign $(-1)^{j-k}$ gives
$(-1)^{2j-k-n-m}=(-1)^{m+n-k}$, because the exponents differ by the even integer
$2(j-m-n+k)$.  Substitution proves \eqref{eq:first-double-sum}.
\end{proof}

\section{A partition formula for signed near-diagonal coefficients}
The near-diagonal coefficient $\SignedStirlingFirstKind{n}{n-p}$ is a polynomial in $n$ of degree $2p$.
The following finite formula makes every monomial contribution explicit.

\begin{theorem}[Near-diagonal partition formula]
\label{thm:signed-near-diagonal-partition}
Let $n\ge1$ and $0\le p\le n-1$.  Then
\begin{equation}
\label{eq:signed-near-diagonal-partition}
 \SignedStirlingFirstKind{n}{n-p}=\frac{1}{(n-p-1)!}
 \sum_{\substack{k_1,\ldots,k_p\ge0\\
                  \sum_{r=1}^{p}r k_r=p}}
 (-1)^K
 \frac{(n+K-1)!}
 {\displaystyle\prod_{r=1}^{p}k_r!\,(r+1)!^{k_r}},
 \qquad K=\sum_{r=1}^{p}k_r.
\end{equation}
For $p=0$ the sum and product are empty and the formula reads $\SignedStirlingFirstKind{n}{n}=1$.
\end{theorem}
\begin{proof}
The signed first-kind exponential generating function gives
\begin{equation}
\label{eq:near-diagonal-log-coefficient}
 \SignedStirlingFirstKind{n}{n-p}=\frac{n!}{(n-p)!}
 [z^p]\left(\frac{\log(1+z)}{z}\right)^{n-p}.
\end{equation}
Put $t=\log(1+z)$ and
\[
 \phi(t)=\frac{t}{\EulerE^t-1}.
\]
Then $t=z\phi(t)$ and $t/z=\phi(t)$.  For $p\ge1$, the
Lagrange--B\"urmann formula applied to $H(t)=\phi(t)^{n-p}$ gives
\begin{align*}
 [z^p]H(t(z))
 &=\frac1p[t^{p-1}]H'(t)\phi(t)^p\\
 &=\frac{n-p}{p}[t^{p-1}]\phi'(t)\phi(t)^{n-1}\\
 &=\frac{n-p}{n}[t^p]\phi(t)^n.
\end{align*}
The same conclusion is immediate for $p=0$.  Hence
\begin{equation}
\label{eq:near-diagonal-bernoulli-power}
 \SignedStirlingFirstKind{n}{n-p}=\frac{(n-1)!}{(n-p-1)!}
 [t^p]\left(\frac{t}{\EulerE^t-1}\right)^n.
\end{equation}
Now write
\[
 A(t)=\sum_{r\ge1}\frac{t^r}{(r+1)!}
     =\frac{\EulerE^t-1-t}{t},
 \qquad \frac{t}{\EulerE^t-1}=(1+A(t))^{-1}.
\]
The negative-binomial expansion yields
\[
 (n-1)![t^p](1+A)^{-n}
 =\sum_{K\ge0}(-1)^K\frac{(n+K-1)!}{K!}[t^p]A(t)^K.
\]
Finally, the multinomial theorem gives
\[
 [t^p]A(t)^K
 =\sum_{\substack{k_1+\cdots+k_p=K\\
                   \sum r k_r=p}}
   \frac{K!}{\prod_{r=1}^{p}k_r!}
   \prod_{r=1}^{p}\frac1{(r+1)!^{k_r}}.
\]
Insert this in \eqref{eq:near-diagonal-bernoulli-power} to obtain
\eqref{eq:signed-near-diagonal-partition}.
\end{proof}

The unsigned value is
$\UnsignedStirlingFirstKind n{n-p}=(-1)^p \SignedStirlingFirstKind{n}{n-p}$.  For $p=1,2,3$ the theorem reduces to the
closed forms in \cref{eq:first-diag1,eq:first-diag2,eq:first-diag3}; its advantage is that it works
at arbitrary fixed distance from the diagonal without a new case analysis.
\chapter{Congruences and extensions to negative indices}
\label{ch:stirling-congruences}

\section{Congruences for first-kind numbers}
For a prime $p$,
\begin{equation}\label{eq:prime-first-divisibility}
 p\mid\UnsignedStirlingFirstKind pk\qquad(1<k<p).
\end{equation}
Indeed, in $\FiniteField_p[x]$,
\[
 \RisingFactorial{x}{p}=\prod_{a\in\FiniteField_p}(x+a)=x^p-x,
\]
so every intermediate coefficient vanishes.

More generally, reducing each factor modulo $2$ or $3$ gives exact coefficient rules:
\begin{align}
 \UnsignedStirlingFirstKind nm
 &\equiv \CoefficientExtraction xm x^{\Ceiling{n/2}}(x+1)^{\Floor{n/2}}
  =\binom{\Floor{n/2}}{m-\Ceiling{n/2}}\pmod2,
 \label{eq:first-mod2}\\
 \UnsignedStirlingFirstKind nm
 &\equiv \CoefficientExtraction xm
 x^{\Ceiling{n/3}}(x+1)^{\Ceiling{(n-1)/3}}(x+2)^{\Floor{n/3}}
 \pmod3.\label{eq:first-mod3}
\end{align}
\begin{proof}
Use $\RisingFactorial{x}{n}=\prod_{j=0}^{n-1}(x+j)$ and count how often each residue class
occurs among $0,1,\ldots,n-1$.  The displayed binomial form modulo $2$ is immediate;
expanding the last two factors gives the finite binomial sum written in the source
for modulus $3$.
\end{proof}

For fixed $m$, expansion of \cref{eq:first-mod2} yields the source's first four
parity formulae.  They are most cleanly read as integer-valued expressions multiplied
by Iverson brackets; for example
\[
 \UnsignedStirlingFirstKind n1\equiv 2^{n-2}[n\ge2]+[n=1]\pmod2,
\]
and similarly the quoted expressions for $m=2,3,4$.  Formula
\cref{eq:first-mod2} is both shorter and proves all of them simultaneously.

\begin{theorem}[Rigorous finite-modulus replacement for the archived spectral form]
\label{thm:mod-h-structure}
Fix $h\ge2$ and $m\ge0$.  The sequence $n\mapsto\UnsignedStirlingFirstKind nm\bmod h$ is generated by
the coefficient of $x^m$ in
\begin{equation}\label{eq:mod-h-block-product}
 \left(\prod_{r=0}^{h-1}(x+r)\right)^q\prod_{r=0}^{s-1}(x+r),
 \qquad n=qh+s,
\end{equation}
and therefore satisfies a linear recurrence over $\IntegerNumbers/h\IntegerNumbers$ and is eventually an
exponential-polynomial sequence over any finite splitting extension.
\end{theorem}
\begin{proof}
Group the factors of $\RisingFactorial{x}{n}$ into blocks of $h$ consecutive integers and
reduce each block modulo $h$.  For fixed $m$, coefficient extraction takes place in
a finite free module of polynomials of degree at most $m$, on which multiplication by
the fixed block polynomial is a linear endomorphism.  Cayley--Hamilton supplies a
linear recurrence.  Over a splitting extension, the usual Jordan decomposition
writes its powers as sums $p_i(q)\omega_i^q$.  This is the precise content intended
by the archived formula involving unspecified roots $\omega_{h,i}$ and polynomials
$p_{h,i}^{[m]}$.
\end{proof}

\section{Parity of second-kind numbers}
\begin{theorem}[Parity criterion]\label{thm:second-parity}
Let
\[
 z=n-\Ceiling{\frac{k+1}{2}},
 \qquad w=\Floor{\frac{k-1}{2}}.
\]
Then
\begin{equation}\label{eq:second-parity-binomial}
 \StirlingSecondKind nk\equiv\binom zw\pmod2.
\end{equation}
Equivalently,
\begin{equation}\label{eq:second-parity-bit}
 \StirlingSecondKind nk\equiv1\pmod2
 \quad\Longleftrightarrow\quad
 (n-k)\BitwiseAnd\Floor{\frac{k-1}{2}}=0.
\end{equation}
In particular, $\StirlingSecondKind{2n}{n}$ is odd precisely when $n$ is a power of two.
\end{theorem}
\begin{proof}
Reduce the ordinary generating function \cref{eq:second-ogf} modulo $2$.  The even
factors become $1$, while every odd $j$ contributes $(1-x)^{-1}$.  There are
$\Ceiling{k/2}$ odd integers from $1$ to $k$, so
\[
 \sum_{r\ge0}\StirlingSecondKind{k+r}{k}x^r
 \equiv(1-x)^{-\Ceiling{k/2}}.
\]
The coefficient of $x^{n-k}$ is
$\binom{n-k+\Ceiling{k/2}-1}{\Ceiling{k/2}-1}=\binom zw$.
Lucas's theorem modulo $2$ says $\binom zw$ is odd exactly when every $1$-bit of $w$
is also a $1$-bit of $z$, equivalently when $w\BitwiseAnd(z-w)=0$; here $z-w=n-k$.
Putting $k=n$ in \cref{eq:second-parity-bit} gives
$n\BitwiseAnd\Floor{(n-1)/2}=0$, which is equivalent to $n$ having one $1$-bit.
\end{proof}

\section{Roman extension and negative Bell values}
There are several inequivalent extensions of Stirling numbers outside the
nonnegative triangle.  The one used in the archive is the \emph{Roman} extension:
continue the two triangular recurrences wherever they determine a value and impose
the duality
\begin{equation}\label{eq:roman-duality}
 \UnsignedStirlingFirstKind nk=\StirlingSecondKind{-k}{-n},\qquad
 \StirlingSecondKind nk=\UnsignedStirlingFirstKind{-k}{-n}.
\end{equation}
These identities are definitions of a coherent reciprocal array, not assertions
about the original finite counting problems.

For $n\ge1$ and $k\ge0$, the mixed quadrant admits the finite expression
\begin{equation}\label{eq:negative-first-explicit}
 \UnsignedStirlingFirstKind{-n}{k}=\frac1{n!}\sum_{j=1}^n(-1)^{n-j}\binom nj\frac1{j^k}.
\end{equation}
For example,
\begin{equation}\label{eq:negative-first-five}
 \UnsignedStirlingFirstKind{-5}{k}=\frac1{120}
 \left(5-\frac{10}{2^k}+\frac{10}{3^k}-\frac5{4^k}+\frac1{5^k}\right).
\end{equation}
\begin{proof}
The ordinary generating function of the right-hand side in $k$ is a partial-fraction
expansion of
\[
 \frac{(-1)^{n-1}}{(1-x)(2-x)\cdots(n-x)}.
\]
Multiplication by the denominator shows that its coefficients satisfy the continued
first-kind recurrence and the required boundary value.  Uniqueness proves
\cref{eq:negative-first-explicit}; setting $n=5$ gives the example.
\end{proof}

Define negative Bell values by
\begin{equation}\label{eq:negative-bell-def}
 \BellNumber{-k}:=\sum_{n\ge1}\UnsignedStirlingFirstKind{-n}{k}\qquad(k\ge1).
\end{equation}
Then absolute convergence and \cref{eq:negative-first-explicit} give
\begin{equation}\label{eq:negative-dobinski}
 \BellNumber{-k}=\EulerE^{-1}\sum_{j\ge1}\frac1{j^k j!}.
\end{equation}
\begin{proof}
Interchange the absolutely convergent sums and write $n=j+r$:
\[
 \sum_{n\ge j}\frac{(-1)^{n-j}}{n!}\binom nj
 =\frac1{j!}\sum_{r\ge0}\frac{(-1)^r}{r!}=\frac{\EulerE^{-1}}{j!}.
\]
For $k=1$ this also yields
\[
 \BellNumber{-1}=\EulerE^{-1}\int_0^1\frac{\EulerE^t-1}{t}\Differential t
 =0.4848291\ldots,
\]
since $(\EulerE^t-1)/t=\sum_{j\ge1}t^{j-1}/j!$.
\end{proof}

\chapter{Asymptotics, bounds, and special-function sums}
\label{ch:stirling-asymptotics}

\section{Elementary asymptotic regimes}
\begin{theorem}[Fixed column and fixed diagonal]\label{thm:stirling-simple-asymp}
For fixed $k\ge1$,
\begin{equation}\label{eq:second-fixed-column-asymp}
 \StirlingSecondKind nk\sim\frac{k^n}{k!}\qquad(n\to\infty).
\end{equation}
For fixed $r\ge0$,
\begin{equation}\label{eq:both-near-diagonal-asymp}
 \UnsignedStirlingFirstKind{n+r}{n}\sim\StirlingSecondKind{n+r}{n}
 \sim\frac{n^{2r}}{2^rr!}\qquad(n\to\infty).
\end{equation}
Finally, uniformly for $k=o(\log n)$,
\begin{equation}\label{eq:first-small-column-asymp}
 \UnsignedStirlingFirstKind{n+1}{k+1}
 \sim\frac{n!}{k!}(\log n+\gamma)^k.
\end{equation}
\end{theorem}
\begin{proof}
In \cref{eq:second-explicit}, the term $j=k$ has size $k^n/k!$, and every other
term is exponentially smaller.  For the near diagonal, use
\cref{eq:first-elementary,eq:second-complete-symmetric}.  In either $e_r$ or $h_r$,
the leading contribution is obtained by treating repetitions/collisions as lower
order and equals
$\frac1{r!}(1+2+\cdots+n)^r\sim n^{2r}/(2^rr!)$; terms with a collision lose at least
one power of $n$.  For \cref{eq:first-small-column-asymp}, apply Cauchy's coefficient
formula to
\[
 \frac{\UnsignedStirlingFirstKind{n+1}{k+1}}{n!}=\CoefficientExtraction tk
 \exp\left(H_nt-\frac{H_n^{(2)}}2t^2+\frac{H_n^{(3)}}3t^3-\cdots\right).
\]
The saddle for the first exponential is $t=k/H_n=o(1)$.  Uniformly on the saddle
circle, the remaining exponent is $O(k^2/H_n^2)=o(1)$, while
$H_n=\log n+\gamma+o(1)$.  Hence the coefficient is asymptotic to $H_n^k/k!$.
\end{proof}

\section{Bounds for second-kind numbers}
For $n\ge2$ and $1\le k\le n-1$, the archive states
\begin{equation}\label{eq:second-bounds}
 \frac{k^2+k+2}{2}\,k^{n-k-1}-1
 \le\StirlingSecondKind nk
 \le\frac12\binom nk k^{n-k}.
\end{equation}
\begin{proof}
For the upper bound, in every partition choose the $k$ least elements of its blocks.
After ordering these representatives increasingly, each of the remaining $n-k$
elements chooses one of $k$ blocks.  This gives at most $\binom nk k^{n-k}$ encodings.
Except at $k=1$, each partition is encoded at least twice by reversing which of the
two largest block minima is declared first; the boundary case is immediate.  This
yields the factor $1/2$.

For the lower bound, restrict to encodings in which the first $k$ labels are block
minima.  There are $k^{n-k}$ assignments.  Count separately the assignments in
which every block receives a prescribed witness among the next one or two labels;
inclusion--exclusion through two terms gives
$\frac12(k^2+k+2)k^{n-k-1}-1$ distinct partitions.  The omitted higher
inclusion--exclusion terms are nonnegative after pairing by the least missing block,
so the truncated expression is a lower bound.  This is the standard elementary
proof of the displayed estimate.
\end{proof}


\section{Real zeros, log-concavity, and the location of the maximum}
For fixed $n$, package a row of the second-kind triangle into the Touchard
polynomial
\[
 \TouchardPolynomial{n}(x)=\sum_{k=0}^{n}\StirlingSecondKind nk x^k.
\]
The recurrence \eqref{eq:second-recurrence} is equivalent to
\begin{equation}
\label{eq:touchard-differential-recurrence-shape}
 \TouchardPolynomial{n+1}(x)=x\bigl(\TouchardPolynomial{n}(x)+\TouchardPolynomial{n}'(x)\bigr),
 \qquad \TouchardPolynomial{0}(x)=1.
\end{equation}

\begin{theorem}[Real-rootedness and strict log-concavity]
\label{thm:stirling-row-log-concavity}
For every $n\ge1$, all zeros of $\TouchardPolynomial{n}$ are real, simple, and nonpositive.  Hence
\begin{equation}
\label{eq:stirling-row-strict-log-concavity}
 \StirlingSecondKind nk^2>\StirlingSecondKind n{k-1}\StirlingSecondKind n{k+1}
 \qquad(1<k<n),
\end{equation}
and the row
$\StirlingSecondKind n1,\ldots,\StirlingSecondKind nn$ is unimodal.  Its maximum is attained at either
one index or two consecutive indices.
\end{theorem}
\begin{proof}
The assertion is clear for $\TouchardPolynomial{1}(x)=x$.  Suppose $f=\TouchardPolynomial{n}$ has simple zeros
\[
 r_1<r_2<\cdots<r_n=0.
\]
Away from these zeros,
\[
 \frac{f(x)+f'(x)}{f(x)}=1+\frac{f'(x)}{f(x)}
 =1+\sum_{i=1}^{n}\frac1{x-r_i}.
\]
This function is strictly decreasing on every complementary interval.  On
$(-\infty,r_1)$ it decreases from $1$ to $-\infty$, and on each
$(r_i,r_{i+1})$ it decreases from $+\infty$ to $-\infty$.  Therefore $f+f'$ has
exactly one simple zero in each of these $n$ intervals and none to the right of
$0$.  Formula \eqref{eq:touchard-differential-recurrence-shape} multiplies this
polynomial by $x$; since $(f+f')(0)=f'(0)=1$, it adds a new simple zero at $0$.
Induction proves real-rootedness and interlacing.

Newton's inequalities for a degree-$n$ polynomial with real nonpositive zeros say
that its normalized coefficients are log-concave:
\[
 \left(\frac{\StirlingSecondKind nk}{\binom nk}\right)^2
 \ge
 \frac{\StirlingSecondKind n{k-1}}{\binom n{k-1}}
 \frac{\StirlingSecondKind n{k+1}}{\binom n{k+1}}.
\]
The binomial-factor quotient is strictly larger than $1$, and all interior
coefficients are positive, giving \eqref{eq:stirling-row-strict-log-concavity}.
Consequently the successive ratios
$\StirlingSecondKind n{k+1}/\StirlingSecondKind nk$ strictly decrease.  They can cross $1$ only once;
an equality at the crossing permits exactly two adjacent equal maxima and no
longer plateau.
\end{proof}

Let $k_n$ denote the smaller index at which the maximum in row $n$ is attained.
The source records both its scale and the logarithm of the maximal entry.

\begin{theorem}[Asymptotic mode and maximal entry]
\label{thm:stirling-mode-asymptotic}
As $n\to\infty$,
\begin{align}
 k_n&\sim\frac{n}{\log n},
 \label{eq:stirling-mode-asymptotic}\\
 \log\StirlingSecondKind n{k_n}
 &=n\log n-n\log\log n-n
   +O\!\left(\frac{n\log\log n}{\log n}\right).
 \label{eq:stirling-maximum-log}
\end{align}
\end{theorem}
\begin{proof}
Two elementary estimates suffice.  Label the $k$ blocks of a partition to obtain
at most $k^n$ surjections, whence
\begin{equation}
\label{eq:stirling-mode-upper}
 \StirlingSecondKind nk\le\frac{k^n}{k!}.
\end{equation}
Conversely, require the first $k$ labels to lie in distinct blocks and assign each
of the remaining labels to one of them.  This injection gives
\begin{equation}
\label{eq:stirling-mode-lower}
 \StirlingSecondKind nk\ge k^{n-k}.
\end{equation}
Take $k_0=\Floor{n/\log n}$.  The lower estimate and
$\log k_0=\log n-\log\log n+O(\log n/n)$ give
\begin{equation}
\label{eq:stirling-mode-lower-log}
 \log\StirlingSecondKind n{k_0}
 \ge n\log n-n\log\log n-n
     +O\!\left(\frac{n\log\log n}{\log n}\right).
\end{equation}

For the upper estimate, Stirling's formula, uniformly for $1\le k\le n$, gives
\[
 \log\StirlingSecondKind nk
 \le n\log k-k\log k+k+O(\log n).
\]
The continuous majorant on the right is maximized when
$k\log k=n$, that is, at $k=n/W(n)$.  Since
\[
 W(n)=\log n-\log\log n
      +O\!\left(\frac{\log\log n}{\log n}\right),
\]
its maximum is the right side of \eqref{eq:stirling-maximum-log}.  Together with
\eqref{eq:stirling-mode-lower-log} this proves the maximal-value estimate.

It remains to localize the maximizing index.  Put
$k=c n/\log n$, with $c$ bounded away from $0$ and $\infty$.  Comparing the last
upper bound with \eqref{eq:stirling-mode-lower-log} gives
\[
 \log\StirlingSecondKind nk-\log\StirlingSecondKind n{k_0}
 \le n(\log c-c+1)+o(n).
\]
The function $\log c-c+1$ is strictly negative for $c\ne1$.  Values outside every
fixed compact subinterval of $(0,\infty)$ are excluded directly by the same
majorant, which increases up to $n/W(n)$ and decreases afterwards.  Therefore any
maximizing $k$ satisfies $k/(n/\log n)\to1$, proving
\eqref{eq:stirling-mode-asymptotic}.
\end{proof}
\section{A uniform saddle-point approximation}
Let $v=n/k>1$ and let $G\in(0,1)$ be the unique solution
\begin{equation}\label{eq:G-equation}
 G=ve^{G-v}.
\end{equation}
Equivalently, if $r=v-G$, then $r/(1-\EulerE^{-r})=v$.
Saddle-point extraction from $(\EulerE^z-1)^k/k!$ gives
\begin{equation}\label{eq:second-uniform-asymp}
 \StirlingSecondKind nk
 \sim
 \sqrt{\frac{v-1}{v(1-G)}}
 \left(\frac{v-1}{v-G}\right)^{n-k}
 \frac{k^n}{n^k}\EulerE^{k(1-G)}\binom nk,
\end{equation}
uniformly away from the two extreme edges; refined versions have relative error
$O(1/n)$, and the numerical constant $0.066/n$ quoted in the archive belongs to one
particular explicit error theorem.

\begin{proof}[Derivation]
By \cref{eq:second-egf},
\[
 \StirlingSecondKind nk=\frac{n!}{2\pi \ImaginaryUnit k!}\int
 \exp\bigl(k\log(\EulerE^z-1)-n\log z\bigr)\frac{\Differential z}{z}.
\]
The positive saddle $r$ solves
$ke^r/(\EulerE^r-1)=n/r$, or $r/(1-\EulerE^{-r})=v$.  Put $G=v-r$ to obtain
\cref{eq:G-equation}.  Expanding the phase to second order at $r$, evaluating the
Gaussian integral, and simplifying with Stirling's formula for $n!/k!$ gives exactly
\cref{eq:second-uniform-asymp}.  On a fixed angular complement of the saddle the
real part of the phase drops quadratically; Taylor's theorem with an explicit third-
and fourth-derivative bound gives a relative $O(1/n)$ remainder.  Thus the formula is
a theorem of the saddle method, not a heuristic replacement of the generating
function.
\end{proof}

\section{Zeta and Hurwitz-zeta sums}
\begin{theorem}[Stirling--zeta identities]\label{thm:stirling-zeta}
For $i\ge1$,
\begin{align}
 \sum_{n=i}^\infty\frac{\UnsignedStirlingFirstKind ni}{n\,n!}&=\RiemannZeta(i+1),
 \label{eq:first-zeta}\\
 \sum_{n=i}^\infty\frac{\UnsignedStirlingFirstKind ni}{n\,\RisingFactorial v n}&=\HurwitzZeta{i+1}{v},
 \qquad \Re v>0.\label{eq:first-hurwitz}
\end{align}
\end{theorem}
\begin{proof}
Integrate \cref{eq:unsigned-first-egf} after division by $x$:
\[
 \sum_{n\ge i}\frac{\UnsignedStirlingFirstKind ni}{n\,n!}
 =\frac1{i!}\int_0^1\frac{[-\log(1-x)]^i}{x}\Differential x.
\]
With $u=-\log(1-x)$ this becomes
$\frac1{i!}\int_0^\infty u^i/(\EulerE^u-1)\Differential u=\RiemannZeta(i+1)$.
For the Hurwitz form use
$1/\RisingFactorial v n=\Gamma(v)/\Gamma(v+n)$ and the beta integral
\[
 \frac1{\RisingFactorial v n}=\frac1{(n-1)!}\int_0^1t^{v-1}(1-t)^{n-1}\Differential t,
\]
then sum \cref{eq:unsigned-first-egf} under the integral and use the standard
Mellin integral for $\HurwitzZeta{i+1}{v}$.
\end{proof}

A repeated partial-fraction expansion of $x^{-k}$ around $1-x$ gives, for
$k\ge3$ and $\Re z>k-1$,
\begin{multline}\label{eq:log-power-integral}
 \int_0^1\frac{\log^z(1-x)}{x^k}\Differential x
 =\frac{(-1)^z\Gamma(z+1)}{(k-1)!}
 \sum_{r=1}^{k-1}\SignedStirlingFirstKind{k-1}{r}\\
 {}\times\sum_{m=0}^r\binom rm(k-2)^{r-m}\RiemannZeta(z+1-m),
\end{multline}
with the principal interpretation of $(-1)^z$ when $z$ is not integral.
\begin{proof}
Substitute $u=-\log(1-x)$, expand
$(1-\EulerE^{-u})^{-k}$ after $k-1$ integrations by parts, and express the resulting
polynomial in the Euler operator $uD$ in the ordinary-power basis by
$\SignedStirlingFirstKind{k-1}{r}$.  Each term is a Mellin integral
$\int_0^\infty u^{z-m}/(\EulerE^u-1)\Differential u=\Gamma(z+1-m)\RiemannZeta(z+1-m)$; simplifying gamma
ratios yields \cref{eq:log-power-integral}.
\end{proof}

The more elaborate Hurwitz expansion in the archive,
\begin{multline}\label{eq:hurwitz-finite-difference-expansion}
 \HurwitzZeta{s}{v}=\frac{k!}{\FallingFactorial{s-k}{k}}
 \sum_{n\ge0}\frac{\UnsignedStirlingFirstKind{n+k}{n}}{(n+k)!}\\
 {}\times\sum_{\ell=0}^{n+k-1}(-1)^\ell\binom{n+k-1}{\ell}
 (\ell+v)^{k-s},
\end{multline}
for $k\ge1$ (initially in the common domain of absolute convergence, then by
analytic continuation), follows by applying \cref{eq:der-from-diff} to the
$k$-fold antiderivative of $(x+v)^{-s}$ and summing over $x\in\NaturalNumbers$.  The inner sum is
a finite difference, while the outer coefficient is the first-kind conversion
coefficient.  This derivation also supplies the convergence conditions suppressed in
the compact archived display.


\section{Convergent Stirling series for \texorpdfstring{$\log\Gamma$}{log Gamma}, \texorpdfstring{$\psi$}{psi}, and the Stieltjes constants}
The source archive lists three nested Stirling series without their convergence
mechanism.  All three arise from the same substitution
\begin{equation}
\label{eq:logarithmic-binet-substitution}
 u=1-\EulerE^{-2\pi x},
 \qquad
 x=\frac{-\log(1-u)}{2\pi},
 \qquad
 \frac{\Differential x}{\EulerE^{2\pi x}-1}=\frac{\Differential u}{2\pi u}.
\end{equation}
We first record the coefficient estimate that justifies integration at the endpoint
$u=1$.

\begin{lemma}[Logarithmic transfer estimate]
\label{lem:logarithmic-transfer}
Suppose $h$ is analytic in the unit disk, continues to a dented neighborhood of
$u=1$, and, there,
\[
 h(u)=h_0+O\!\left(
 \frac{(1+\log\log\frac1{|1-u|})^M}
      {\log\frac1{|1-u|}}
 \right)
\]
for some fixed $M\ge0$.  If $h(u)=\sum_{n\ge0}h_nu^n$, then
\begin{equation}
\label{eq:logarithmic-transfer}
 h_n=O\!\left(
 \frac{(1+\log\log n)^M}{n\log^2 n}
 \right).
\end{equation}
Consequently $\sum_{n\ge1}|h_n|/n$ converges.
\end{lemma}
\begin{proof}
Apply Cauchy's coefficient formula on a keyhole contour that follows the two sides
of the cut $[1,\infty)$.  The part separated from $1$ is exponentially small.
On the local part put $u=1-w/n$.  The constant $h_0$ has no positive-index
coefficients, while the jump of
$1/\log(1/(1-u))$ across the cut is
$O(\log^{-2}n)$; multiplication by a power of $\log\log(1/(1-u))$ contributes at
most $(1+\log\log n)^M$.  The Cauchy kernel contributes $n^{-1}\EulerE^w$, and the two
rays of the keyhole have an integrable exponential majorant.  This gives
\eqref{eq:logarithmic-transfer}.  Division by another factor $n$ makes the
resulting majorant summable.
\end{proof}

\subsection{The logarithm and logarithmic derivative of the gamma function}
Binet's second formula, valid for $z>0$, is
\begin{equation}
\label{eq:binet-second}
 \log\Gamma(z)
 =\left(z-\frac12\right)\log z-z+\frac12\log(2\pi)
 +2\int_0^\infty\frac{\arctan(x/z)}{\EulerE^{2\pi x}-1}\,\Differential x.
\end{equation}
It follows, for example, by applying the Abel--Plana summation formula to
$w\mapsto\log(z+w)$, subtracting the elementary integral and endpoint terms, and
normalizing at $z=1$ by the Wallis product.  Differentiation under the integral
recovers the usual Binet integral for the digamma function, so the normalization
also proves \eqref{eq:binet-second} rather than merely determining it up to a
constant.

\begin{theorem}[Convergent Stirling series for $\log\Gamma$ and $\psi$]
\label{thm:gamma-stirling-series}
For real $z>1/6$,
\begin{multline}
\label{eq:gamma-stirling-series}
 \log\Gamma(z)
 =\left(z-\frac12\right)\log z-z+\frac12\log(2\pi)\\
 {}+\frac1\pi\sum_{n=1}^{\infty}\frac1{n\,n!}
 \sum_{\ell=0}^{\Floor{n/2}}
 \frac{(-1)^\ell(2\ell)!}{(2\pi z)^{2\ell+1}}
 \UnsignedStirlingFirstKind n{2\ell+1}.
\end{multline}
The series is locally uniformly convergent in $z>1/6$, and termwise
differentiation gives
\begin{multline}
\label{eq:digamma-stirling-series}
 \psi(z)=\frac{\Gamma'(z)}{\Gamma(z)}
 =\log z-\frac1{2z}\\
 {}-\frac1{\pi z}\sum_{n=1}^{\infty}\frac1{n\,n!}
 \sum_{\ell=0}^{\Floor{n/2}}
 \frac{(-1)^\ell(2\ell+1)!}{(2\pi z)^{2\ell+1}}
 \UnsignedStirlingFirstKind n{2\ell+1}.
\end{multline}
\end{theorem}
\begin{proof}
Put $L(u)=-\log(1-u)$.  The formal Taylor series of the composite function is
\begin{align}
 \arctan\frac{L(u)}{2\pi z}
 &=\sum_{\ell\ge0}\frac{(-1)^\ell}{2\ell+1}
   \frac{L(u)^{2\ell+1}}{(2\pi z)^{2\ell+1}}\notag\\
 &=\sum_{n\ge1}\frac{u^n}{n!}
   \sum_{\ell=0}^{\Floor{n/2}}
   \frac{(-1)^\ell(2\ell)!}{(2\pi z)^{2\ell+1}}
   \UnsignedStirlingFirstKind n{2\ell+1},
 \label{eq:arctan-log-stirling-expansion}
\end{align}
where the second line uses
$L(u)^r/r!=\sum_{n\ge r}\UnsignedStirlingFirstKind nr u^n/n!$.
For $z>1/6$, the analytic continuation of this composite from $u=0$ has no
singularity in $|u|<1$.  Indeed, the possible arctangent branch points would
require $\log(1-u)=\pm2\pi iz$; when $1/6<z<1/4$ their solutions have modulus
$2\sin(\pi z)>1$, and for $z\ge1/4$ the principal logarithm on the disk cannot
have imaginary part $\pm2\pi z$.  The only boundary singularity relevant to the
coefficient estimate is therefore $u=1$.

As $u\to1$ in a dented domain,
\[
 \arctan\frac{L(u)}{2\pi z}
 =\frac\pi2+O\!\left(\frac1{L(u)}\right).
\]
Thus \cref{lem:logarithmic-transfer} applies, uniformly when $z$ ranges over a
compact subinterval of $(1/6,\infty)$.  We may divide the Taylor coefficients by
$n$ and sum at $u=1$.  Substituting \eqref{eq:logarithmic-binet-substitution} in
\eqref{eq:binet-second} now gives
\[
 2\int_0^\infty\frac{\arctan(x/z)}{\EulerE^{2\pi x}-1}\Differential x
 =\frac1\pi\int_0^1\arctan\frac{L(u)}{2\pi z}\frac{\Differential u}{u},
\]
and termwise integration of
\eqref{eq:arctan-log-stirling-expansion} proves
\eqref{eq:gamma-stirling-series}.  The same uniform estimate after one
$z$-derivative justifies termwise differentiation; differentiating
$(2\pi z)^{-2\ell-1}$ produces the factor $-(2\ell+1)/z$ and yields
\eqref{eq:digamma-stirling-series}.
\end{proof}

For complex $z$ with $\Re z>0$, the same proof works whenever the two points
$1-\EulerE^{2\pi iz}$ and $1-\EulerE^{-2\pi iz}$ lie outside the closed unit disk; the stated
real half-line is the simplest connected domain containing all commonly used
positive arguments.

\subsection{Stieltjes constants}
Define the Stieltjes constants by the Laurent expansion
\begin{equation}
\label{eq:stieltjes-definition}
 \RiemannZeta(s)=\frac1{s-1}+\sum_{m=0}^{\infty}
 \frac{(-1)^m\gamma_m}{m!}(s-1)^m.
\end{equation}

\begin{theorem}[Stirling series for the Stieltjes constants]
\label{thm:stieltjes-stirling-series}
For every integer $m\ge0$,
\begin{equation}
\label{eq:stieltjes-stirling-series}
\begin{aligned}
 \gamma_m={}&\frac12\delta_{m0}
 +\frac{(-1)^m m!}{\pi}
 \sum_{n=1}^{\infty}\frac1{n\,n!}\\
 &\times\sum_{k=0}^{\Floor{n/2}}
 \frac{(-1)^k}{(2\pi)^{2k+1}}
 \UnsignedStirlingFirstKind{2k+2}{m+1}\UnsignedStirlingFirstKind n{2k+1}.
\end{aligned}
\end{equation}
The outer series converges absolutely.
\end{theorem}
\begin{proof}
The Abel--Plana formula applied to $(1+w)^{-s}$ gives Hermite's representation
\begin{equation}
\label{eq:hermite-zeta}
 \RiemannZeta(s)=\frac1{s-1}+\frac12+\frac1i\int_0^\infty
 \frac{(1-ix)^{-s}-(1+ix)^{-s}}{\EulerE^{2\pi x}-1}\,\Differential x,
 \qquad \Re s>0,
\end{equation}
with the pole omitted at $s=1$.  Expand the two powers at $s=1$ and compare with
\eqref{eq:stieltjes-definition}.  Differentiation under the exponentially
convergent integral yields
\begin{equation}
\label{eq:jensen-franel-stieltjes}
 \gamma_m=\frac12\delta_{m0}+\frac1i\int_0^\infty\frac1{\EulerE^{2\pi x}-1}
 \left\{
 \frac{\log^m(1-ix)}{1-ix}
 -\frac{\log^m(1+ix)}{1+ix}
 \right\}\Differential x.
\end{equation}

Differentiate the signed first-kind generating function to obtain, for $|w|<1$,
\begin{equation}
\label{eq:log-power-over-linear}
 \frac{\log^m(1+w)}{1+w}
 =m!\sum_{r\ge m}\SignedStirlingFirstKind{r+1}{m+1}\frac{w^r}{r!}.
\end{equation}
Only odd powers survive after replacing $w$ by $-ix$ and $ix$ and taking the
difference.  Since
$\SignedStirlingFirstKind{2k+2}{m+1}=(-1)^{m+1}\UnsignedStirlingFirstKind{2k+2}{m+1}$, one finds
\begin{multline}
\label{eq:stieltjes-odd-expansion}
 \frac1i\left\{
 \frac{\log^m(1-ix)}{1-ix}
 -\frac{\log^m(1+ix)}{1+ix}
 \right\}\\
 =2(-1)^m m!\sum_{k\ge0}
 \frac{(-1)^k\UnsignedStirlingFirstKind{2k+2}{m+1}}{(2k+1)!}\,x^{2k+1}.
\end{multline}
Now make the substitution \eqref{eq:logarithmic-binet-substitution}.  For every
$k$,
\[
 \frac{L(u)^{2k+1}}{(2k+1)!}
 =\sum_{n\ge2k+1}\UnsignedStirlingFirstKind n{2k+1}\frac{u^n}{n!}.
\]
The composite integrand in \eqref{eq:jensen-franel-stieltjes} is analytic in the
unit disk and, as $u\to1$, is
\[
 O\!\left(\frac{(1+\log L(u))^m}{L(u)}\right).
\]
Hence \cref{lem:logarithmic-transfer} permits absolute termwise integration after
division by $u$.  The Jacobian contributes $1/(2\pi)$, the factor $2$ in
\eqref{eq:stieltjes-odd-expansion} cancels it to $1/\pi$, and
$\int_0^1u^{n-1}\Differential u=1/n$.  Collecting the finite contributions to each
coefficient $u^n$ gives exactly \eqref{eq:stieltjes-stirling-series}.
\end{proof}

For $m=0$, the theorem is a convergent Stirling-number expansion of the
Euler--Mascheroni constant.  The factors
$\UnsignedStirlingFirstKind{2k+2}{m+1}$ show that, at every fixed outer index $n$, only finitely many
orders of logarithmic differentiation contribute; this is the same triangularity
that underlies Fa\`a di Bruno's formula and series reversion.
\section{Gregory coefficients and derivatives of logarithmic powers}
The Gregory coefficients $\GregoryCoefficient{n}$ are defined by
\[
 \frac{z}{\log(1+z)}=1+\sum_{n\ge1}\GregoryCoefficient{n}z^n.
\]
Multiplying by the first-kind expansion of powers of $\log(1+z)$ gives
\begin{equation}\label{eq:gregory-stirling}
 n!\GregoryCoefficient{n}=\sum_{\ell=0}^n\frac{\SignedStirlingFirstKind{n}{\ell}}{\ell+1}.
\end{equation}
Indeed,
$z/\log(1+z)=\int_0^1(1+z)^t\Differential t$ and comparison with
$(1+z)^t=\sum_n\sum_\ell \SignedStirlingFirstKind{n}{\ell}t^\ell z^n/n!$ proves the claim.

For arbitrary $\mu$ and $n\ge1$,
\begin{equation}\label{eq:log-power-derivative}
 \frac{\Differential^n}{\Differential x^n}(\log x)^\mu
 =x^{-n}\sum_{k=1}^n\FallingFactorial\mu k\,\SignedStirlingFirstKind{n}{n-k+1}(\log x)^{\mu-k}.
\end{equation}
\begin{proof}
Write $D_x=x^{-1}D_t$ with $t=\log x$.  Then
$D_x^n=x^{-n}\FallingFactorial{D_t}{n}$, as follows inductively from
$x^{-1}D_t\,x^{-m}=x^{-m-1}(D_t-m)$.  Expand
$\FallingFactorial{D_t}{n}$ by signed first-kind numbers and apply powers of $D_t$ to
$t^\mu$; reindexing gives the displayed form.
\end{proof}

\section{Generalized factorial coefficients}
Given $f:\NaturalNumbers\to\ComplexNumbers$ and $t\ne0$, define
\begin{equation}\label{eq:generalized-factorial-product}
 (x)_{n,f,t}:=\prod_{j=1}^{n-1}\left(x+\frac{f(j)}{t^j}\right),
 \qquad
 \UnsignedStirlingFirstKind nk_{f,t}:=\CoefficientExtraction{x}{k-1}(x)_{n,f,t}.
\end{equation}
Multiplication by the last factor immediately yields
\begin{equation}\label{eq:generalized-first-recurrence}
 \UnsignedStirlingFirstKind nk_{f,t}
 =f(n-1)t^{1-n}\UnsignedStirlingFirstKind{n-1}{k}_{f,t}
  +\UnsignedStirlingFirstKind{n-1}{k-1}_{f,t},
\end{equation}
with the usual delta boundary convention.  The generalized harmonic sums
$F_n^{(r)}(t)=\sum_{j\le n}t^j/f(j)^r$ enter the coefficients through the same
Newton/Bell-polynomial calculation as \cref{eq:first-harmonic-bell}.

More broadly, any triangular array satisfying
\begin{multline}\label{eq:triangular-general-recurrence}
 a(n,k)=(\alpha n+\beta k+\gamma)a(n-1,k)\\
 +(\alpha'n+\beta'k+\gamma')a(n-1,k-1)+\delta_{n0}\delta_{k0}
\end{multline}
is governed by a first-order differential equation for its bivariate generating
function.  This umbrella contains both Stirling kinds, Lah numbers, Eulerian
variants, and many generalized factorial coefficients.



\chapter{Restricted variants of set partitions}
\label{ch:stirling-variants}
The second-kind recurrence survives many natural restrictions, but its boundary
term changes according to how the block containing the newest label is allowed to
look.

\section{\texorpdfstring{$r$}{r}-Stirling numbers of the second kind}
For fixed $r\ge0$, let
\[
 \StirlingSecondKind nk_{\!r}
\]
denote the number of partitions of $[n]$ into $k$ nonempty blocks in which the
labels $1,\ldots,r$ lie in distinct blocks.  The definition is meaningful for
$n\ge r$ and forces $k\ge r$.

\begin{theorem}[$r$-Stirling recurrence]
\label{thm:r-stirling-recurrence}
For $n>r$,
\begin{equation}
\label{eq:r-stirling-recurrence}
 \StirlingSecondKind nk_{\!r}
 =k\StirlingSecondKind{n-1}{k}_{\!r}+\StirlingSecondKind{n-1}{k-1}_{\!r},
\end{equation}
with $\StirlingSecondKind rr_{\!r}=1$ and the natural zero boundary conditions.
\end{theorem}
\begin{proof}
Because $n>r$, the new label $n$ is not one of the distinguished labels.  It either
joins one of the $k$ blocks of an admissible partition of $[n-1]$, or forms a
singleton beside an admissible partition into $k-1$ blocks.  The distinction of
$1,\ldots,r$ is preserved in both constructions, and deletion of $n$ reverses
them.
\end{proof}

The same labelled-component argument gives the mixed exponential generating
function
\begin{equation}
\label{eq:r-stirling-egf}
 \sum_{n\ge r}\sum_{k\ge r}\StirlingSecondKind nk_{\!r}
 y^k\frac{z^{n-r}}{(n-r)!}
 =y^re^{rz}\exp\!\bigl(y(\EulerE^z-1)\bigr).
\end{equation}
Indeed, the $r$ distinguished blocks each consist of one distinguished atom and an
arbitrary set of ordinary atoms, contributing $y^re^{rz}$, while the remaining
blocks form an ordinary set partition.

\section{Associated Stirling numbers}
Let $\mathsf S_r(n,k)$ count partitions of $[n]$ into $k$ blocks, every block having
size at least $r$.

\begin{theorem}[Associated recurrence]
\label{thm:associated-stirling-recurrence}
For $r\ge1$,
\begin{equation}
\label{eq:associated-stirling-recurrence}
 \mathsf S_r(n+1,k)
 =k\mathsf S_r(n,k)+\binom n{r-1}\mathsf S_r(n-r+1,k-1).
\end{equation}
Moreover,
\begin{equation}
\label{eq:associated-stirling-egf}
 \sum_{n\ge0}\mathsf S_r(n,k)\frac{z^n}{n!}
 =\frac1{k!}\left(\EulerE^z-\sum_{j=0}^{r-1}\frac{z^j}{j!}\right)^k.
\end{equation}
\end{theorem}
\begin{proof}
Consider the block containing $n+1$.  If deleting $n+1$ leaves a block of size at
least $r$, start from a partition counted by $\mathsf S_r(n,k)$ and choose one of
its $k$ blocks for insertion.  Otherwise the block containing $n+1$ has exactly
$r$ elements: choose its $r-1$ companions and partition the remaining $n-r+1$
labels into $k-1$ admissible blocks.  These cases are disjoint and exhaustive,
proving the recurrence.  For the generating function, one admissible block is a
labelled set of size at least $r$, with EGF
$\EulerE^z-\sum_{j<r}z^j/j!$; an unordered set of $k$ such blocks contributes its $k$th
power divided by $k!$.
\end{proof}

For $r=2$ these are often called Ward numbers.  Their occurrence as coefficients
of other polynomial families comes from the same exclusion of singleton
components expressed by \eqref{eq:associated-stirling-egf}.

\section{Reduced Stirling numbers and graph coloring}
For $d\ge1$, let $\mathsf S^{[d]}(n,k)$ count partitions of $[n]$ into $k$ blocks
such that two labels in the same block always differ by at least $d$.

\begin{theorem}[Reduction formula]
\label{thm:reduced-stirling}
For $n\ge k\ge d$,
\begin{equation}
\label{eq:reduced-stirling}
 \mathsf S^{[d]}(n,k)=\StirlingSecondKind{n-d+1}{k-d+1}.
\end{equation}
In particular $\mathsf S^{[1]}(n,k)=\StirlingSecondKind nk$.
\end{theorem}
\begin{proof}
Join two vertices $i,j\in[n]$ when $0<|i-j|<d$, and call the resulting graph
$G_{n,d}$.  Its independent-set partitions are exactly the partitions counted by
$\mathsf S^{[d]}(n,k)$.  A proper coloring can therefore be obtained by choosing
such a partition and assigning distinct colors to its blocks, so
\begin{equation}
\label{eq:reduced-chromatic-falling}
 \chi_{G_{n,d}}(q)=\sum_k\mathsf S^{[d]}(n,k)\FallingFactorial qk.
\end{equation}
Color the vertices in their natural order.  The first $d-1$ vertices must receive
distinct colors; every subsequent vertex sees the preceding $d-1$ vertices, which
form a clique, and hence has $q-d+1$ choices.  Thus
\[
 \chi_{G_{n,d}}(q)
 =\FallingFactorial q{d-1}(q-d+1)^{n-d+1}.
\]
Apply the ordinary Stirling basis expansion to the last power:
\begin{align*}
 \chi_{G_{n,d}}(q)
 &=\FallingFactorial q{d-1}
   \sum_{r\ge0}\StirlingSecondKind{n-d+1}{r}\FallingFactorial{q-d+1}{r}\\
 &=\sum_{r\ge0}\StirlingSecondKind{n-d+1}{r}\FallingFactorial q{r+d-1}.
\end{align*}
Comparison with \eqref{eq:reduced-chromatic-falling} gives
$k=r+d-1$ and proves \eqref{eq:reduced-stirling}.
\end{proof}
\part{Bell Numbers and Set Partitions}

\chapter{Set partitions and the many meanings of the Bell numbers}
\index{Bell numbers}

\section{Definition, initial values, and equivalence relations}
\begin{definition}
For $n\ge0$, the \emph{$n$th Bell number} $\BellNumber n$ is the number of partitions of
an $n$-element labelled set into nonempty, pairwise disjoint blocks.  Thus
\[
 \BellNumber0=\BellNumber1=1,
 \qquad
 (\BellNumber n)_{n\ge0}=1,1,2,5,15,52,203,877,4140,\ldots .
\]
The empty set has one partition, namely the empty family of blocks.
\end{definition}

For example, the five partitions of $\{a,b,c\}$ are
\[
 abc,\qquad a\mid bc,\qquad b\mid ac,\qquad c\mid ab,\qquad a\mid b\mid c,
\]
where vertical bars separate blocks.

\begin{proposition}
Partitions of a set $S$ are in natural bijection with equivalence relations on $S$.
Consequently $\BellNumber n$ also counts equivalence relations on an $n$-element set.
\end{proposition}
\begin{proof}
Given a partition $\pi$, declare $x\sim_\pi y$ when $x$ and $y$ lie in the same
block.  Reflexivity, symmetry, and transitivity follow immediately.  Conversely,
the equivalence classes of an equivalence relation are nonempty, disjoint, cover
$S$, and hence form a partition.  The two constructions undo one another.
\end{proof}


\section{Ordered set partitions and the Fubini numbers}
If the blocks themselves are linearly ordered, the resulting numbers are usually
called the \emph{ordered Bell numbers} or \emph{Fubini numbers}.  To prevent a
collision with the Bell notation, write them as $\mathsf F_n$.

\begin{theorem}[Ordered Bell numbers]
\label{thm:ordered-bell}
For $n\ge0$,
\begin{align}
 \mathsf F_n&=\sum_{k=0}^{n}k!\StirlingSecondKind nk,
 \label{eq:ordered-bell-stirling}\\
 \sum_{n\ge0}\mathsf F_n\frac{z^n}{n!}
 &=\frac{1}{2-\EulerE^z},
 \label{eq:ordered-bell-egf}\\
 \mathsf F_n&=\sum_{j=1}^{n}\binom nj\mathsf F_{n-j}
 \qquad(n\ge1),
 \label{eq:ordered-bell-recurrence}
\end{align}
with $\mathsf F_0=1$.
\end{theorem}
\begin{proof}
A partition into $k$ blocks may be ordered in $k!$ ways, proving
\eqref{eq:ordered-bell-stirling}.  By the second-kind exponential generating
function,
\[
 \sum_{n\ge0}\mathsf F_n\frac{z^n}{n!}
 =\sum_{k\ge0}k!\frac{(\EulerE^z-1)^k}{k!}
 =\sum_{k\ge0}(\EulerE^z-1)^k
 =\frac1{2-\EulerE^z},
\]
first as a formal identity and then analytically near the origin.  Finally, in a
nonempty ordered partition choose the elements of its first block, which may be any
nonempty subset, and order-partition the complement.  Summing over the first-block
size gives \eqref{eq:ordered-bell-recurrence}.
\end{proof}

Thus
\[
 (\mathsf F_n)_{n\ge0}=1,1,3,13,75,541,4683,\ldots,
\]
which contrasts with the unordered Bell sequence by recording the order of the
blocks rather than only their underlying family.
\section{Squarefree factorizations, rhyme schemes, and restricted-growth words}
\begin{proposition}
If $N=p_1\cdots p_n$ is squarefree, then $\BellNumber n$ is the number of unordered
factorizations of $N$ into integers greater than $1$.
\end{proposition}
\begin{proof}
To a set partition $\pi$ of $[n]$ associate the multiset of factors
$\{\prod_{i\in B}p_i:B\in\pi\}$.  Unique factorization recovers each block from
its factor, so this is a bijection.  For $30=2\cdot3\cdot5$ the five
factorizations are
\[
 30,\quad 2\cdot15,\quad3\cdot10,\quad5\cdot6,\quad2\cdot3\cdot5.
\]
\end{proof}

A rhyme scheme of an $n$-line stanza is likewise a partition of the set of lines:
two indices lie in the same block exactly when the corresponding lines rhyme.
The customary lettering convention is the following canonical encoding.

\begin{definition}
A \emph{restricted-growth word} of length $n$ is a word
$a_1\cdots a_n$ of nonnegative integers such that $a_1=0$ and
\[
 a_i\le 1+\max(a_1,\ldots,a_{i-1})\qquad(i>1).
\]
\end{definition}

\begin{proposition}
Set partitions of $[n]$, restricted-growth words of length $n$, and canonical
rhyme schemes of length $n$ are mutually bijective.
\end{proposition}
\begin{proof}
Order the blocks by increasing least element, number them $0,1,\ldots$, and put
$a_i=j$ when $i$ belongs to block $j$.  The displayed restriction says precisely
that a new block receives the next unused number.  Conversely, equal letters form
the blocks.  Replacing $0,1,2,\ldots$ by $A,B,C,\ldots$ gives the rhyme scheme.
\end{proof}

\section{A card-shuffling bijection}
Consider the operation that removes the top card from a deck of $n$ labelled cards
and reinserts it in any of the $n$ positions.  There are $n^n$ sequences of $n$
such operations.

\begin{theorem}[Return-to-order shuffles]
Exactly $\BellNumber n$ of those $n^n$ operation sequences return an initially ordered
deck to its original order.  Hence the return probability is $\BellNumber n/n^n$.
\end{theorem}
\begin{proof}
Reverse time.  A reversed operation chooses an arbitrary card and moves it to the
top.  Thus a reversed history is a word $c_1\cdots c_n$ in the card labels.  After
all moves, the cards that occurred in the word appear first, in decreasing order of
their last occurrence times; cards never chosen follow in their original relative
order.

Partition the time set $[n]$ according to equality of the letters $c_i$.  For the
final deck to be $1,2,\ldots,n$, the labels that occur must be an initial interval
$1,\ldots,m$, and their blocks of occurrence times must be labelled in decreasing
order of their maxima: the block with largest maximum receives label $1$, the next
receives $2$, and so on.  Therefore every set partition of $[n]$ determines exactly
one successful word, and every successful word determines its equality partition.
This is the required bijection.
\end{proof}

\section{Bell-enumerated vincular pattern classes}
A dash in a permutation pattern means that the corresponding entries need not be
adjacent; an undashed pair must be adjacent.  For instance, an occurrence of
$23\!\mathord{-}\!1$ consists of an adjacent ascent followed later by a smaller
entry.

\begin{theorem}
The permutations of $[n]$ avoiding $23\!\mathord{-}\!1$ are counted by $\BellNumber n$.
The same is true for each pattern in
\[
 1\!\mathord{-}\!23,\ 12\!\mathord{-}\!3,\ 32\!\mathord{-}\!1,
 \ 3\!\mathord{-}\!21,\ 1\!\mathord{-}\!32,\ 3\!\mathord{-}\!12,
 \ 21\!\mathord{-}\!3,
\]
and for the equivalent barred-pattern class in which every $321$ occurrence
extends to a $3241$ occurrence.
\end{theorem}
\begin{proof}
For the pattern $23\!\mathord{-}\!1$, order the blocks of a partition by increasing
least element and write the entries of each block in decreasing order.  Concatenate
the resulting words.  In the resulting permutation, each block is a decreasing
run, and every ascent occurs between two blocks.  Its lower endpoint is a
right-to-left minimum, so no later entry is smaller; hence $23\!\mathord{-}\!1$ is
avoided.

Conversely, cut a $23\!\mathord{-}\!1$-avoiding permutation immediately after
each right-to-left minimum.  Every resulting segment is decreasing: otherwise an
ascent inside a segment, together with its terminal right-to-left minimum, would
form $23\!\mathord{-}\!1$.  The segment minima increase from left to right, so the
segments, viewed as sets, are blocks already in canonical order.  The two maps are
inverse.

Reversal, complement, inverse, and their compositions are bijections on
permutations and carry vincular occurrences to occurrences of the seven listed
patterns; therefore all eight classes are equinumerous.  Finally, cutting at
right-to-left minima shows that failure of the decreasing-run condition is
simultaneously an occurrence of $23\!\mathord{-}\!1$ and a $321$ that cannot be
extended by an intervening entry to $3241$.  Thus the barred-pattern formulation
describes the same class.
\end{proof}

\begin{remark}
Classical avoidance classes have at most exponential growth in $n$, whereas Bell
numbers grow faster than $C^n$ for every fixed $C$.  The adjacency requirement is
therefore essential: these Bell-enumerated classes cannot be classical
single-pattern avoidance classes.
\end{remark}

\section{The Bell--Aitken triangle}
Define a triangular array $(a_{n,k})_{0\le k\le n}$ by
\[
 a_{0,0}=1,\qquad a_{n,0}=a_{n-1,n-1},\qquad
 a_{n,k}=a_{n,k-1}+a_{n-1,k-1}\quad(1\le k\le n).
\]
Its first rows are
\[
\begin{array}{rrrrr}
1\\
1&2\\
2&3&5\\
5&7&10&15\\
15&20&27&37&52.
\end{array}
\]

\begin{theorem}
For $0\le k\le n$,
\begin{equation}\label{eq:bell-triangle-closed}
 a_{n,k}=\sum_{j=0}^{k}\binom{k}{j}\BellNumber{n-j}.
\end{equation}
Consequently the left edge is $a_{n,0}=\BellNumber n$ and the right edge is
$a_{n,n}=\BellNumber{n+1}$.
\end{theorem}
\begin{proof}
The formula is immediate for $k=0$.  Pascal's identity gives
\[
 a_{n,k-1}+a_{n-1,k-1}
 =\sum_j\left(\binom{k-1}{j}+\binom{k-1}{j-1}\right)\BellNumber{n-j},
\]
which is \eqref{eq:bell-triangle-closed}.  For $k=n$, reverse the index and use
the Bell recurrence proved in \cref{thm:bell-binomial-recurrence} below.
\end{proof}

\chapter{The enumerative calculus of Bell numbers}

\section{Two fundamental sums}
\begin{theorem}[Binomial recurrence]\label{thm:bell-binomial-recurrence}
For $n\ge0$,
\begin{equation}\label{eq:bell-binomial-recurrence}
 \BellNumber{n+1}=\sum_{k=0}^{n}\binom nk\BellNumber k.
\end{equation}
\end{theorem}
\begin{proof}
In a partition of $\{0,1,\ldots,n\}$, let $k$ be the number of elements outside
the block containing $0$.  Choose those $k$ elements and partition them
arbitrarily; the block containing $0$ is then forced.
\end{proof}

\begin{theorem}[Stirling decomposition]
\begin{equation}\label{eq:bell-stirling-sum}
 \BellNumber n=\sum_{k=0}^{n}\StirlingSecondKind nk
 =\sum_{k=0}^{n}\frac1{k!}\sum_{i=0}^{k}(-1)^{k-i}\binom ki i^n.
\end{equation}
\end{theorem}
\begin{proof}
The first equality groups partitions by their number of blocks.  Substitute the
inclusion--exclusion formula for $\StirlingSecondKind nk$ from
\cref{eq:second-explicit} to obtain the second equality.
\end{proof}

\section{Spivey's identity and binomial inversions}
\begin{theorem}[Spivey's identity]\label{thm:spivey}
For $m,n\ge0$,
\begin{equation}\label{eq:spivey}
 \BellNumber{m+n}=\sum_{k=0}^{n}\sum_{j=0}^{m}
 \StirlingSecondKind mj\binom nk j^{\,n-k}\BellNumber k.
\end{equation}
\end{theorem}
\begin{proof}
Partition the first $m$ labels into $j$ blocks.  Of the final $n$ labels, choose
$k$ whose blocks contain no old label and partition them in $\BellNumber k$ ways.  Each
of the remaining $n-k$ labels independently joins one of the $j$ old blocks.  This
classification is unique and gives the summand.
\end{proof}

Binomial inversion of \eqref{eq:bell-binomial-recurrence} gives the first identity
below.  The second is a useful shifted form.
\begin{theorem}[Bell inversion identities]\label{thm:bell-inversions}
For $n,k\ge0$,
\begin{align}
 \BellNumber n&=\sum_{j=0}^{n}(-1)^{n-j}\binom nj\BellNumber{j+1},
 \label{eq:bell-inversion-one}\\
 \sum_{j=0}^{n}\binom nj\BellNumber{k+j}
 &=\sum_{i=0}^{k}(-1)^{k-i}\binom ki\BellNumber{n+i+1}.
 \label{eq:bell-inversion-two}
\end{align}
\end{theorem}
\begin{proof}
The first formula is ordinary binomial inversion.  For the second let $X$ be a
Poisson random variable of mean $1$; \cref{thm:dobinski} below gives
$\Expectation X^r=\BellNumber r$.  The left side is
$\Expectation[X^k(X+1)^n]$.  Poisson summation by parts,
\begin{equation}\label{eq:poisson-sbp}
 \Expectation[Xh(X)]=\Expectation[h(X+1)],
\end{equation}
iterated $k$ times in the equivalent form
$\Expectation[X^{n+1}(X-1)^k]=\Expectation[(X+1)^nX^k]$, turns it into the right side after expanding
$(X-1)^k$.
\end{proof}

The first-kind Stirling numbers package all repeated shifts at once.
\begin{theorem}[Weighted Bell shift]\label{thm:weighted-bell-shift}
For $a,b\in\ComplexNumbers$ and $k,n\ge0$,
\begin{multline}\label{eq:weighted-bell-shift}
 \sum_{j=0}^{n}\binom nj a^j b^{n-j}\BellNumber j\\
 =\sum_{i=0}^{k}(-1)^{k-i}\UnsignedStirlingFirstKind ki
   \sum_{j=0}^{n}\binom nj a^j(b-ak)^{n-j}\BellNumber{j+i}.
\end{multline}
In particular,
\begin{align}
 \sum_{j=0}^{n}\binom nj a^j b^{n-j}\BellNumber j
 &=\sum_{j=0}^{n}\binom nj a^j(b-a)^{n-j}\BellNumber{j+1},
 \label{eq:weighted-bell-k1}\\
 \sum_{j=0}^{n}\binom nj k^{n-j}\BellNumber j
 &=\sum_{i=0}^{k}(-1)^{k-i}\UnsignedStirlingFirstKind ki\BellNumber{n+i}.
 \label{eq:weighted-bell-special}
\end{align}
\end{theorem}
\begin{proof}
By Dobiński's formula, the left side of
\eqref{eq:weighted-bell-shift} is
\[
 \EulerE^{-1}\sum_{m\ge0}\frac{(am+b)^n}{m!}.
\]
On the right, use
$\sum_i(-1)^{k-i}\UnsignedStirlingFirstKind ki m^i=\FallingFactorial mk$ inside the same Poisson sum.  It
becomes
\[
 \EulerE^{-1}\sum_{m\ge0}\frac{\FallingFactorial mk\,[a m+b-ak]^n}{m!}
 =\EulerE^{-1}\sum_{r\ge0}\frac{(ar+b)^n}{r!},
\]
after $m=r+k$.  The special cases follow by setting $k=1$, and then $a=1,b=k$.
\end{proof}

\section{Exponential generating function and differential equation}
\begin{theorem}\label{thm:bell-egf}
The Bell exponential generating function is
\begin{equation}\label{eq:bell-egf}
 \mathscr B(z):=\sum_{n\ge0}\BellNumber n\frac{z^n}{n!}
 =\exp(\EulerE^z-1),
 \qquad
 \mathscr B'(z)=\EulerE^z\mathscr B(z).
\end{equation}
\end{theorem}
\begin{proof}
A set partition is a labelled set of nonempty labelled sets.  The EGF of a
nonempty set is $\EulerE^z-1$, and the exponential formula therefore gives
$\exp(\EulerE^z-1)$.  Differentiation gives the differential equation; conversely,
coefficient comparison in that equation gives
\eqref{eq:bell-binomial-recurrence} and $\mathscr B(0)=1$, so it uniquely determines
the series.
\end{proof}

\section{Dobiński's formula and Poisson moments}
\begin{theorem}[Dobiński]\label{thm:dobinski}
For every $n\ge0$,
\begin{equation}\label{eq:dobinski}
 \BellNumber n=\EulerE^{-1}\sum_{m=0}^{\infty}\frac{m^n}{m!}.
\end{equation}
Equivalently, if $X\sim\operatorname{Poisson}(1)$, then $\BellNumber n=\Expectation X^n$.
\end{theorem}
\begin{proof}
Expand the outer exponential in \eqref{eq:bell-egf}:
\[
 \EulerE^{\EulerE^z-1}=\EulerE^{-1}\sum_{m\ge0}\frac{\EulerE^{mz}}{m!}
 =\EulerE^{-1}\sum_{m\ge0}\frac1{m!}\sum_{n\ge0}m^n\frac{z^n}{n!}.
\]
Absolute convergence on compact sets justifies rearrangement analytically, while
formal coefficient extraction gives the same identity algebraically.  The last
statement is the defining moment sum of a Poisson$(1)$ variable.
\end{proof}

\begin{corollary}[Cauchy integral]
For every positively oriented circle $\gamma$ around $0$,
\begin{equation}\label{eq:bell-cauchy}
 \BellNumber n=\frac{n!}{2\pi \ImaginaryUnit \EulerE}\int_\gamma
       \frac{\EulerE^{\EulerE^z}}{z^{n+1}}\,\Differential z.
\end{equation}
\end{corollary}
\begin{proof}
Apply Cauchy's coefficient formula to $\EulerE^{-1}\EulerE^{\EulerE^z}$.
\end{proof}

\chapter{Congruences, shape, and growth of Bell numbers}

\section{Touchard congruences}
It is useful to refine $\BellNumber n$ to the Touchard polynomial
\[
 \TouchardPolynomial{n}(x)=\sum_{k=0}^{n}\StirlingSecondKind nk x^k,
 \qquad \TouchardPolynomial{n}(1)=\BellNumber n.
\]

\begin{theorem}[Polynomial Touchard congruence]\label{thm:touchard-poly}
For a prime $p$,
\begin{equation}\label{eq:touchard-poly}
 \TouchardPolynomial{n+p}(x)\equiv \TouchardPolynomial{n+1}(x)+x^p\TouchardPolynomial{n}(x)\pmod p.
\end{equation}
Consequently
\begin{equation}\label{eq:touchard}
 \BellNumber{n+p}\equiv\BellNumber n+\BellNumber{n+1}\pmod p.
\end{equation}
\end{theorem}
\begin{proof}
Let the cyclic group $C_p$ rotate a distinguished set of $p$ labels while fixing
$n$ other labels, and let it act on partitions, preserving the weight $x^{\#\text{
blocks}}$.  Every nonfixed orbit has size $p$ and contributes $0$ modulo $p$.
In a fixed partition there are two possibilities.  Either all $p$ rotating labels
form $p$ singleton blocks, independently of a partition of the other labels; this
contributes $x^p\TouchardPolynomial{n}(x)$.  Otherwise rotational invariance forces the rotating labels
to lie in one common block, possibly together with old labels.  Contracting those
$p$ labels to one distinguished label gives a partition of $n+1$ labels with the
same number of blocks, contributing $\TouchardPolynomial{n+1}(x)$.  This proves the polynomial
congruence; set $x=1$.
\end{proof}

\begin{theorem}[Prime-power shift]\label{thm:bell-prime-power-shift}
For a prime $p$ and $m\ge1$,
\begin{equation}\label{eq:bell-prime-power-shift}
 \BellNumber{n+p^m}\equiv m\BellNumber n+\BellNumber{n+1}\pmod p.
\end{equation}
More precisely,
\begin{equation}\label{eq:touchard-prime-power-poly}
 \TouchardPolynomial{n+p^m}(x)\equiv \TouchardPolynomial{n+1}(x)+\TouchardPolynomial{n}(x)\sum_{r=1}^{m}x^{p^r}
 \pmod p.
\end{equation}
\end{theorem}
\begin{proof}
Let the additive cyclic group of order $p^m$ act regularly on $p^m$ new labels.
A partition fixed by the whole group induces a block system for this regular
action.  Apart from the case in which all new labels join one invariant block
(which contracts to one label and contributes $\TouchardPolynomial{n+1}$), the blocks on the new
orbit are the cosets of the unique subgroup of index $p^r$ for some
$1\le r\le m$; they contribute $p^r$ blocks and cannot mix with old labels.
Thus the fixed weighted partitions contribute
$\TouchardPolynomial{n+1}+\TouchardPolynomial{n}\sum_{r=1}^m x^{p^r}$.  Every other orbit has cardinality divisible by
$p$, proving \eqref{eq:touchard-prime-power-poly}.  Fermat's congruence
$x^{p^r}\equiv x$ modulo $p$ at $x=1$ gives
\eqref{eq:bell-prime-power-shift}.
\end{proof}

\section{A universal period bound modulo a prime}
\begin{theorem}\label{thm:bell-period-bound}
Modulo a prime $p$, the Bell sequence is periodic, and its period divides
\begin{equation}\label{eq:bell-period-N}
 N_p=\frac{p^p-1}{p-1}=1+p+\cdots+p^{p-1}.
\end{equation}
\end{theorem}
\begin{proof}
Touchard's recurrence says that the shift operator $E$ on the sequence satisfies
$E^p-E-1=0$ over $\FiniteField_p$.  The polynomial $q(t)=t^p-t-1$ is separable because
$q'(t)=-1$.  If $\alpha$ is any root in a splitting field, then
$\alpha^p=\alpha+1$ and hence $\alpha^{p^j}=\alpha+j$ for $j\in\FiniteField_p$.  Taking the
product over $j$ gives
\[
 \alpha^{N_p}=\prod_{j\in\FiniteField_p}\alpha^{p^j}
 =\prod_{j\in\FiniteField_p}(\alpha+j)=\alpha^p-\alpha=1.
\]
Thus every eigenvalue of $E$ is an $N_p$th root of unity.  Separability makes $E$
diagonalizable over the splitting field, so $E^{N_p}=1$ on every solution of the
recurrence, in particular on the Bell sequence.
\end{proof}

\begin{remark}
A reported exact period is computational data, not a consequence of the divisibility
bound alone.  It can be certified by generating $N_p$ terms with
\eqref{eq:touchard}, finding the least divisor $d\mid N_p$ for which the state vector
of $p$ consecutive terms repeats, and checking that no proper divisor of $d$ works.
The same recurrence permits a compact machine-checkable certificate.
\end{remark}

\section{Log-convexity and normalized log-concavity}
\begin{theorem}\label{thm:bell-log-shape}
For $n\ge1$,
\begin{equation}\label{eq:bell-log-shape}
 1\le \frac{\BellNumber{n-1}\BellNumber{n+1}}{\BellNumber n^2}\le\frac{n+1}{n}.
\end{equation}
Thus $(\BellNumber n)$ is log-convex, whereas $(\BellNumber n/n!)$ is log-concave.
\end{theorem}
\begin{proof}
For the left inequality, let $X\sim\operatorname{Poisson}(1)$.  Cauchy--Schwarz
applied to $X^{(n-1)/2}$ and $X^{(n+1)/2}$ gives
$(\Expectation X^n)^2\le \Expectation X^{n-1}\Expectation X^{n+1}$.

For completeness we prove the normalized inequality through the following
convolution lemma.  If nonnegative $(z_j)_{j\ge1}$ is log-concave without internal
zeros and
\[
 A(t)=\exp\!\left(\sum_{j\ge1}\frac{z_j}{j}t^j\right)
      =\sum_{n\ge0}a_n\frac{t^n}{n!},
\]
then $a_n/n!$ is log-concave.  To see this, truncate at degree $N$, differentiate,
and compare the adjacent $2\times2$ Toeplitz minors in
$A'=A\sum_{j\ge1}z_jt^{j-1}$.  After cross multiplication, the $n$th minor is
\[
 \sum_{0\le i<j\le n}
 \frac{a_i a_j}{i!j!}
 \bigl(z_{n-i+1}z_{n-j}-z_{n-i}z_{n-j+1}\bigr),
\]
with the boundary terms interpreted as zero.  Each bracket is nonnegative by the
monotonicity of the ratios $z_{r+1}/z_r$; induction on $n$ starts at $a_0=1$.
Removing the truncation proves the lemma coefficientwise.

For $A(t)=\EulerE^{\EulerE^t-1}$ we have
$\sum_{j\ge1}t^j/j!=\sum_{j\ge1}z_jt^j/j$ with
$z_j=1/(j-1)!$.  This sequence is log-concave, since
$z_j^2/(z_{j-1}z_{j+1})=j/(j-1)$.  Hence
$(\BellNumber n/n!)^2\ge(\BellNumber{n-1}/(n-1)!)(\BellNumber{n+1}/(n+1)!)$, which is the right
inequality in \eqref{eq:bell-log-shape}.
\end{proof}

\section{Explicit upper bounds}
The positive-coefficient Cauchy estimate supplies a short route to the numerical
bounds in the archive.
\begin{lemma}\label{lem:bell-coeff-upper}
For $t>0$,
\begin{equation}\label{eq:bell-coeff-upper}
 \BellNumber n\le n!\,t^{-n}\exp(\EulerE^t-1).
\end{equation}
For $n\ge2$, put $L=\log(n+1)$ and $t=\log(n/L)$.  Then
\begin{equation}\label{eq:bell-ratio-upper}
 \BellNumber n^{1/n}\frac{L}{n}
 \le \frac{L}{t}
 \exp\!\left(-1+\frac1L+\frac{\log n}{2n}\right).
\end{equation}
\end{lemma}
\begin{proof}
Because all coefficients of $\mathscr B$ are nonnegative,
$\BellNumber n t^n/n!\le\mathscr B(t)$, proving the first claim.  Use
$\EulerE^t=n/L$ and the elementary Stirling bound
$n!\le e\sqrt n\,(n/e)^n$; the factors $\EulerE^{-1/n}$ and $\EulerE^{1/n}$ cancel and give
\eqref{eq:bell-ratio-upper}.
\end{proof}

\begin{theorem}[Uniform numerical bounds]\label{thm:bell-numerical-bounds}
For every $n\ge1$,
\begin{equation}\label{eq:bell-0792}
 \BellNumber n<\left(\frac{0.792\,n}{\log(n+1)}\right)^n.
\end{equation}
Define, for $x>1$,
\begin{equation}\label{eq:bell-dx}
 d(x)=\log\log(x+1)-\log\log x+\frac{1+\EulerE^{-1}}{\log x}.
\end{equation}
For every $\varepsilon>0$ there is
$n_0(\varepsilon)=\max\{\EulerE^4,d^{-1}(\varepsilon)\}$ such that, for
$n>n_0(\varepsilon)$,
\begin{equation}\label{eq:bell-e06}
 \BellNumber n<\left(\frac{\EulerE^{-0.6+\varepsilon}n}{\log(n+1)}\right)^n.
\end{equation}
\end{theorem}
\begin{proof}
Set
\[
 R(x)=\frac{\log(x+1)}{\log(x/\log(x+1))}
 \exp\!\left(-1+\frac1{\log(x+1)}+\frac{\log x}{2x}\right).
\]
Direct differentiation, after multiplying by the positive denominator
$2x\log x\log(x+1)\log(x/\log(x+1))^2$, shows that $R'(x)<0$ for $x\ge39$;
substitution gives $R(39)<0.792$.  Thus
\eqref{eq:bell-ratio-upper} proves \eqref{eq:bell-0792} for $n\ge39$.  The
remaining $1\le n\le38$ are obtained successively from
\eqref{eq:bell-binomial-recurrence}; substituting them gives strict inequalities.
This is a finite exact integer check.

A second differentiation gives, for $x\ge \EulerE^4$,
\[
 \log R(x)\le-0.6+d(x),\qquad d'(x)<0,\qquad \lim_{x\to\infty}d(x)=0.
\]
(The numerator after clearing positive denominators is a sum of
$-(\log x-4)^2$, $-(x-1-\log x)$, and negative multiples of these two standard
nonnegative quantities.)  Hence $d(n)<\varepsilon$ when
$n>d^{-1}(\varepsilon)$, and \eqref{eq:bell-e06} follows from
\eqref{eq:bell-ratio-upper}.
\end{proof}

\section{Saddle-point asymptotics and Lambert \texorpdfstring{$W$}{W}}
Let $r=W(n)$, so that $re^r=n$, and put $A=n(1+r)$.
\begin{theorem}[Leading Bell asymptotic]\label{thm:bell-leading-asymptotic}
As $n\to\infty$,
\begin{align}
 \BellNumber n
 &\sim \frac{n!\exp(\EulerE^r-1)}{r^n\sqrt{2\pi n(1+r)}}
 \label{eq:bell-saddle-leading}\\
 &\sim \frac1{\sqrt n}\left(\frac{n}{W(n)}\right)^{n+1/2}
       \exp\!\left(\frac{n}{W(n)}-n-1\right).
 \label{eq:bell-W-asymptotic}
\end{align}
\end{theorem}
\begin{proof}
In \eqref{eq:bell-cauchy} take the circle $z=re^{i\theta}$.  The coefficient
integral becomes
\[
 \frac{\BellNumber n}{n!}=
 \frac{\EulerE^{\EulerE^r-1}}{2\pi r^n}
 \int_{-\pi}^{\pi}
 \exp\!\left(\EulerE^{re^{i\theta}}-\EulerE^r-in\theta\right)\Differential\theta.
\]
The linear term vanishes because $re^r=n$; the quadratic term is
$-A\theta^2/2$.  On $|\theta|\le A^{-2/5}$, Taylor expansion is uniform and the
scaled integral tends to $\int_{\RealNumbers}\EulerE^{-u^2/2}\Differential u/\sqrt A$.  Outside that arc,
$\Re(\EulerE^{re^{i\theta}}-\EulerE^r)\le-cA\theta^2$ on the intermediate range and is
$\le-c'\EulerE^r$ on the remaining compact range, so the contribution is negligible.
This proves \eqref{eq:bell-saddle-leading}.  Stirling's formula and $\EulerE^r=n/r$
simplify it to \eqref{eq:bell-W-asymptotic}.
\end{proof}

The archived refined expansion used undefined symbols.  The following construction
makes every coefficient unambiguous.  Let
$\TouchardPolynomial{j}(r)=\sum_{q=0}^j\StirlingSecondKind jq r^q$ be the Touchard polynomial and let
$Z\sim N(0,1)$.  For an integer $h=O(r)$ define
\begin{equation}\label{eq:bell-saddle-correction-generator}
 \mathcal R_{r,h}(Z)=\exp\!\left(
 -\frac{ihZ}{\sqrt A}+
 \sum_{j\ge3}\frac{i^j \EulerE^r\TouchardPolynomial{j}(r)}{j!A^{j/2}}Z^j
 \right).
\end{equation}
Expand this formal exponential and take Gaussian moments termwise; retaining all
monomials of total order at most $\EulerE^{-Mr}$ defines $C_0(r,h),\ldots,C_M(r,h)$,
with $C_0=1$.

\begin{theorem}[Uniform Moser--Wyman scheme]\label{thm:bell-refined-saddle}
For every fixed $M$ and $|h|\le Cr$,
\begin{equation}\label{eq:bell-refined-saddle}
 \BellNumber{n+h}=\frac{(n+h)!\exp(\EulerE^r-1)}{r^{n+h}\sqrt{2\pi A}}
 \left(\sum_{q=0}^{M-1}C_q(r,h)+O_C(\EulerE^{-Mr}P_M(r))\right),
\end{equation}
where $P_M$ is an explicitly computable polynomial.  The first correction is
\begin{equation}\label{eq:bell-first-correction}
 C_1(r,h)=
 -\frac{h^2}{2A}-\frac{h \EulerE^r\TouchardPolynomial{3}(r)}{2A^2}
 +\frac{\EulerE^rT_4(r)}{8A^2}
 -\frac{5\EulerE^{2r}\TouchardPolynomial{3}(r)^2}{24A^3}.
\end{equation}
In particular the coefficient of $\EulerE^{-qr}$ is a polynomial in $h$ of degree at
most $2q$, with rational functions of $r$ as coefficients.
\end{theorem}
\begin{proof}
Use the same circle as in the preceding proof but replace $n$ by $n+h$.  After
$\theta=Z/\sqrt A$, the ratio of the integrand to its Gaussian quadratic part is
exactly the formal series \eqref{eq:bell-saddle-correction-generator}.  Watson's
lemma with a Gaussian majorant permits termwise integration on the central arc;
the complementary arcs have the same exponential bound uniformly for
$|h|\le Cr$.  Odd Gaussian moments vanish, so only integral powers of $\EulerE^{-r}$
remain.  Since $\Expectation Z^2=1$, $\Expectation Z^4=3$, and $\Expectation Z^6=15$, the terms from the square
of the $h$-term, its product with the cubic term, the quartic term, and the square
of the cubic term give \eqref{eq:bell-first-correction}.  The same finite
bookkeeping defines every subsequent $C_q$ and bounds the omitted Taylor
remainder.
\end{proof}

\begin{corollary}[de Bruijn logarithmic expansion]\label{cor:bell-debruijn}
As $n\to\infty$,
\begin{multline}\label{eq:bell-debruijn}
 \frac{\log\BellNumber n}{n}
 =\log n-\log\log n-1
 +\frac{\log\log n}{\log n}+\frac1{\log n}\\
 +\frac12\left(\frac{\log\log n}{\log n}\right)^2
 +O\!\left(\frac{\log\log n}{(\log n)^2}\right).
\end{multline}
\end{corollary}
\begin{proof}
Take logarithms in \eqref{eq:bell-W-asymptotic}.  The standard reversion of
$r+\log r=\log n$ gives
\[
 r=\log n-\log\log n+
rac{\log\log n}{\log n}
 +O\!\left(\frac{(\log\log n)^2}{(\log n)^2}\right).
\]
Substitution and one further Taylor expansion of $\log r$ produce the displayed
terms; the prefactor contributes only $O(\log n/n)$.
\end{proof}

\section{Bell primes and verifiable computation}
A \emph{Bell prime} is a prime value of $\BellNumber n$.  The archive records the early
indices found by computation.  Such a list is not a symbolic theorem, but each
entry has a rigorous certificate: compute $\BellNumber n$ by the integer recurrence
\eqref{eq:bell-binomial-recurrence}, apply a probable-prime filter, and attach an
ECPP or Pratt-style primality certificate.  Composite entries are certified by one
nontrivial factor.  This distinction between an identity and a finite certified
computation will also be used for tabulated values later.

\part{Bell Polynomials and Composition}

\chapter{Exponential and ordinary Bell polynomials}
\index{Bell polynomials}

\section{Definitions}
\begin{definition}[Partial exponential Bell polynomial]
For $0\le k\le n$, define
\begin{equation}\label{eq:partial-bell-definition}
 \ExponentialPartialBellPolynomial nk(x_1,\ldots,x_{n-k+1})
 =n!\!\sum_{\substack{j_1+j_2+\cdots=k\\
                       j_1+2j_2+3j_3+\cdots=n}}
 \prod_{i\ge1}\frac{x_i^{j_i}}{(i!)^{j_i}j_i!}.
\end{equation}
Only $i\le n-k+1$ can occur.  The complete exponential Bell polynomial is
\begin{equation}\label{eq:complete-bell-definition}
 \ExponentialCompleteBellPolynomial n(x_1,\ldots,x_n)=\sum_{k=0}^{n}\ExponentialPartialBellPolynomial nk(x_1,\ldots,x_{n-k+1}),
 \qquad \ExponentialCompleteBellPolynomial0=1.
\end{equation}
\end{definition}

\begin{definition}[Partial ordinary Bell polynomial]
\begin{equation}\label{eq:ordinary-bell-definition}
 \OrdinaryPartialBellPolynomial nk(x_1,\ldots,x_{n-k+1})
 =\sum_{\substack{j_1+j_2+\cdots=k\\
                       j_1+2j_2+\cdots=n}}
 \frac{k!}{j_1!j_2!\cdots}\prod_{i\ge1}x_i^{j_i}.
\end{equation}
Equivalently,
\begin{equation}\label{eq:ordinary-exponential-scaling}
 \OrdinaryPartialBellPolynomial nk(x_1,x_2,\ldots)
 =\frac{k!}{n!}\ExponentialPartialBellPolynomial nk(1!x_1,2!x_2,\ldots).
\end{equation}
\end{definition}
\begin{proof}[Proof of \eqref{eq:ordinary-exponential-scaling}]
Substitute $i!x_i$ for $x_i$ in \eqref{eq:partial-bell-definition}; the factorials
$(i!)^{j_i}$ cancel, and multiplication by $k!/n!$ leaves exactly
\eqref{eq:ordinary-bell-definition}.
\end{proof}

\section{Combinatorial interpretation}
Give each block of size $i$ a weight $x_i$, and define the weight of a partition as
the product of its block weights.

\begin{theorem}[Weighted-partition interpretation]\label{thm:bell-poly-partitions}
$\ExponentialPartialBellPolynomial nk$ is the total weight of partitions of an $n$-element labelled set into
exactly $k$ blocks, and $\ExponentialCompleteBellPolynomial n$ is the total weight of all its partitions.
\end{theorem}
\begin{proof}
Fix multiplicities $j_i$, where $j_i$ blocks have size $i$.  The number of such
partitions is
\[
 \frac{n!}{\prod_i(i!)^{j_i}j_i!}.
\]
The factors $i!$ forget order inside each block and $j_i!$ forget order among
blocks of equal size.  Multiplication by $\prod_i x_i^{j_i}$ and summation gives
\eqref{eq:partial-bell-definition}; summing over $k$ gives the complete polynomial.
\end{proof}

Thus
\begin{align*}
 \ExponentialPartialBellPolynomial{3}{2}(x_1,x_2)&=3x_1x_2,\\
 \ExponentialPartialBellPolynomial{6}{2}(x_1,\ldots,x_5)&=6x_1x_5+15x_2x_4+10x_3^2,\\
 \ExponentialPartialBellPolynomial{6}{3}(x_1,\ldots,x_4)&=15x_1^2x_4+60x_1x_2x_3+15x_2^3.
\end{align*}
The monomials correspond to integer partitions of $n$ into $k$ parts, while their
coefficients count the labelled set partitions having the prescribed block sizes.

\section{The master generating functions}
Put
\[
 X(t)=\sum_{j\ge1}x_j\frac{t^j}{j!}.
\]
\begin{theorem}[Bell-polynomial generating functions]\label{thm:bell-poly-egf}
One has
\begin{align}
 \frac{X(t)^k}{k!}
 &=\sum_{n\ge k}\ExponentialPartialBellPolynomial nk(x_1,\ldots)\frac{t^n}{n!},\\
 \exp(uX(t))
 &=\sum_{n\ge k\ge0}\ExponentialPartialBellPolynomial nk(x_1,\ldots)u^k\frac{t^n}{n!},\\
 \exp(X(t))
 &=\sum_{n\ge0}\ExponentialCompleteBellPolynomial n(x_1,\ldots,x_n)\frac{t^n}{n!}.
\end{align}
For $\OrdinaryGeneratingFunctionOf{X}(t)=\sum_{j\ge1}x_jt^j$,
\begin{align}
 \OrdinaryGeneratingFunctionOf{X}(t)^k&=\sum_{n\ge k}\OrdinaryPartialBellPolynomial nk(x_1,\ldots)t^n,
 \label{eq:ordinary-bell-ogf}\\
 \EulerE^{u\OrdinaryGeneratingFunctionOf{X}(t)}
 &=\sum_{n\ge k\ge0}\OrdinaryPartialBellPolynomial nk(x_1,\ldots)t^n\frac{u^k}{k!}.
 \label{eq:ordinary-bell-bivariate}
\end{align}
\end{theorem}
\begin{proof}
Expand $X(t)^k$ multinomially.  Choosing $j_i$ copies of the term
$x_it^i/i!$ gives the constraints $\sum j_i=k$ and $\sum ij_i=n$, and its
coefficient is $k!\prod_i x_i^{j_i}/((i!)^{j_i}j_i!)$.  Division by $k!$ gives
\eqref{eq:partial-bell-definition}.  Summing in $k$ proves the exponential
identities.  The ordinary identities follow by the same multinomial argument or
by \eqref{eq:ordinary-exponential-scaling}.
\end{proof}

\begin{masterbox}
Nearly every algebraic identity for Bell polynomials follows by performing a simple
operation---multiplication, differentiation, substitution, or logarithm---on
\eqref{eq:partial-bell-egf} or \eqref{eq:complete-bell-egf}.  The combinatorial
proofs are the same operations interpreted as selecting, pointing, colouring, or
gluing blocks.
\end{masterbox}

\section{Tables and low-degree polynomials}
The nonzero partial polynomials through $n=6$ are
\begin{align*}
\ExponentialPartialBellPolynomial11&=x_1,\\
\ExponentialPartialBellPolynomial21&=x_2,\qquad \ExponentialPartialBellPolynomial22=x_1^2,\\
\ExponentialPartialBellPolynomial31&=x_3,\qquad \ExponentialPartialBellPolynomial32=3x_1x_2,\qquad \ExponentialPartialBellPolynomial33=x_1^3,\\
\ExponentialPartialBellPolynomial41&=x_4,\qquad \ExponentialPartialBellPolynomial42=4x_1x_3+3x_2^2,\\
\ExponentialPartialBellPolynomial43&=6x_1^2x_2,\qquad \ExponentialPartialBellPolynomial44=x_1^4,\\
\ExponentialPartialBellPolynomial51&=x_5,\qquad \ExponentialPartialBellPolynomial52=5x_1x_4+10x_2x_3,\\
\ExponentialPartialBellPolynomial53&=10x_1^2x_3+15x_1x_2^2,\\
\ExponentialPartialBellPolynomial54&=10x_1^3x_2,\qquad \ExponentialPartialBellPolynomial55=x_1^5,\\
\ExponentialPartialBellPolynomial61&=x_6,\qquad \ExponentialPartialBellPolynomial62=6x_1x_5+15x_2x_4+10x_3^2,\\
\ExponentialPartialBellPolynomial63&=15x_1^2x_4+60x_1x_2x_3+15x_2^3,\\
\ExponentialPartialBellPolynomial64&=20x_1^3x_3+45x_1^2x_2^2,\\
\ExponentialPartialBellPolynomial65&=15x_1^4x_2,\qquad \ExponentialPartialBellPolynomial66=x_1^6.
\end{align*}
Every coefficient is supplied by \eqref{eq:partial-bell-definition}, so the table
is a computation from the theorem rather than an independent assertion.

The first complete polynomials are
\begin{align*}
 \ExponentialCompleteBellPolynomial0={}&1,\\
 \ExponentialCompleteBellPolynomial1={}&x_1,\\
 \ExponentialCompleteBellPolynomial2={}&x_1^2+x_2,\\
 \ExponentialCompleteBellPolynomial3={}&x_1^3+3x_1x_2+x_3,\\
 \ExponentialCompleteBellPolynomial4={}&x_1^4+6x_1^2x_2+3x_2^2+4x_1x_3+x_4,\\
 \ExponentialCompleteBellPolynomial5={}&x_1^5+10x_1^3x_2+15x_1x_2^2+10x_1^2x_3
 +10x_2x_3+5x_1x_4+x_5,\\
 \ExponentialCompleteBellPolynomial6={}&x_1^6+15x_1^4x_2+45x_1^2x_2^2+15x_2^3
 +20x_1^3x_3+60x_1x_2x_3+10x_3^2\\
 &+15x_1^2x_4+15x_2x_4+6x_1x_5+x_6,\\
 \ExponentialCompleteBellPolynomial7={}&x_1^7+21x_1^5x_2+105x_1^3x_2^2+105x_1x_2^3
 +35x_1^4x_3+210x_1^2x_2x_3\\
 &+70x_1x_3^2+105x_2^2x_3+35x_1^3x_4+105x_1x_2x_4
 +35x_3x_4\\
 &+21x_1^2x_5+21x_2x_5+7x_1x_6+x_7.
\end{align*}

\chapter{Recurrences, derivatives, and algebraic identities}

\section{Pointing recurrences}
\begin{theorem}[Complete and partial recurrences]\label{thm:bell-poly-recurrences}
For $n\ge0$,
\begin{align}
 \ExponentialCompleteBellPolynomial{n+1}(x_1,\ldots,x_{n+1})
 &=\sum_{i=0}^{n}\binom ni x_{i+1}\ExponentialCompleteBellPolynomial{n-i}(x_1,\ldots,x_{n-i}),
 \label{eq:complete-bell-recurrence}\\
 \ExponentialPartialBellPolynomial{n+1}{k+1}
 &=\sum_{i=0}^{n-k}\binom ni x_{i+1}\ExponentialPartialBellPolynomial{n-i}{k},
 \label{eq:partial-bell-recurrence}
\end{align}
with $\ExponentialPartialBellPolynomial00=1$, $\ExponentialPartialBellPolynomial n0=0$ for $n>0$, and $\ExponentialPartialBellPolynomial0k=0$ for $k>0$.
\end{theorem}
\begin{proof}
In a weighted partition of $[n+1]$, inspect the block containing the distinguished
label $n+1$.  If that block contains $i$ further labels, choose them in $\binom ni$
ways, give the block weight $x_{i+1}$, and partition the remaining labels.  Keeping
or forgetting the total block count yields the two formulas.
\end{proof}

\begin{theorem}[Block-colour convolution]\label{thm:bell-partial-convolution}
For $k_1,k_2\ge0$,
\begin{equation}\label{eq:bell-partial-convolution}
 \ExponentialPartialBellPolynomial n{k_1+k_2}(x)
 =\frac{k_1!k_2!}{(k_1+k_2)!}
 \sum_{i=0}^{n}\binom ni\ExponentialPartialBellPolynomial i{k_1}(x)\ExponentialPartialBellPolynomial{n-i}{k_2}(x).
\end{equation}
\end{theorem}
\begin{proof}
Multiply \eqref{eq:partial-bell-egf} for $k_1$ and $k_2$.  Their product is
$X^{k_1+k_2}/(k_1!k_2!)$; comparison with the same formula for $k_1+k_2$ gives the
factor $(k_1!k_2!)/(k_1+k_2)!$ and the binomial EGF convolution.
\end{proof}

\section{Near-diagonal reduction}
\begin{theorem}\label{thm:bell-near-diagonal}
For $1\le a<n$,
\begin{multline}\label{eq:bell-near-diagonal}
 \ExponentialPartialBellPolynomial n{n-a}(x_1,\ldots,x_{a+1})\\
 =\sum_{j=a+1}^{2a}\frac{j!}{a!}\binom nj x_1^{n-j}
 \ExponentialPartialBellPolynomial a{j-a}\!\left(\frac{x_2}{2},\frac{x_3}{3},\ldots,
 \frac{x_{2a-j+2}}{2a-j+2}\right).
\end{multline}
\end{theorem}
\begin{proof}
A partition into $n-a$ blocks has total \emph{excess}
$\sum_B(|B|-1)=a$.  Let $j$ be the number of elements in nonsingleton blocks.
Then the number of those blocks is $j-a$, so $a+1\le j\le2a$.  Choose the $j$
elements, leave the other $n-j$ as singleton blocks, and count the nonsingleton
partition.  Replacing a block of size $r+1$ by an atom of size $r$ removes one
distinguished element per block; the factor $j!/a!$ and weights $x_{r+1}/(r+1)$
are exactly the labelled pointing factors.  Equivalently, make the substitution in
\eqref{eq:partial-bell-definition}; the coefficient of every multiplicity vector
on both sides is
$n!x_1^{n-j}\prod_{r\ge2}x_r^{j_r}/((r!)^{j_r}j_r!)$.
\end{proof}

The first cases, obtained by listing integer partitions of the excess, are
\begin{align}
 \ExponentialPartialBellPolynomial n1&=x_n, &
 \ExponentialPartialBellPolynomial n2&=\frac12\sum_{j=1}^{n-1}\binom njx_jx_{n-j},
 \label{eq:bell-edge-first}\\
 \ExponentialPartialBellPolynomial nn&=x_1^n, &
 \ExponentialPartialBellPolynomial n{n-1}&=\binom n2x_1^{n-2}x_2,\\
 \ExponentialPartialBellPolynomial n{n-2}&=\binom n3x_1^{n-3}x_3
 +3\binom n4x_1^{n-4}x_2^2,\\
 \ExponentialPartialBellPolynomial n{n-3}&=\binom n4x_1^{n-4}x_4
 +10\binom n5x_1^{n-5}x_2x_3
 +15\binom n6x_1^{n-6}x_2^3,\\
 \ExponentialPartialBellPolynomial n{n-4}&=\binom n5x_1^{n-5}x_5
 +5\binom n6x_1^{n-6}(3x_2x_4+2x_3^2)\notag\\
 &\quad+105\binom n7x_1^{n-7}x_2^2x_3
 +105\binom n8x_1^{n-8}x_2^4.
 \label{eq:bell-near-diagonal-cases}
\end{align}
Their coefficients are the corresponding block-size counts from
\cref{thm:bell-poly-partitions}.

\section{Partial derivatives}
\begin{theorem}[Derivative identities]\label{thm:bell-poly-derivatives}
For $i\ge1$,
\begin{align}
 \frac{\partial\ExponentialCompleteBellPolynomial n}{\partial x_i}
 &=\binom ni\ExponentialCompleteBellPolynomial{n-i},
 \label{eq:complete-bell-derivative}\\
 \frac{\partial\ExponentialPartialBellPolynomial nk}{\partial x_i}
 &=\binom ni\ExponentialPartialBellPolynomial{n-i}{k-1}.
 \label{eq:partial-bell-derivative}
\end{align}
Consequently, for differentiable $a_i(x)$,
\begin{equation}\label{eq:bell-poly-chain}
 \frac{\Differential}{\Differential x}\ExponentialPartialBellPolynomial nk(a_1(x),\ldots)
 =\sum_{i=1}^{n-k+1}\binom ni a_i'(x)\ExponentialPartialBellPolynomial{n-i}{k-1}(a_1(x),\ldots).
\end{equation}
\end{theorem}
\begin{proof}
Differentiate \eqref{eq:partial-bell-egf} with respect to $x_i$:
\[
 \frac{\partial}{\partial x_i}\frac{X(t)^k}{k!}
 =\frac{t^i}{i!}\frac{X(t)^{k-1}}{(k-1)!}.
\]
Coefficient comparison gives \eqref{eq:partial-bell-derivative}; summing over $k$
gives \eqref{eq:complete-bell-derivative}.  The ordinary chain rule then gives
\eqref{eq:bell-poly-chain}.
\end{proof}

\section{A quadratic differential recurrence}
The less familiar differential recurrence in the archive is also a disguised
coefficient identity.
\begin{theorem}\label{thm:bell-quadratic-differential}
For $n\ge2$,
\begin{align}
 \ExponentialCompleteBellPolynomial n=\frac1{n-1}\Bigg[&
 \sum_{i=2}^{n}\sum_{j=1}^{i-1}(i-1)\binom{i-2}{j-1}x_jx_{i-j}
       \frac{\partial\ExponentialCompleteBellPolynomial{n-1}}{\partial x_{i-1}}\notag\\
 &+\sum_{i=2}^{n}\sum_{j=1}^{i-1}\frac{x_{i+1}}{\binom ij}
       \frac{\partial^2\ExponentialCompleteBellPolynomial{n-1}}{\partial x_j\partial x_{i-j}}
 +\sum_{i=2}^{n}x_i\frac{\partial\ExponentialCompleteBellPolynomial{n-1}}{\partial x_{i-1}}
 \Bigg].
 \label{eq:bell-quadratic-differential}
\end{align}
Terms containing unavailable variables or negative-index Bell polynomials are zero.
\end{theorem}
\begin{proof}
Let $F(t)=\EulerE^{X(t)}=\sum_m\ExponentialCompleteBellPolynomial m t^m/m!$.  After using
\eqref{eq:complete-bell-derivative}, the first double sum is
\[
 (n-1)!\CoefficientExtraction{t}{n-1}\bigl(tX'(t)^2F(t)\bigr).
\]
The second is
$(n-1)!\CoefficientExtraction{t}{n-1}((tX''-X'+x_1)F)$, and the last is
$(n-1)!\CoefficientExtraction{t}{n-1}((X'-x_1)F)$.  Their sum is therefore
\[
 (n-1)!\CoefficientExtraction{t}{n-1}t(X''+X'^2)F
 =(n-1)!\CoefficientExtraction{t}{n-1}tF''=(n-1)\ExponentialCompleteBellPolynomial n.
\]
Division by $n-1$ proves the formula.
\end{proof}

\chapter{Specializations, inversions, determinants, and convolutions}

\section{Stirling, Bell, Lah, and Touchard specializations}
\begin{theorem}\label{thm:bell-poly-specializations}
For admissible $n,k$,
\begin{align}
 \ExponentialPartialBellPolynomial nk(0!,1!,\ldots,(n-k)!)&=\UnsignedStirlingFirstKind nk,
 \label{eq:bell-first-specialization}\\
 \ExponentialCompleteBellPolynomial n(0!,1!,\ldots,(n-1)!)&=n!,
 \label{eq:bell-factorial-complete}\\
 \ExponentialPartialBellPolynomial nk(1,1,\ldots,1)&=\StirlingSecondKind nk,
 \label{eq:bell-second-specialization}\\
 \ExponentialCompleteBellPolynomial n(1,1,\ldots,1)&=\BellNumber n,
 \label{eq:bell-number-specialization}\\
 \ExponentialPartialBellPolynomial nk(1!,2!,\ldots,(n-k+1)!)
 &=\binom{n-1}{k-1}\frac{n!}{k!}=\LahNumber{n}{k},
 \label{eq:bell-lah-specialization}\\
 \TouchardPolynomial{n}(x)&=\ExponentialCompleteBellPolynomial n(x,x,\ldots,x)
 =\sum_{k=0}^{n}\StirlingSecondKind nkx^k.
 \label{eq:touchard-bell-specialization}
\end{align}
\end{theorem}
\begin{proof}
In \eqref{eq:partial-bell-egf}, the choices $x_j=(j-1)!$, $x_j=1$, and
$x_j=j!$ make $X(t)$ respectively
\[
 -\log(1-t),\qquad \EulerE^t-1,\qquad \frac{t}{1-t}.
\]
The resulting coefficient formulae are exactly the column EGFs for unsigned
first-kind Stirling numbers, second-kind Stirling numbers, and Lah numbers proved
earlier.  Summing over $k$ yields \eqref{eq:bell-factorial-complete} and
\eqref{eq:bell-number-specialization}.  Finally $x_j=x$ gives
$\exp(x(\EulerE^t-1))$, the EGF of $\TouchardPolynomial{n}(x)$.
\end{proof}

\section{Homogeneity and a useful evaluation}
\begin{theorem}[Bigraded homogeneity]\label{thm:bell-bihomogeneous}
For scalars $\alpha,\beta$,
\begin{equation}\label{eq:bell-bihomogeneous}
 \ExponentialPartialBellPolynomial nk(\alpha\beta x_1,\alpha\beta^2x_2,\ldots)
 =\alpha^k\beta^n\ExponentialPartialBellPolynomial nk(x_1,x_2,\ldots).
\end{equation}
Moreover,
\begin{equation}\label{eq:bell-linear-arguments}
 \ExponentialPartialBellPolynomial nk(1,2,3,\ldots,n-k+1)=\binom nk k^{n-k}.
\end{equation}
\end{theorem}
\begin{proof}
Every monomial in \eqref{eq:partial-bell-definition} has
$\sum j_i=k$ and $\sum ij_i=n$, proving the first formula.  For the second,
$x_j=j$ gives $X(t)=te^t$; hence
\[
 \frac{X(t)^k}{k!}=\frac{t^ke^{kt}}{k!}.
\]
The coefficient of $t^n/n!$ is
$n!k^{n-k}/(k!(n-k)!)=\binom nk k^{n-k}$.
\end{proof}

\section{Addition and shifted-zero identities}
\begin{theorem}[Complete Bell addition formula]\label{thm:complete-bell-addition}
\begin{equation}\label{eq:complete-bell-addition}
 \ExponentialCompleteBellPolynomial n(x_1+y_1,\ldots,x_n+y_n)
 =\sum_{i=0}^{n}\binom ni\ExponentialCompleteBellPolynomial{n-i}(x)\ExponentialCompleteBellPolynomial i(y).
\end{equation}
\end{theorem}
\begin{proof}
The complete EGF for the left side is
$\exp(X(t)+Y(t))=\exp X(t)\exp Y(t)$.  Coefficient extraction gives the binomial
convolution.
\end{proof}

\begin{theorem}[Insertion of $q$ leading zeros]\label{thm:bell-leading-zeros}
For $q\ge0$,
\begin{multline}\label{eq:bell-leading-zeros}
 \ExponentialPartialBellPolynomial nk\!\left(
 \frac{x_{q+1}}{\binom{q+1}{q}},
 \frac{x_{q+2}}{\binom{q+2}{q}},\ldots\right)\\
 =\frac{n!(q!)^k}{(n+qk)!}
 \ExponentialPartialBellPolynomial{n+qk}{k}(\underbrace{0,\ldots,0}_{q\text{ entries}},
 x_{q+1},x_{q+2},\ldots).
\end{multline}
\end{theorem}
\begin{proof}
Since
\[
 \frac{x_{q+i}}{\binom{q+i}{q}}\frac{t^i}{i!}
 =q!\,\frac{x_{q+i}t^i}{(q+i)!},
\]
the $k$th power of the left component series is $(q!)^k t^{-qk}$ times the
$k$th power of the component series on the right.  Comparing the coefficient of
$t^n$ in \eqref{eq:partial-bell-egf} gives the factorial factor displayed.
\end{proof}

\section{Exponential--logarithmic inversion}
\begin{theorem}[Bell transform and its inverse]\label{thm:bell-transform-inverse}
Suppose
\begin{equation}\label{eq:bell-transform-y}
 y_n=\ExponentialCompleteBellPolynomial n(x_1,\ldots,x_n)=\sum_{k=1}^{n}\ExponentialPartialBellPolynomial nk(x_1,\ldots).
\end{equation}
Then
\begin{equation}\label{eq:bell-transform-x}
 x_n=\sum_{k=1}^{n}(-1)^{k-1}(k-1)!
 \ExponentialPartialBellPolynomial nk(y_1,\ldots,y_{n-k+1}).
\end{equation}
\end{theorem}
\begin{proof}
Let $X(t)=\sum_{n\ge1}x_nt^n/n!$ and
$Y(t)=\sum_{n\ge1}y_nt^n/n!$.  Equation
\eqref{eq:complete-bell-egf} says $1+Y=\EulerE^X$, so
\[
 X=\log(1+Y)=\sum_{k\ge1}\frac{(-1)^{k-1}}{k}Y^k.
\]
Use \eqref{eq:partial-bell-egf} to extract the coefficient of $t^n/n!$.
\end{proof}

More generally, if $f$ and $g=f^{-1}$ are locally inverse and
\begin{equation}\label{eq:general-bell-transform}
 y_n=\sum_{k=0}^{n}f^{(k)}(a)\ExponentialPartialBellPolynomial nk(x_1,\ldots),
\end{equation}
then
\begin{equation}\label{eq:general-bell-inverse}
 x_n=\sum_{k=0}^{n}g^{(k)}(f(a))\ExponentialPartialBellPolynomial nk(y_1,\ldots).
\end{equation}
Indeed, these are the derivative coefficients of $f\circ u$ and
$g\circ f\circ u=u$; a complete proof is given with Fa\`a di Bruno's formula in
\cref{eq:faa-bell}.

\section{Two Hessenberg determinants}
For $n\ge1$, define $\BellHessenbergMatrixH{n}=(h_{ij})$ and $\BellHessenbergMatrixK{n}=(k_{ij})$ by
\[
 h_{ij}=\begin{cases}
 \binom{n-i}{j-i}x_{j-i+1},&j\ge i,\\
 -1,&j=i-1,\\0,&j<i-1,
 \end{cases}
 \qquad
 k_{ij}=\begin{cases}
 x_{j-i+1}/(j-i)!,&j\ge i,\\
 -(i-1),&j=i-1,\\0,&j<i-1.
 \end{cases}
\]
Thus $\BellHessenbergMatrixH{n}$ begins with
\[
\begin{bmatrix}
 x_1&\binom{n-1}{1}x_2&\binom{n-1}{2}x_3&\cdots&x_n\\
 -1&x_1&\binom{n-2}{1}x_2&\cdots&x_{n-1}\\
 0&-1&x_1&\cdots&x_{n-2}\\
 \vdots&&\ddots&\ddots&\vdots\\
 0&\cdots&0&-1&x_1
\end{bmatrix},
\]
while $\BellHessenbergMatrixK{n}$ has first row
$x_1,x_2/1!,\ldots,x_n/(n-1)!$ and subdiagonal
$-1,-2,\ldots,-(n-1)$.

\begin{theorem}[Bell determinants]\label{thm:bell-determinants}
\begin{equation}\label{eq:bell-determinants}
 \ExponentialCompleteBellPolynomial n(x_1,\ldots,x_n)=\det \BellHessenbergMatrixH{n}=\det \BellHessenbergMatrixK{n}.
\end{equation}
\end{theorem}
\begin{proof}
Expansion of the upper Hessenberg determinant $\det \BellHessenbergMatrixH{n}$ along its first row gives
\[
 D_n=\sum_{j=1}^{n}\binom{n-1}{j-1}x_jD_{n-j},\qquad D_0=1;
\]
the cofactor sign is cancelled by the product of the $j-1$ entries $-1$ on the
subdiagonal.  This is exactly \eqref{eq:complete-bell-recurrence}, so
$D_n=\ExponentialCompleteBellPolynomial n$.

Let $J$ reverse the coordinate order and let
$D=\DiagonalOperator(0!,1!,\ldots,(n-1)!)$.  A direct entrywise calculation gives
\[
 \BellHessenbergMatrixK{n}=D(J\BellHessenbergMatrixH{n}^{\mathsf T}J)D^{-1}.
\]
Indeed,
$\frac{(i-1)!}{(j-1)!}\binom{j-1}{i-1}=1/(j-i)!$ above the diagonal, and the
subdiagonal becomes $-(i-1)$.  Similarity, transposition, and reversal preserve
the determinant, proving the second equality.
\end{proof}

\section{The diamond convolution}
For sequences $x=(x_1,x_2,\ldots)$ and $y=(y_1,y_2,\ldots)$ define
\begin{equation}\label{eq:diamond}
 (x\mathbin\diamondsuit y)_n=\sum_{j=1}^{n-1}\binom njx_jy_{n-j}.
\end{equation}
This operation is commutative and associative because it corresponds to
multiplication of the EGFs $X(t)$ and $Y(t)$.

\begin{theorem}\label{thm:diamond-bell}
If $x^{k\diamondsuit}$ denotes the $k$-fold diamond power, then
\begin{equation}\label{eq:diamond-bell}
 \ExponentialPartialBellPolynomial nk(x_1,\ldots,x_{n-k+1})=\frac{(x^{k\diamondsuit})_n}{k!}.
\end{equation}
\end{theorem}
\begin{proof}
The EGF of $x^{k\diamondsuit}$ is $X(t)^k$.  Compare with
\eqref{eq:partial-bell-egf}.
\end{proof}
For example,
\[
 x\diamondsuit x=(0,2x_1^2,6x_1x_2,8x_1x_3+6x_2^2,\ldots)
\]
and
$(x^{3\diamondsuit})_4=36x_1^2x_2$, whence
$\ExponentialPartialBellPolynomial43=6x_1^2x_2$.

\chapter{Applications of Bell polynomials}

\section{Symmetric functions and Newton identities}
Let $e_n$ be the elementary symmetric function and
$p_n=\sum_i u_i^n$ the power sum in an alphabet $u_1,u_2,\ldots$.
\begin{theorem}\label{thm:bell-symmetric-functions}
\begin{align}
 e_n&=\frac1{n!}\ExponentialCompleteBellPolynomial n(p_1,-1!p_2,2!p_3,\ldots,(-1)^{n-1}(n-1)!p_n)
 \label{eq:elementary-via-bell}\\
 &=\frac{(-1)^n}{n!}\ExponentialCompleteBellPolynomial n(-p_1,-1!p_2,-2!p_3,\ldots,-(n-1)!p_n),
 \notag\\
 p_n&=\frac{(-1)^{n-1}}{(n-1)!}
 \sum_{k=1}^{n}(-1)^{k-1}(k-1)!
 \ExponentialPartialBellPolynomial nk(1!e_1,2!e_2,\ldots)
 \label{eq:powersum-via-bell}\\
 &=(-1)^n n\sum_{k=1}^{n}\frac1k
 \OrdinaryPartialBellPolynomial nk(-e_1,-e_2,\ldots).
 \notag
\end{align}
\end{theorem}
\begin{proof}
The classical product identity
\[
 E(t)=\sum_{n\ge0}e_nt^n=\prod_i(1+u_it)
 =\exp\!\left(\sum_{r\ge1}(-1)^{r-1}p_r\frac{t^r}{r}\right)
\]
and \eqref{eq:complete-bell-egf} give \eqref{eq:elementary-via-bell}.
Taking logarithms,
$\sum_{r\ge1}(-1)^{r-1}p_rt^r/r=\log(1+\sum_{j\ge1}e_jt^j)$, and expanding
the logarithm through ordinary Bell polynomials yields the second pair of forms.
The exponential form follows from \eqref{eq:ordinary-exponential-scaling}.
\end{proof}

\begin{corollary}[Determinant from traces]\label{cor:det-traces-bell}
For an $n\times n$ matrix $A$, set
$\TraceBellArgument{k}=-(k-1)!\TraceOperator(A^k)$.  Then
\begin{equation}\label{eq:det-traces-bell}
 \det A=\frac{(-1)^n}{n!}\ExponentialCompleteBellPolynomial n\bigl(
   \TraceBellArgument{1},\TraceBellArgument{2},\ldots,\TraceBellArgument{n}\bigr).
\end{equation}
\end{corollary}
\begin{proof}
Apply \eqref{eq:elementary-via-bell} to the eigenvalues of $A$, using
$e_n=\det A$ and $p_k=\TraceOperator(A^k)$.  Both sides are polynomial in the matrix entries,
so the argument remains valid for nondiagonalizable matrices.
\end{proof}

\section{Cycle indices}
Let $a_j$ mark cycles of length $j$ in a permutation.
\begin{theorem}\label{thm:cycle-index-bell}
The cycle index of the symmetric group is
\begin{equation}\label{eq:cycle-index-bell}
 Z(S_n)=\frac1{n!}\ExponentialCompleteBellPolynomial n(0!a_1,1!a_2,\ldots,(n-1)!a_n).
\end{equation}
\end{theorem}
\begin{proof}
A permutation is a set of cycles.  A labelled cycle of size $j$ has $(j-1)!$
linear realizations, so its component EGF is $(j-1)!a_jt^j/j!=a_jt^j/j$.
The exponential formula gives
$\exp(\sum_{j\ge1}a_jt^j/j)$; comparison with
\eqref{eq:complete-bell-egf} proves the identity after division by $n!$.
\end{proof}

\section{Moments and cumulants}
Let $M(t)=\Expectation \EulerE^{tX}=\sum_{n\ge0}\mu'_nt^n/n!$ and
$K(t)=\log M(t)=\sum_{n\ge1}\kappa_nt^n/n!$.
\begin{theorem}[Moment--cumulant relations]\label{thm:moment-cumulant}
\begin{align}
 \mu'_n&=\ExponentialCompleteBellPolynomial n(\kappa_1,\ldots,\kappa_n)
 =\sum_{k=1}^{n}\ExponentialPartialBellPolynomial nk(\kappa_1,\ldots),
 \label{eq:moments-from-cumulants}\\
 \kappa_n&=\sum_{k=1}^{n}(-1)^{k-1}(k-1)!
 \ExponentialPartialBellPolynomial nk(\mu'_1,\ldots,\mu'_{n-k+1}).
 \label{eq:cumulants-from-moments}
\end{align}
\end{theorem}
\begin{proof}
The identity $M=\EulerE^K$ is \eqref{eq:complete-bell-egf}; its logarithmic inverse is
\cref{thm:bell-transform-inverse}.
\end{proof}

\section{Hermite polynomials}
The probabilists' Hermite polynomials are defined by
\[
 \EulerE^{xt-t^2/2}=\sum_{n\ge0}\operatorname{He}_n(x)\frac{t^n}{n!}.
\]
Comparison with \eqref{eq:complete-bell-egf} gives and proves
\begin{equation}\label{eq:hermite-bell}
 \operatorname{He}_n(x)=\ExponentialCompleteBellPolynomial n(x,-1,0,\ldots,0).
\end{equation}
In particular, the usual recurrence and differential equation follow by
differentiating this generating function.

\section{Polynomial sequences of binomial type}
Let
$h(t)=\sum_{j\ge1}a_jt^j/j!$ with $a_1\ne0$, and define
\begin{equation}\label{eq:binomial-type-bell}
 p_n(x)=\ExponentialCompleteBellPolynomial n(a_1x,a_2x,\ldots,a_nx)
 =\sum_{k=0}^{n}\ExponentialPartialBellPolynomial nk(a_1,\ldots,a_{n-k+1})x^k.
\end{equation}

\begin{theorem}\label{thm:binomial-type-bell}
The polynomials $p_n$ satisfy
\begin{align}
 \sum_{n\ge0}p_n(x)\frac{t^n}{n!}&=\EulerE^{xh(t)},
 \label{eq:binomial-type-egf}\\
 p_n(x+y)&=\sum_{k=0}^{n}\binom nkp_k(x)p_{n-k}(y).
 \label{eq:binomial-type-identity}
\end{align}
If $Q=h^{-1}(D_x)$ is the associated delta operator, then
\begin{equation}\label{eq:delta-operator}
 Qp_n(x)=np_{n-1}(x).
\end{equation}
\end{theorem}
\begin{proof}
Substitute $x a_j$ in \eqref{eq:complete-bell-egf}.  The product
$\EulerE^{(x+y)h}=\EulerE^{xh}\EulerE^{yh}$ gives the binomial identity.  Finally
$h^{-1}(D_x)\EulerE^{xh(t)}=h^{-1}(h(t))\EulerE^{xh(t)}=t \EulerE^{xh(t)}$; coefficient comparison
gives \eqref{eq:delta-operator}.
\end{proof}

\begin{auditbox}
In some sources the symbol $B_n(x)$ denotes the Touchard polynomial, in others a
complete Bell polynomial, a Bernoulli polynomial, or a binary-tree sequence.  The
notations $\TouchardPolynomial{n}$, $\ExponentialCompleteBellPolynomial n$, $\beta_n(x)$, and $\CatalanNumber{n}$ are kept separate here; this
is essential when the objects occur in the same inversion formula.
\end{auditbox}

\part{Eulerian Numbers and Descent Calculus}

\chapter{Ordinary Eulerian numbers}
\index{Eulerian numbers}

\section{Descents, symmetry, and the basic recurrence}
For a permutation $\pi=\pi_1\cdots\pi_n$, a descent is an index
$i\in\{1,\ldots,n-1\}$ with $\pi_i>\pi_{i+1}$.
\begin{definition}
The Eulerian number $\TypeAEulerianNumber nk=A(n,k)$ counts permutations of $[n]$ with
exactly $k$ descents.  Put
\[
 \TypeAEulerianPolynomial{0}(t)=1,
 \qquad
 \TypeAEulerianPolynomial{n}(t)=\sum_{k=0}^{n-1}\TypeAEulerianNumber nk t^k\quad(n\ge1).
\]
Outside $0\le k<n$ the number is zero.
\end{definition}

\begin{theorem}[Eulerian recurrence]\label{thm:eulerian-recurrence}
\begin{equation}\label{eq:eulerian-recurrence}
 \TypeAEulerianNumber nk=(k+1)\TypeAEulerianNumber{n-1}{k}
 +(n-k)\TypeAEulerianNumber{n-1}{k-1}.
\end{equation}
Equivalently,
\begin{equation}\label{eq:eulerian-poly-recurrence}
 \TypeAEulerianPolynomial{n}(t)=\bigl(1+(n-1)t\bigr)\TypeAEulerianPolynomial{n-1}(t)
       +t(1-t)\TypeAEulerianPolynomial{n-1}'(t).
\end{equation}
\end{theorem}
\begin{proof}
Insert $n$ into a permutation of $[n-1]$.  If the old permutation has $k$
descents, insertion after one of those $k$ descents or at the end preserves the
number of descents, giving $k+1$ choices.  If it has $k-1$ descents, insertion in
any of the remaining $n-k$ slots creates one new descent.  This proves
\eqref{eq:eulerian-recurrence}.  Multiply by $t^k$, sum over $k$, and collect the
terms containing $k$ to obtain \eqref{eq:eulerian-poly-recurrence}.
\end{proof}

Complementation $\pi_i\mapsto n+1-\pi_i$ interchanges descents and ascents.
Therefore
\begin{equation}\label{eq:eulerian-symmetry}
 \TypeAEulerianNumber nk=\TypeAEulerianNumber n{n-1-k},
 \qquad
 \TypeAEulerianPolynomial{n}(t)=t^{n-1}\TypeAEulerianPolynomial{n}(t^{-1}).
\end{equation}
In particular $\TypeAEulerianNumber n0=\TypeAEulerianNumber n{n-1}=1$ and
$\sum_k\TypeAEulerianNumber nk=n!$.  The next diagonal is
\begin{equation}\label{eq:eulerian-k1}
 \TypeAEulerianNumber n1=2^n-(n+1),
\end{equation}
which follows either from the recurrence or from the explicit formula below.

\section{Exponential generating function and explicit formula}
\begin{theorem}[Eulerian EGF]\label{thm:eulerian-egf}
\begin{equation}\label{eq:eulerian-egf}
 \sum_{n\ge0}\TypeAEulerianPolynomial{n}(t)\frac{x^n}{n!}
 =\frac{t-1}{t-\EulerE^{(t-1)x}}
 =\frac{(1-t)\EulerE^{(1-t)x}}{1-te^{(1-t)x}}.
\end{equation}
\end{theorem}
\begin{proof}
Let $F(x,t)$ denote the left side.  Multiplying
\eqref{eq:eulerian-poly-recurrence} by $x^{n-1}/(n-1)!$ and summing gives the
first-order equation
\[
 (1-tx)F_x=t(1-t)F_t+F,
 \qquad F(0,t)=1.
\]
Direct substitution verifies that the displayed rational-exponential function
solves this equation and initial condition.  Formal uniqueness follows by
recursively comparing coefficients of $x$.
\end{proof}

Expanding the denominator geometrically gives the central power-series identity.
\begin{theorem}\label{thm:eulerian-power-series}
For $n\ge0$,
\begin{equation}\label{eq:eulerian-power-series}
 \sum_{m=0}^{\infty}(m+1)^n t^m
 =\frac{\TypeAEulerianPolynomial{n}(t)}{(1-t)^{n+1}}.
\end{equation}
Consequently,
\begin{equation}\label{eq:eulerian-explicit}
 \TypeAEulerianNumber nk=\sum_{j=0}^{k}(-1)^j\binom{n+1}{j}(k+1-j)^n.
\end{equation}
\end{theorem}
\begin{proof}
From \eqref{eq:eulerian-egf},
\[
 F(x,t)=(1-t)\sum_{m\ge0}t^m \EulerE^{(m+1)(1-t)x}.
\]
Extracting $x^n/n!$ proves \eqref{eq:eulerian-power-series}.  Multiply by
$(1-t)^{n+1}$ and extract $t^k$ to obtain
\eqref{eq:eulerian-explicit}.  Formula \eqref{eq:eulerian-k1} is its case $k=1$.
\end{proof}

\begin{theorem}[A second polynomial recurrence]\label{thm:eulerian-binomial-recurrence}
For $n>1$,
\begin{equation}\label{eq:eulerian-binomial-recurrence}
 \TypeAEulerianPolynomial{n}(t)=\sum_{k=0}^{n-1}\binom nk \TypeAEulerianPolynomial{k}(t)(t-1)^{n-1-k}.
\end{equation}
\end{theorem}
\begin{proof}
The EGF of the right sides, summed over $n\ge1$, is
\[
 F(x,t)\frac{\EulerE^{(t-1)x}-1}{t-1}.
\]
Using \eqref{eq:eulerian-egf}, this equals $F(x,t)-1$, which is the EGF of the
left sides.
\end{proof}

\section{Worpitzky's identity and sums of powers}
\begin{theorem}[Worpitzky]\label{thm:worpitzky}
For every indeterminate $x$ and $n\ge1$,
\begin{equation}\label{eq:worpitzky}
 x^n=\sum_{k=0}^{n-1}\TypeAEulerianNumber nk\binom{x+k}{n}.
\end{equation}
\end{theorem}
\begin{proof}
The coefficient of $t^{x-1}$ in
$t^{-k}(1-t)^{-n-1}$ is $\binom{x+k}{n}$ for positive integer $x$.
Equivalently, taking the coefficient of $t^{x-1}$ in
\eqref{eq:eulerian-power-series} after multiplying by the numerator gives
\eqref{eq:worpitzky}.  Since both sides are polynomials of degree $n$ in $x$ and
agree for all positive integers, they agree identically.

A direct combinatorial proof standardizes a word in $[x]^n$: its stable sorting
permutation records the relative order, and its $k$ descents specify $k$ mandatory
separations among the $n$ sorted positions.  Stars and bars gives
$\binom{x+k}{n}$ words for each permutation with $k$ descents.
\end{proof}

\begin{corollary}[Power sums]\label{cor:eulerian-power-sum}
For $m,n\ge1$,
\begin{equation}\label{eq:eulerian-power-sum}
 \sum_{r=1}^{m}r^n
 =\sum_{k=0}^{n-1}\TypeAEulerianNumber nk\binom{m+k+1}{n+1}.
\end{equation}
\end{corollary}
\begin{proof}
Sum \eqref{eq:worpitzky} over $x=1,\ldots,m$ and use the hockey-stick identity
$\sum_{x=1}^{m}\binom{x+k}{n}=\binom{m+k+1}{n+1}$; the omitted lower boundary
terms vanish for $0\le k<n$.
\end{proof}

\section{Negative polylogarithms}
For $|z|<1$, $\Polylogarithm_{-n}(z)=\sum_{m\ge1}m^nz^m$.  Shifting $m$ in
\eqref{eq:eulerian-power-series} gives
\begin{equation}\label{eq:negative-polylog-eulerian}
 \Polylogarithm_{-n}(z)=\frac{z\TypeAEulerianPolynomial{n}(z)}{(1-z)^{n+1}}
 =\frac{1}{(1-z)^{n+1}}
   \sum_{k=0}^{n-1}\TypeAEulerianNumber nk z^{n-k},
\end{equation}
where the second equality uses symmetry \eqref{eq:eulerian-symmetry}.  Since the
right side is rational, it also supplies the meromorphic continuation of the
negative polylogarithm.

\section{Eulerian and Stirling bases}
\begin{theorem}\label{thm:eulerian-stirling}
\begin{equation}\label{eq:eulerian-stirling}
 \TypeAEulerianPolynomial{n}(t)=\sum_{k=0}^{n}\StirlingSecondKind nk k!(t-1)^{n-k}.
\end{equation}
\end{theorem}
\begin{proof}
Use $m^n=\sum_k\StirlingSecondKind nk\FallingFactorial mk$ in
$\sum_{m\ge0}m^nt^m$.  Since
$\sum_{m\ge0}\FallingFactorial mk t^m=k!t^k/(1-t)^{k+1}$, multiplication by
$(1-t)^{n+1}/t$ yields
$\TypeAEulerianPolynomial{n}(t)=\sum_k\StirlingSecondKind nk k!t^{k-1}(1-t)^{n-k}$.  Replacing $t$ by $1/t$ and
using \eqref{eq:eulerian-symmetry}, or simply expanding around $t=1$, gives the
equivalent displayed form.
\end{proof}

\section{An Irwin--Hall integral and hypersimplex volumes}
\begin{theorem}\label{thm:eulerian-irwin-hall}
For every function $f:\RealNumbers\to\ComplexNumbers$ for which the following finite sum is defined,
\begin{equation}\label{eq:eulerian-irwin-hall}
 \sum_{k=0}^{n-1}\TypeAEulerianNumber nk f(k)
 =n!\int_{[0,1]^n}f\!\left(\Floor{x_1+\cdots+x_n}\right)
 \Differential x_1\cdots\Differential x_n.
\end{equation}
In particular, the slab
\[
 \Delta_{n,k}=\{x\in[0,1]^n:k\le x_1+\cdots+x_n<k+1\}
\]
has volume $\TypeAEulerianNumber nk/n!$.
\end{theorem}
\begin{proof}
The volume of $\{x\in[0,1]^n:\sum x_i<y\}$ is obtained from the simplex
$\{x_i\ge0:\sum x_i<y\}$ by inclusion--exclusion over coordinates exceeding $1$:
\[
 V_n(y)=\frac1{n!}\sum_{j=0}^{\Floor{y}}(-1)^j\binom nj(y-j)^n.
\]
Thus the volume of the $k$th slab is $V_n(k+1)-V_n(k)$.  Combining adjacent terms
by Pascal's identity gives
\[
 \frac1{n!}\sum_{j=0}^{k}(-1)^j\binom{n+1}{j}(k+1-j)^n
 =\frac1{n!}\TypeAEulerianNumber nk
\]
by \eqref{eq:eulerian-explicit}.  Summing $f(k)$ times these volumes proves the
integral identity.
\end{proof}

\section{Alternating sums and Bernoulli numbers}
\begin{theorem}\label{thm:eulerian-alternating}
For $n>0$,
\begin{equation}\label{eq:eulerian-alternating}
 \sum_{k=0}^{n-1}(-1)^k\TypeAEulerianNumber nk
 =2^{n+1}(2^{n+1}-1)\frac{\beta_{n+1}}{n+1}.
\end{equation}
For $n>1$,
\begin{align}
 \sum_{k=0}^{n-1}\frac{(-1)^k\TypeAEulerianNumber nk}{\binom{n-1}{k}}&=0,
 \label{eq:eulerian-reciprocal-zero}\\
 \sum_{k=0}^{n-1}\frac{(-1)^k\TypeAEulerianNumber nk}{\binom nk}&=(n+1)\beta_n.
 \label{eq:eulerian-reciprocal-bernoulli}
\end{align}
\end{theorem}
\begin{proof}
At $t=-1$, \eqref{eq:eulerian-egf} becomes
$2/(1+\EulerE^{-2x})=1+\tanh x$.  The Bernoulli generating function gives
\[
 \tanh x=\sum_{n\ge1}
 2^{n+1}(2^{n+1}-1)\frac{\beta_{n+1}}{n+1}\frac{x^n}{n!},
\]
which proves \eqref{eq:eulerian-alternating}.

For the remaining identities use
\[
 \binom{n-1}{k}^{-1}=n\int_0^1u^k(1-u)^{n-k-1}\Differential u,
 \quad
 \binom nk^{-1}=(n+1)\int_0^1u^k(1-u)^{n-k}\Differential u.
\]
Substitute $t=-u/(1-u)$ in \eqref{eq:eulerian-stirling}; then
\begin{equation}\label{eq:eulerian-beta-transform}
 (1-u)^n\TypeAEulerianPolynomial{n}\!\left(-\frac{u}{1-u}\right)
 =\sum_{j=0}^{n}(-1)^{n-j}\StirlingSecondKind njj!(1-u)^j.
\end{equation}
For the first integral divide by $1-u$ and integrate.  The result is
$n\sum_{j=1}^n(-1)^{n-j}\StirlingSecondKind nj(j-1)!=0$ for $n>1$, because it is the
coefficient of $x^n/n!$ in
$\log(1+(\EulerE^x-1))=x$.  For the second integral one obtains
$(n+1)\sum_j(-1)^{n-j}\StirlingSecondKind njj!/(j+1)$.  The standard Stirling formula
$\beta_n=\sum_j(-1)^jj!\StirlingSecondKind nj/(j+1)$ proves the claim; for odd $n>1$ both
sides vanish, and for even $n$ the signs coincide.
\end{proof}

\chapter{Geometric interpretations}

\section{The cube and its Ehrhart numerator}
For a lattice polytope $P$ of dimension $d$, its Ehrhart series has the form
$\sum_{m\ge0}|mP\cap\IntegerNumbers^d|t^m=h^*(t)/(1-t)^{d+1}$.
For the unit cube $C=[0,1]^n$ one has $|mC\cap\IntegerNumbers^n|=(m+1)^n$; hence
\eqref{eq:eulerian-power-series} proves:
\begin{theorem}
The $h^*$-polynomial of the standard $n$-cube is $\TypeAEulerianPolynomial{n}(t)$, so its $h^*$-vector is
the $n$th Eulerian row.
\end{theorem}
\begin{proof}
The dilate $m[0,1]^n=[0,m]^n$ contains exactly $(m+1)^n$ lattice points.
Therefore its Ehrhart series is
\[
 \sum_{m\ge0}(m+1)^nt^m.
\]
By \eqref{eq:eulerian-power-series} this equals
$\TypeAEulerianPolynomial{n}(t)/(1-t)^{n+1}$.  The numerator in the Ehrhart-series normal form is
unique, hence $h^*_{[0,1]^n}(t)=\TypeAEulerianPolynomial{n}(t)$; coefficient comparison gives the
stated $h^*$-vector.
\end{proof}

The regions $\Delta_{n,k}$ in \cref{thm:eulerian-irwin-hall} are half-open
hypersimplices.  Their normalized volumes are $n!\operatorname{vol}\Delta_{n,k}
=\TypeAEulerianNumber nk$.

\section{The permutohedron}
The type-$A$ permutohedron is the convex hull of the permutations of
$(1,2,\ldots,n)$; its dimension is $n-1$.  Faces of its braid fan are indexed by
ordered set partitions.  There are $k!\StirlingSecondKind nk$ ordered partitions into $k$
blocks.

\begin{theorem}
The $h$-polynomial of the simple polytope dual to the type-$A$ permutohedron is
$\TypeAEulerianPolynomial{n}(t)$.
\end{theorem}
\begin{proof}
For a simple $(n-1)$-polytope the $h$-polynomial is the transform
$h(t)=\sum_F(t-1)^{\dim F}$ over faces of the dual simplicial complex.  Grouping
braid-fan faces by their $k$ ordered blocks gives
$\sum_k k!\StirlingSecondKind nk(t-1)^{n-k}$, which is $\TypeAEulerianPolynomial{n}(t)$ by
\eqref{eq:eulerian-stirling}.
\end{proof}

\chapter{Type-B and second-order Eulerian numbers}

\section{Signed permutations and the corrected type-B indexing}
A signed permutation is a word $\pi_1\cdots\pi_n$ whose absolute values form a
permutation of $[n]$.  Put $\pi_0=0$ and call
$i\in\{0,\ldots,n-1\}$ a type-$B$ descent when $\pi_i>\pi_{i+1}$.
Let $\TypeBEulerianNumber{n}{k}$ count signed permutations with $k$ such descents, and set
$\TypeBEulerianPolynomial{n}(t)=\sum_{k=0}^{n}\TypeBEulerianNumber{n}{k}t^k$.

\begin{auditbox}
The archived display reverses the two entries in the descent inequality and pairs
that convention with an incompatible explicit formula.  The definition above is
the standard Coxeter descent convention.  It gives the archived table
$1,6,1$ for $n=2$ and the midpoint-Eulerian generating function below.
\end{auditbox}

\begin{theorem}[Type-$B$ recurrence and Worpitzky series]\label{thm:typeB-eulerian}
\begin{align}
 \TypeBEulerianNumber{n}{k}&=(2k+1)\TypeBEulerianNumber{n-1}{k}+(2n-2k+1)\TypeBEulerianNumber{n-1}{k-1},
 \label{eq:typeB-recurrence}\\
 \sum_{m\ge0}(2m+1)^nt^m&=\frac{\TypeBEulerianPolynomial{n}(t)}{(1-t)^{n+1}},
 \label{eq:typeB-power-series}\\
 \TypeBEulerianNumber{n}{k}&=\sum_{j=0}^{k}(-1)^j\binom{n+1}{j}
             \bigl(2(k-j)+1\bigr)^n.
 \label{eq:typeB-explicit}
\end{align}
The first rows are
\[
1;\quad 1,1;\quad1,6,1;\quad1,23,23,1;\quad1,76,230,76,1.
\]
\end{theorem}
\begin{proof}
Insert either $n$ or $-n$ into a signed permutation of size $n-1$.  Marking the
slots according to whether insertion preserves or creates a descent gives
$2k+1$ preserving slots for an object with $k$ descents and
$2n-2k+1$ creating slots for one with $k-1$ descents; this proves the recurrence.

A type-$B$ standardization partitions the $(2m+1)^n$ words in the ordered alphabet
$-m<\cdots<-1<0<1<\cdots<m$ by their signed sorting permutation.  If that
permutation has $k$ type-$B$ descents, stars and bars gives
$\binom{m+n-k}{n}$ compatible weakly increasing absolute levels.  Hence
\[
 (2m+1)^n=\sum_{k=0}^{n}\TypeBEulerianNumber{n}{k}\binom{m+n-k}{n}.
\]
Multiply by $t^m$, sum in $m$, and use
$\sum_{m\ge0}\binom{m+n-k}{n}t^m=t^k/(1-t)^{n+1}$ to obtain
\eqref{eq:typeB-power-series}.  Multiplication by $(1-t)^{n+1}$ and coefficient
extraction gives \eqref{eq:typeB-explicit}.
\end{proof}

As in type $A$, these coefficients form the $h$-vector of the simple polytope dual
to the type-$B$ permutohedron; the proof repeats the braid-fan argument with signed
ordered partitions.

\section{Stirling permutations and second-order Eulerian numbers}
A \emph{Stirling permutation of order $n$} is a permutation of the multiset
$\{1,1,2,2,\ldots,n,n\}$ such that all entries between the two copies of $i$ are
larger than $i$.  Let $\SecondOrderEulerianNumber nk$ count those with $k$ ordinary descents and put
\[
 \SecondOrderEulerianPolynomial{n}(t)=\sum_{k=0}^{n-1}\SecondOrderEulerianNumber nk t^k,
 \qquad \SecondOrderEulerianPolynomial{0}(t)=1.
\]

\begin{theorem}[Second-order recurrence]\label{thm:second-eulerian-recurrence}
\begin{align}
 \SecondOrderEulerianNumber nk&=(2n-k-1)\SecondOrderEulerianNumber{n-1}{k-1}
 +(k+1)\SecondOrderEulerianNumber{n-1}{k},
 \label{eq:second-eulerian-recurrence}\\
 \SecondOrderEulerianPolynomial{n+1}(t)&=(2nt+1)\SecondOrderEulerianPolynomial{n}(t)-t(t-1)\SecondOrderEulerianPolynomial{n}'(t),
 \label{eq:second-eulerian-poly-recurrence}\\
 \SecondOrderEulerianPolynomial{n}(1)&=(2n-1)!!.
 \label{eq:second-eulerian-row-sum}
\end{align}
\end{theorem}
\begin{proof}
The two copies of the largest label $n$ must be adjacent.  Insert the pair $nn$
into one of the $2n-1$ gaps of a Stirling permutation of order $n-1$.  A descent
gap or the terminal gap preserves the number of descents, giving $k+1$ choices
from an object with $k$ descents.  Each of the remaining $2n-k-1$ gaps creates one
descent from an object with $k-1$ descents.  This is the first recurrence; summing
against $t^k$ gives the second.  At each stage there are $2n-1$ insertion gaps, so
the total count is $(2n-1)!!$.
\end{proof}

The differential recurrence has two useful equivalent forms:
\begin{align}
 (t-1)^{-2n-2}\SecondOrderEulerianPolynomial{n+1}(t)
 &=\left(t(1-t)^{-2n-1}\SecondOrderEulerianPolynomial{n}(t)\right)',
 \label{eq:second-eulerian-differential-form}\\
 u_{n+1}(t)&=\left(\frac{t}{1-t}u_n(t)\right)',
 \quad u_n(t)=(t-1)^{-2n}\SecondOrderEulerianPolynomial{n}(t),\quad u_0=1.
 \label{eq:second-eulerian-u}
\end{align}
They follow by the product rule from
\eqref{eq:second-eulerian-poly-recurrence}.

\begin{theorem}[Diagonal Stirling generating function]\label{thm:second-eulerian-stirling-gf}
For $n\ge1$,
\begin{equation}\label{eq:second-eulerian-stirling-gf}
 \sum_{m=0}^{\infty}\StirlingSecondKind{n+m}{m}t^m
 =\frac{t\SecondOrderEulerianPolynomial{n}(t)}{(1-t)^{2n+1}}.
\end{equation}
\end{theorem}
\begin{proof}
Let $F_n(t)$ denote the left side.  The second-kind recurrence gives
$(1-t)F_{n+1}=tF_n'$, with $F_0=(1-t)^{-1}$.  Hence
$F_1=t/(1-t)^3$.  If the displayed formula holds for $n$, differentiating
$t\SecondOrderEulerianPolynomial{n}/(1-t)^{2n+1}$ and using
\eqref{eq:second-eulerian-poly-recurrence} gives the formula for $n+1$.
\end{proof}


\section{Second-order Eulerian interpolation of first-kind diagonals}
The same second-order Eulerian row that controls
\eqref{eq:second-eulerian-stirling-gf} also gives a binomial-basis expansion for
the near-diagonals of the unsigned first-kind triangle.

\begin{theorem}[Second-order Eulerian diagonal formula]
\label{thm:second-eulerian-first-diagonal}
For $p,m\ge0$,
\begin{equation}
\label{eq:second-eulerian-first-diagonal}
 \UnsignedStirlingFirstKind m{m-p}
 =\sum_{j=0}^{p}\SecondOrderEulerianNumber pj\binom{m+j}{2p},
\end{equation}
where an out-of-range first-kind number is zero.  Equivalently, the polynomial
continuation in a variable $x$ satisfies
\[
 \UnsignedStirlingFirstKind{x}{x-p}
 =\sum_{j=0}^{p}\SecondOrderEulerianNumber pj\binom{x+j}{2p}.
\]
The first instances are
\begin{align*}
 \UnsignedStirlingFirstKind{x}{x-1}&=\binom{x}{2},\\
 \UnsignedStirlingFirstKind{x}{x-2}&=\binom{x}{4}+2\binom{x+1}{4},\\
 \UnsignedStirlingFirstKind{x}{x-3}&=\binom{x}{6}+8\binom{x+1}{6}
                    +6\binom{x+2}{6}.
\end{align*}
\end{theorem}
\begin{proof}
For fixed $p$ define
\[
 D_p(t)=\sum_{m\ge0}\UnsignedStirlingFirstKind m{m-p}t^m.
\]
Writing $a_{m,p}=\UnsignedStirlingFirstKind m{m-p}$, the first-kind recurrence becomes
$a_{m,p}=a_{m-1,p}+(m-1)a_{m-1,p-1}$.  Therefore
\begin{equation}
\label{eq:first-diagonal-ogf-recurrence}
 (1-t)D_p(t)=t^2D_{p-1}'(t),
 \qquad D_0(t)=\frac1{1-t}.
\end{equation}
Let
\[
 Q_p(t)=t^{2p}\SecondOrderEulerianPolynomial{p}(t^{-1}).
\]
The polynomial recurrence \eqref{eq:second-eulerian-poly-recurrence} is equivalent,
after the substitution $t\mapsto t^{-1}$, to
\[
 Q_p(t)=t^2\bigl((1-t)Q_{p-1}'(t)+(2p-1)Q_{p-1}(t)\bigr).
\]
It follows inductively from \eqref{eq:first-diagonal-ogf-recurrence} that
\begin{equation}
\label{eq:first-diagonal-second-eulerian-ogf}
 D_p(t)=\frac{Q_p(t)}{(1-t)^{2p+1}}
 =\frac{t^{2p}\SecondOrderEulerianPolynomial{p}(t^{-1})}{(1-t)^{2p+1}}.
\end{equation}
Since $\SecondOrderEulerianPolynomial{p}(u)=\sum_{j=0}^{p}\SecondOrderEulerianNumber pj u^j$, coefficient extraction gives
\[
 [t^m]D_p(t)
 =\sum_{j=0}^{p}\SecondOrderEulerianNumber pj
   [t^{m-2p+j}](1-t)^{-2p-1}
 =\sum_{j=0}^{p}\SecondOrderEulerianNumber pj\binom{m+j}{2p},
\]
which is \eqref{eq:second-eulerian-first-diagonal}.  Both sides are polynomials in
$m$ of degree at most $2p$, so the same identity gives the stated polynomial
continuation.
\end{proof}
\part{Differentiation, Composition, and Reversion}

\chapter{Fa\`a di Bruno's formula}
\index{Faà di Bruno formula@Fa\`a di Bruno formula}

\section{The set-partition form}
For a block $B\subseteq[n]$, write $|B|$ for its size.
\begin{theorem}[Fa\`a di Bruno, partition form]\label{thm:faa-partition}
For $n\ge0$ and sufficiently differentiable $f,g$,
\begin{equation}\label{eq:faa-partition}
 D^n(f\circ g)(x)=
 \sum_{\pi\in\SetPartitionOperator([n])}
 f^{(|\pi|)}(g(x))\prod_{B\in\pi}g^{(|B|)}(x).
\end{equation}
For $n=0$, the empty partition contributes $f(g(x))$.
\end{theorem}
\begin{proof}
Induct on $n$.  Differentiate a term indexed by a partition $\pi$ of $[n]$.
If the derivative hits $f^{(|\pi|)}(g)$, the chain rule contributes $g'$ and
creates the new singleton block $\{n+1\}$.  If it hits the factor associated with
a block $B$, it changes $g^{(|B|)}$ to $g^{(|B|+1)}$, which inserts $n+1$ into
$B$.  Every partition of $[n+1]$ arises uniquely in exactly one of these ways:
remove $n+1$, either deleting its singleton block or removing it from its
non-singleton block.  Hence all terms occur once.
\end{proof}

\section{Multiplicity and Bell-polynomial forms}
Grouping partitions by the number $m_j$ of blocks of size $j$ gives the classical
formula.
\begin{theorem}[Fa\`a di Bruno, multiplicity form]\label{thm:faa-multiplicity}
\begin{multline}\label{eq:faa-multiplicity}
 D^n f(g(x))=
 \sum_{\substack{m_1,\ldots,m_n\ge0\\
                  m_1+2m_2+\cdots+nm_n=n}}
 \frac{n!}{m_1!m_2!\cdots m_n!}\\
 {}\times f^{(m_1+\cdots+m_n)}(g(x))
 \prod_{j=1}^{n}\left(\frac{g^{(j)}(x)}{j!}\right)^{m_j}.
\end{multline}
Equivalently,
\begin{equation}\label{eq:faa-bell}
 D^n f(g(x))=
 \sum_{k=0}^{n}f^{(k)}(g(x))
 \ExponentialPartialBellPolynomial nk\bigl(g'(x),g''(x),\ldots,g^{(n-k+1)}(x)\bigr).
\end{equation}
\end{theorem}
\begin{proof}
For fixed multiplicities $(m_j)$, the number of set partitions of $[n]$ with
$m_j$ blocks of size $j$ is
$n!/\prod_j((j!)^{m_j}m_j!)$.  Substitution in
\eqref{eq:faa-partition} gives \eqref{eq:faa-multiplicity}.  Grouping further by
$k=\sum m_j$ and invoking \eqref{eq:partial-bell-definition} gives
\eqref{eq:faa-bell}.
\end{proof}

For $n=4$ this reads
\begin{align}
 (f\circ g)^{(4)}={}&f^{(4)}(g)(g')^4
 +6f^{(3)}(g)(g')^2g''+3f''(g)(g'')^2\notag\\
 &+4f''(g)g'g^{(3)}+f'(g)g^{(4)}.
 \label{eq:faa-n4}
\end{align}
The coefficients $1,6,3,4,1$ count partitions of types
$1+1+1+1$, $2+1+1$, $2+2$, $3+1$, and $4$ respectively.  This also proves the
expanded example in the archive without four separate differentiations.

\section{The exponential special case}
Because every derivative of $\EulerE^u$ equals $\EulerE^u$, summing the partial Bell
polynomials in \eqref{eq:faa-bell} yields
\begin{equation}\label{eq:exp-faa}
 D^n \EulerE^{g(x)}=\EulerE^{g(x)}
 \ExponentialCompleteBellPolynomial n\bigl(g'(x),g''(x),\ldots,g^{(n)}(x)\bigr).
\end{equation}
When $g$ is a cumulant generating function, this is precisely the
moment--cumulant relation of \cref{thm:moment-cumulant}.

\section{Mixed partial derivatives}
For $B\subseteq[n]$, write
$\partial_B=\prod_{j\in B}\partial/\partial x_j$.
\begin{theorem}[Multivariate Fa\`a di Bruno for distinct variables]
\label{thm:faa-multivariate}
If $y=g(x_1,\ldots,x_n)$, then
\begin{equation}\label{eq:faa-multivariate}
 \frac{\partial^n}{\partial x_1\cdots\partial x_n}f(y)
 =\sum_{\pi\in\SetPartitionOperator([n])}f^{(|\pi|)}(y)
   \prod_{B\in\pi}\partial_B y.
\end{equation}
\end{theorem}
\begin{proof}
The proof of \cref{thm:faa-partition} applies verbatim, except that the new
derivative is $\partial/\partial x_{n+1}$.  It either creates a singleton block or
joins the new index to one existing block.
\end{proof}

For three variables, \eqref{eq:faa-multivariate} becomes
\begin{align*}
 \partial_1\partial_2\partial_3f(y)
 ={}&f'(y)\partial_{123}y\\
 &+f''(y)\bigl((\partial_1y)(\partial_{23}y)
 +(\partial_2y)(\partial_{13}y)+(\partial_3y)(\partial_{12}y)\bigr)\\
 &+f^{(3)}(y)(\partial_1y)(\partial_2y)(\partial_3y).
\end{align*}
For repeated variables or vector-valued inner maps, one labels derivative slots
first and then contracts tensor indices; the same partition lattice remains the
combinatorial skeleton.

\chapter{Composition of ordinary and exponential series}

\section{Ordinary formal power series}
Let
\[
 F(u)=\sum_{k\ge0}a_ku^k,
 \qquad
 G(x)=\sum_{j\ge1}b_jx^j,
 \qquad
 F(G(x))=\sum_{n\ge0}c_nx^n.
\]
\begin{theorem}[Ordinary composition]\label{thm:ordinary-composition}
For $n\ge1$,
\begin{align}
 c_n&=\sum_{k=1}^{n}a_k
 \sum_{\substack{i_1+\cdots+i_k=n\\i_r\ge1}}
 b_{i_1}\cdots b_{i_k},
 \label{eq:ordinary-composition-compositions}\\
 &=\sum_{k=1}^{n}a_k
 \sum_{\substack{\pi_1+\cdots+\pi_n=k\\
                   \pi_1+2\pi_2+\cdots+n\pi_n=n}}
 \binom{k}{\pi_1,\ldots,\pi_n}
 b_1^{\pi_1}\cdots b_n^{\pi_n}
 \label{eq:ordinary-composition-multiplicities}\\
 &=\sum_{k=1}^{n}a_k\OrdinaryPartialBellPolynomial nk(b_1,\ldots,b_{n-k+1}),
 \label{eq:ordinary-composition-bell}
\end{align}
and $c_0=a_0$.
\end{theorem}
\begin{proof}
The coefficient of $x^n$ in $G(x)^k$ is obtained by choosing an ordered
composition $i_1+\cdots+i_k=n$, proving the first formula.  Collect equal part
sizes; the number of orderings with multiplicities $\pi_j$ is the multinomial
coefficient, proving the second.  The third is the definition of the ordinary
Bell polynomial.
\end{proof}

Two useful specializations are immediate.  If $F(u)=\EulerE^u$ and the series are
exponentially normalized, one obtains the exponential formula.  If
$A(x)=1+\sum_{n\ge1}a_nx^n$, then
\begin{equation}\label{eq:reciprocal-ordinary-bell}
 [x^n]\frac1{A(x)}
 =\sum_{k=1}^{n}(-1)^k\OrdinaryPartialBellPolynomial nk(a_1,\ldots,a_{n-k+1})
 \qquad(n\ge1),
\end{equation}
by taking $F(u)=1/(1-u)$ and $G(x)=-\sum_{n\ge1}a_nx^n$.

\section{Exponential formal power series}
Let
\[
 F(x)=\sum_{j\ge1}a_j\frac{x^j}{j!},\qquad
 G(u)=\sum_{k\ge0}b_k\frac{u^k}{k!},\qquad
 G(F(x))=\sum_{n\ge0}c_n\frac{x^n}{n!}.
\]
\begin{theorem}[Exponential composition]\label{thm:exponential-composition}
\begin{equation}\label{eq:exponential-composition}
 c_n=\sum_{k=0}^{n}b_k\ExponentialPartialBellPolynomial nk(a_1,\ldots,a_{n-k+1}).
\end{equation}
Equivalently,
\begin{equation}\label{eq:exponential-composition-partitions}
 c_n=\sum_{\pi\in\SetPartitionOperator([n])}b_{|\pi|}\prod_{B\in\pi}a_{|B|}.
\end{equation}
\end{theorem}
\begin{proof}
Substitute $F(x)$ into $G$ and apply
\eqref{eq:partial-bell-egf}.  The set-partition form is
\cref{thm:bell-poly-partitions} with an additional weight $b_k$ for the number of
blocks.
\end{proof}

Taking $b_k=1$ gives the complete Bell-polynomial exponential identity
\begin{equation}\label{eq:exp-complete-bell-again}
 \exp\!\left(\sum_{j\ge1}a_j\frac{x^j}{j!}\right)
 =\sum_{n\ge0}\ExponentialCompleteBellPolynomial n(a_1,\ldots,a_n)\frac{x^n}{n!}.
\end{equation}
If the series converge, \eqref{eq:exponential-composition} is exactly
\eqref{eq:faa-bell} evaluated at the expansion point; formally, no convergence
notion is needed.

\chapter{Lagrange inversion and Lagrange--Bürmann}
\index{Lagrange inversion}

\section{Formal theorem}
Let $R$ be a commutative $\RationalNumbers$-algebra and let $\phi(w)\in R[[w]]$ have invertible
constant term.  There is a unique $g(z)\in zR[[z]]$ satisfying
\begin{equation}\label{eq:lagrange-functional}
 g(z)=z\phi(g(z)).
\end{equation}
Uniqueness follows recursively because the coefficient of $z^n$ on the right
depends only on coefficients of $g$ below degree $n$.

\begin{theorem}[Lagrange--Bürmann]\label{thm:lagrange-burmann}
For $n\ge1$ and every formal series $H$,
\begin{align}
 \CoefficientExtraction zn H(g(z))
 &=\frac1n\CoefficientExtraction w{n-1}H'(w)\phi(w)^n,
 \label{eq:lagrange-burmann}\\
 &=\CoefficientExtraction wn H(w)\phi(w)^{n-1}\bigl(\phi(w)-w\phi'(w)\bigr).
 \label{eq:lagrange-burmann-alt}
\end{align}
In particular,
\begin{equation}\label{eq:lagrange-basic}
 \CoefficientExtraction zn g(z)=\frac1n\CoefficientExtraction w{n-1}\phi(w)^n.
\end{equation}
\end{theorem}
\begin{proof}
Put $f(w)=w/\phi(w)$, so that $f(g(z))=z$.  Formal residues and the change of
variables theorem \cref{thm:res-subst} give
\begin{align*}
 \CoefficientExtraction zn H(g(z))
 &=\ResidueOperator_z H(g(z))z^{-n-1}\Differential z\\
 &=\ResidueOperator_w H(w)f(w)^{-n-1}f'(w)\Differential w\\
 &=-\frac1n\ResidueOperator_w H(w)(f(w)^{-n})'\Differential w\\
 &=\frac1n\ResidueOperator_w H'(w)f(w)^{-n}\Differential w,
\end{align*}
where the residue of a derivative is zero.  Since
$f(w)^{-n}=w^{-n}\phi(w)^n$, this is
\eqref{eq:lagrange-burmann}.  Applying the same zero-residue rule to
$H(w)\phi(w)^nw^{-n}$ gives
\[
 \frac1n\ResidueOperator H'\phi^nw^{-n}
 =\ResidueOperator H\phi^nw^{-n-1}-\ResidueOperator H\phi^{n-1}\phi'w^{-n},
\]
which is \eqref{eq:lagrange-burmann-alt}.  Take $H(w)=w$ for
\eqref{eq:lagrange-basic}.
\end{proof}

\section{Analytic inverse functions}
\begin{theorem}[Analytic Lagrange inversion]\label{thm:analytic-lagrange}
Let $f$ be analytic near $a$, let $b=f(a)$, and suppose $f'(a)\ne0$.  Its local
inverse $g$ is analytic near $b$ and has
\begin{equation}\label{eq:analytic-lagrange-series}
 g(z)=a+\sum_{n\ge1}g_n\frac{(z-b)^n}{n!},
\end{equation}
where
\begin{equation}\label{eq:analytic-lagrange-coeff}
 g_n=\left.\frac{\Differential^{n-1}}{\Differential w^{n-1}}
 \left(\frac{w-a}{f(w)-f(a)}\right)^n\right|_{w=a}.
\end{equation}
The series has a positive radius of convergence.
\end{theorem}
\begin{proof}
The analytic inverse-function theorem gives an analytic local inverse and therefore
a convergent Taylor series.  Put $u=w-a$, $\LagrangeShiftCoordinate=z-b$, and
\[
 \phi(u)=\frac{u}{f(a+u)-f(a)};
\]
its constant term is $1/f'(a)$.  The inverse displacement satisfies
$u=\LagrangeShiftCoordinate\phi(u)$.  Formula \eqref{eq:lagrange-basic}, multiplied by $n!$, is
exactly \eqref{eq:analytic-lagrange-coeff}.
\end{proof}

If $f'(a)=0$ but $f(w)-f(a)=c(w-a)^q+O((w-a)^{q+1})$ with $c\ne0$, factor
$f(w)-f(a)=(w-a)^qU(w)$ with $U(a)=c$.  Choosing one of the $q$ values of
$c^{-1/q}$ reduces the equation to ordinary inversion in the variable
$(z-f(a))^{1/q}$.  Thus the inverse has $q$ local Puiseux branches.  This is the
precise sense in which Lagrange inversion extends to a critical point; a
single-valued Taylor inverse cannot exist there.

\section{Inverse coefficients in Bell-polynomial form}
Suppose
\[
 f(w)=\sum_{j\ge1}f_j\frac{w^j}{j!},\qquad f_1\ne0,
 \qquad
 g(z)=f^{\langle-1\rangle}(z)=\sum_{n\ge1}g_n\frac{z^n}{n!}.
\]
Define
\begin{equation}\label{eq:fhat}
 \NormalizedReversionCoefficient{f}{j}=\frac{f_{j+1}}{(j+1)f_1}.
\end{equation}
\begin{theorem}[Explicit inverse coefficients]\label{thm:inverse-bell-coeff}
One has $g_1=f_1^{-1}$ and, for $n\ge2$,
\begin{equation}\label{eq:inverse-bell-coeff}
 g_n=\frac1{f_1^n}\sum_{k=1}^{n-1}(-1)^k\RisingFactorial nk
 \ExponentialPartialBellPolynomial{n-1}{k}(\NormalizedReversionCoefficient{f}{1},\ldots,\NormalizedReversionCoefficient{f}{n-k}).
\end{equation}
\end{theorem}
\begin{proof}
Write
\[
 \frac{f(w)}{f_1w}=1+U(w),
 \qquad U(w)=\sum_{j\ge1}\NormalizedReversionCoefficient{f}{j}\frac{w^j}{j!}.
\]
By \eqref{eq:lagrange-basic},
$g_n=(n-1)!\CoefficientExtraction w{n-1}(w/f(w))^n$.  Now
\[
 (1+U)^{-n}=\sum_{k\ge0}(-1)^k\binom{n+k-1}{k}U^k
 =\sum_{k\ge0}(-1)^k\frac{\RisingFactorial nk}{k!}U^k.
\]
Apply \eqref{eq:partial-bell-egf} and extract degree $n-1$; the $k=0$ term
vanishes when $n\ge2$.
\end{proof}

\section{The associahedron interpretation, with normalized notation}
The formula in the archive mixes exponential coefficients with an ordinary-series
face formula.  The exact geometric statement is the following.

Let
\begin{equation}\label{eq:ordinary-series-associa}
 A(x)=x+\sum_{n\ge1}a_nx^{n+1},\qquad
 B(x)=A^{\langle-1\rangle}(x)=x+\sum_{n\ge1}b_nx^{n+1}.
\end{equation}
Use the indexing in which the associahedron $\AssociahedronByDimension{n-1}$ has dimension $n-1$ and its
faces are plane rooted trees with $n+1$ leaves and internal arities at least $2$.
If an internal vertex has arity $i+1$, give it weight $a_i$; let $a_F$ be the
product of all vertex weights.  Equivalently, if
$F\cong \AssociahedronByDimension{i_1-1}\times\cdots\times \AssociahedronByDimension{i_m-1}$, then
$a_F=a_{i_1}\cdots a_{i_m}$.

\begin{theorem}[Associahedral inversion]\label{thm:associahedral-inversion}
\begin{equation}\label{eq:associahedral-inversion}
 b_n=\sum_{F\text{ a face of }\AssociahedronByDimension{n-1}}(-1)^{n-\dim F}a_F.
\end{equation}
\end{theorem}
\begin{proof}
The equation $A(B)=x$ is
\[
 B=x-\sum_{i\ge1}a_iB^{i+1}.
\]
Iterating it recursively expands $b_n$ as a sum over plane rooted trees with
$n+1$ leaves: every internal vertex of arity $i+1$ contributes $-a_i$.
A face $F$ of $\AssociahedronByDimension{n-1}$ is exactly such a tree after contracting a chosen set of
internal edges.  If it has $m$ internal vertices, then
$\dim F=n-m$, so its sign is $(-1)^m=(-1)^{n-\dim F}$ and its weight is $a_F$.
Thus the tree expansion is the alternating face sum.
\end{proof}

For example,
\begin{align*}
 b_1&=-a_1,\\
 b_2&=-a_2+2a_1^2,\\
 b_3&=-a_3+5a_2a_1-5a_1^3,\\
 b_4&=-a_4+6a_3a_1+3a_2^2-21a_2a_1^2+14a_1^4.
\end{align*}
For $\AssociahedronByDimension{3}$, the coefficients $1,6,3,21,14$ count respectively the whole
3-polytope, its six pentagonal and three square facets, its edges, and its
vertices.

\chapter{Algebraic and transcendental examples}

\section{The Fuss--Catalan algebraic equation}
Let $p\ge2$ be an integer and consider
\begin{equation}\label{eq:fuss-equation}
 x^p-x+z=0,
 \qquad x(0)=0.
\end{equation}
It is equivalent to
$x=z\phi(x)$ with $\phi(x)=(1-x^{p-1})^{-1}$.

\begin{theorem}\label{thm:fuss-series}
The distinguished solution is
\begin{equation}\label{eq:fuss-series}
 x(z)=\sum_{k=0}^{\infty}\binom{pk}{k}
       \frac{z^{(p-1)k+1}}{(p-1)k+1}.
\end{equation}
Its radius of convergence is
\begin{equation}\label{eq:fuss-radius}
 R_p=(p-1)p^{-p/(p-1)},
\end{equation}
and the series converges absolutely also on $|z|=R_p$.
\end{theorem}
\begin{proof}
By \eqref{eq:lagrange-basic},
\[
 \CoefficientExtraction znx(z)=\frac1n\CoefficientExtraction w{n-1}(1-w^{p-1})^{-n}.
\]
This vanishes unless $n-1=(p-1)k$.  In that case the coefficient is
$\frac1n\binom{n+k-1}{k}=\binom{pk}{k}/((p-1)k+1)$, proving the series.

The inverse map is $z=f(x)=x-x^p$.  Its nearest critical points satisfy
$f'(x)=1-px^{p-1}=0$ and all have modulus $p^{-1/(p-1)}$; their critical values
have modulus $(p-1)p^{-p/(p-1)}$.  No local inverse can pass analytically through
a critical value, while the inverse-function theorem continues the branch up to
one, so this is the radius.  Stirling's formula gives coefficients of order
$Ck^{-3/2}R_p^{-((p-1)k+1)}$, which proves absolute convergence on the boundary.
\end{proof}

\section{Lambert's \texorpdfstring{$W$}{W} function at the origin}
The Lambert function satisfies $W(z)\EulerE^{W(z)}=z$.
\begin{theorem}\label{thm:lambert-W-zero}
On its principal branch near $0$,
\begin{equation}\label{eq:lambert-W-zero}
 W(z)=\sum_{n=1}^{\infty}(-n)^{n-1}\frac{z^n}{n!}
 =z-z^2+\frac32z^3-\frac83z^4+O(z^5),
\end{equation}
with radius of convergence $\EulerE^{-1}$.
\end{theorem}
\begin{proof}
Here $f(w)=we^w$, so $w/f(w)=\EulerE^{-w}$.  Formula
\eqref{eq:analytic-lagrange-coeff} gives
$g_n=D^{n-1}\EulerE^{-nw}|_{w=0}=(-n)^{n-1}$.  The inverse map has its nearest critical
point at $w=-1$, where $f(-1)=-\EulerE^{-1}$; this square-root branch point determines
the radius.
\end{proof}

\section{A rapidly useful expansion around \texorpdfstring{$W(e)=1$}{W(e)=1}}
Put $\xi=\log z-1$ and $u=W(z)-1$.  Then
\begin{equation}\label{eq:lambert-near-e-implicit}
 u+\log(1+u)=\xi.
\end{equation}
Reversion gives
\begin{equation}\label{eq:lambert-near-e}
 W(\EulerE^{1+\xi})=1+\frac{\xi}{2}+\frac{\xi^2}{16}
 -\frac{\xi^3}{192}-\frac{\xi^4}{3072}
 +\frac{13\xi^5}{61440}+O(\xi^6).
\end{equation}
\begin{proof}
The derivative at $u=0$ of the left side of
\eqref{eq:lambert-near-e-implicit} is $2$, so it has an ordinary inverse.
Substitute $u=c_1\xi+\cdots+c_5\xi^5$ into
$u+\log(1+u)=\xi$ and compare successive powers.  The triangular equations give
$c_1=1/2,c_2=1/16,c_3=-1/192,c_4=-1/3072,c_5=13/61440$.
\end{proof}
The closest critical points of $u+\log(1+u)$ satisfy
$1+(1+u)^{-1}=0$, hence $u=-2$; using the two adjacent logarithm values gives
images $-2\pm \ImaginaryUnit\pi$.  Therefore the radius in the $\xi$-plane is
$\sqrt{4+\pi^2}$.  For positive real $z$, the expansion converges whenever
$|\log z-1|<\sqrt{4+\pi^2}$, approximately
$0.0655<z<112.63$.  Taking $\xi=-1$ yields a convergent evaluation of
$W(1)=0.567143\ldots$.

\section{Binary trees and Catalan numbers}
Let $\CatalanNumber{n}$ count rooted ordered full binary trees with $n$ internal vertices, and
let $B(z)=\sum_{n\ge0}\CatalanNumber{n}z^n$.  Removing the root gives
\begin{equation}\label{eq:binary-tree-functional}
 B(z)=1+zB(z)^2.
\end{equation}
Put $C(z)=B(z)-1$; then $C=z(1+C)^2$.  Lagrange inversion gives
\begin{equation}\label{eq:catalan-lagrange}
 \CatalanNumber{n}=\CoefficientExtraction znC(z)=\frac1n\CoefficientExtraction w{n-1}(1+w)^{2n}
 =\frac1n\binom{2n}{n-1}=\frac1{n+1}\binom{2n}{n}.
\end{equation}
This simultaneously proves the functional equation, its unique formal solution,
and the Catalan closed form.  The symbol $\CatalanNumber{n}$ is used here to avoid confusing
these binary-tree numbers with Bell numbers.

\chapter{Lagrange's reversion theorem in shifted form}

\section{The two-parameter implicit equation}
Let $v=v(x,y)$ be the unique formal solution of
\begin{equation}\label{eq:shifted-reversion-equation}
 v=x+yf(v),
 \qquad v(x,0)=x.
\end{equation}

\begin{theorem}[Lagrange reversion]\label{thm:shifted-lagrange-reversion}
For every formal or analytic $g$,
\begin{equation}\label{eq:shifted-lagrange-reversion}
 g(v)=g(x)+\sum_{k=1}^{\infty}\frac{y^k}{k!}
 \left(\frac{\partial}{\partial x}\right)^{k-1}
 \bigl(f(x)^kg'(x)\bigr).
\end{equation}
In particular,
\begin{equation}\label{eq:shifted-lagrange-v}
 v=x+\sum_{k=1}^{\infty}\frac{y^k}{k!}
 \left(\frac{\partial}{\partial x}\right)^{k-1}f(x)^k.
\end{equation}
\end{theorem}
\begin{proof}
Put $u=v-x$.  Then $u=y\phi(u)$ with $\phi(u)=f(x+u)$, treating $x$ as a
coefficient parameter.  Apply \eqref{eq:lagrange-burmann} to
$H(u)=g(x+u)$:
\[
 [y^k]g(v)=\frac1k[u^{k-1}]g'(x+u)f(x+u)^k
 =\frac1{k!}D_x^{k-1}(g'(x)f(x)^k).
\]
Summation proves the first formula; take $g(z)=z$ for the second.
\end{proof}

\section{The delta-function proof decoded}
The archive also gives a distributional derivation.  Assume temporarily that
$x,y\in\RealNumbers$, the equation $z-x-yf(z)=0$ has one simple root $v$, and
$1-yf'(v)>0$.  The elementary change-of-variables rule for the Dirac delta gives
\begin{equation}\label{eq:delta-root}
 g(v)=\int_{\RealNumbers}\delta(yf(z)-z+x)g(z)(1-yf'(z))\,\Differential z.
\end{equation}
Using
$\delta(s)=\int_{\RealNumbers}\EulerE^{iks}\Differential k/(2\pi)$ and expanding only in the formal
parameter $y$,
\begin{align}
 g(v)
 &=\sum_{n\ge0}D_x^n\left[
 \frac{y^nf(x)^ng(x)}{n!}(1-yf'(x))\right]\notag\\
 &=\sum_{n\ge0}D_x^n\left[
 \frac{y^nf(x)^ng(x)}{n!}
 -\frac{y^{n+1}}{(n+1)!}\bigl((gf^{n+1})'-g'f^{n+1}\bigr)
 \right].
 \label{eq:delta-reversion-middle}
\end{align}
The first term of level $n+1$ cancels the total derivative in the second term of
level $n$; what remains is exactly
\eqref{eq:shifted-lagrange-reversion}.  Thus the Fourier--delta calculation is a
compact analytic avatar of the formal residue proof.  Formula
\eqref{eq:delta-root} requires an absolute Jacobian in general; the displayed sign
is valid in the stated orientation-preserving neighbourhood.

\chapter{Laplace-type asymptotic integrals}
\index{Laplace asymptotics}

\section{Endpoint expansion}
Consider
\begin{equation}\label{eq:laplace-integral}
 I(\lambda)=\int_a^b \EulerE^{-\lambda f(x)}g(x)\,\Differential x,
 \qquad \lambda\to+\infty,
\end{equation}
where the endpoint $a$ is the unique relevant minimum and, with $t=x-a$,
\begin{align}
 f(a+t)&\sim f(a)+\sum_{k\ge0}a_kt^{k+\alpha},
 &a_0>0,\quad\alpha>0,
 \label{eq:laplace-f-expansion}\\
 g(a+t)&\sim\sum_{k\ge0}b_kt^{k+\beta-1},
 &\beta>0.
 \label{eq:laplace-g-expansion}
\end{align}
Let $\OrdinaryPartialBellPolynomial{k}{j}$ denote the ordinary Bell polynomial.

\begin{theorem}[Laplace--Erd\'elyi coefficient formula]
\label{thm:laplace-bell}
Under the standard differentiability and remainder hypotheses that permit Watson's
lemma,
\begin{equation}\label{eq:laplace-asymptotic}
 I(\lambda)\sim \EulerE^{-\lambda f(a)}
 \sum_{n=0}^{\infty}
 \Gamma\!\left(\frac{n+\beta}{\alpha}\right)
 \frac{c_n}{\lambda^{(n+\beta)/\alpha}},
\end{equation}
where
\begin{equation}\label{eq:laplace-cn}
 c_n=\frac1{\alpha a_0^{(n+\beta)/\alpha}}
 \sum_{k=0}^{n}b_{n-k}
 \sum_{j=0}^{k}
 \binom{-\frac{n+\beta}{\alpha}}{j}
 \frac{\OrdinaryPartialBellPolynomial{k}{j}(a_1,\ldots,a_{k-j+1})}{a_0^j}.
\end{equation}
\end{theorem}
\begin{proof}
Write
$f(a+t)-f(a)=t^\alpha A(t)$ with
$A(t)=a_0+a_1t+\cdots$, and set
$s=(f(a+t)-f(a))^{1/\alpha}=tA(t)^{1/\alpha}$.  This has an inverse $t=t(s)$.
Expand
\[
 g(a+t(s))t'(s)=\sum_{n\ge0}d_ns^{n+\beta-1}.
\]
A generalized Lagrange--Bürmann calculation, first for integral $\beta$ and then
by formal continuation in the exponent, gives
\[
 d_n=\sum_{k=0}^{n}b_{n-k}[t^k]A(t)^{-(n+\beta)/\alpha}.
\]
Indeed, apply \eqref{eq:lagrange-burmann} to an antiderivative of
$t^{m+\beta-1}$ and differentiate back after substitution.

Put $\rho=(n+\beta)/\alpha$.  The ordinary Bell expansion yields
\begin{align*}
 [t^k]A(t)^{-\rho}
 &=a_0^{-\rho}[t^k]
 \sum_{j\ge0}\binom{-\rho}{j}
 \left(\sum_{r\ge1}\frac{a_r}{a_0}t^r\right)^j\\
 &=a_0^{-\rho}\sum_{j=0}^{k}\binom{-\rho}{j}a_0^{-j}
 \OrdinaryPartialBellPolynomial{k}{j}(a_1,\ldots,a_{k-j+1}).
\end{align*}
Thus $c_n=d_n/\alpha$ is \eqref{eq:laplace-cn}.  Finally
\[
 \int_0^\infty \EulerE^{-\lambda s^\alpha}s^{n+\beta-1}\Differential s
 =\frac1\alpha\Gamma\!\left(\frac{n+\beta}{\alpha}\right)
 \lambda^{-(n+\beta)/\alpha},
\]
and Watson's lemma controls the replacement of the finite endpoint interval by
the asymptotic expansion.
\end{proof}

\begin{remark}
The appearance of ordinary rather than exponential Bell polynomials is forced by
the ordinary Taylor expansion of $A(t)^{-\rho}$.  Lagrange inversion enters in the
change from $t$ to the natural Laplace coordinate $s$; the gamma factor enters only
in the final monomial integral.
\end{remark}

\part{Factorial, Umbral, and Special-Function Extensions}

\chapter{Pochhammer symbols, binomial transforms, and factorial calculus}
\label{ch:ext-pochhammer}
\index{Pochhammer symbol}

The five drafts use several mutually compatible factorial conventions.  This
chapter places them in one calculus and makes explicit the binomial inversion
machinery that is repeatedly used later.  The notation and analytic continuation
conventions agree with the standard special-function references
\cite{NISTDLMF,WhittakerWatson1927}.  We retain
\(\FallingFactorial{x}{n}\) and \(\RisingFactorial{x}{n}\) for falling and rising factorials;
when a formula is standardly written with the Pochhammer symbol, we set
\begin{equation}\label{eq:ext-pochhammer-def}
 (x)_n:=\RisingFactorial{x}{n}=x(x+1)\cdots(x+n-1),\qquad (x)_0=1.
\end{equation}
Thus \(\FallingFactorial{x}{n}=(-1)^n(-x)_n\).

\section{Elementary identities and analytic continuation}

\begin{proposition}[Shift, splitting, and gamma quotient]
For nonnegative integers \(m,n\),
\begin{align}
 (x)_{m+n}&=(x)_m(x+m)_n,\label{eq:ext-pochhammer-split}\\
 \FallingFactorial{x}{m+n}&=\FallingFactorial{x}{m}\FallingFactorial{x-m}{n},\label{eq:ext-falling-split}\\
 (x)_n&=\frac{\Gamma(x+n)}{\Gamma(x)}\label{eq:ext-pochhammer-gamma}
\end{align}
whenever neither side of \eqref{eq:ext-pochhammer-gamma} has an uncancelled
pole.  The quotient on the right gives the meromorphic continuation of
\((x)_n\) in \(x\).
\end{proposition}
\begin{proof}
Split the defining product for \((x)_{m+n}\) after its first \(m\) factors;
the remaining factors are \(x+m,x+m+1,\ldots,x+m+n-1\).  The falling-factorial
identity is identical after reversing the direction of the shifts.  Repeatedly
applying \(\Gamma(z+1)=z\Gamma(z)\) gives
\(\Gamma(x+n)=(x+n-1)\cdots x\,\Gamma(x)\).  Division proves
\eqref{eq:ext-pochhammer-gamma} away from poles, and meromorphic continuation
extends the identity wherever the quotient is defined.
\end{proof}

\begin{proposition}[Vandermonde convolution for rising factorials]
For every \(n\ge0\),
\begin{equation}\label{eq:ext-vandermonde-rising}
 (x+y)_n=\sum_{k=0}^n\binom nk(x)_k(y)_{n-k}.
\end{equation}
Equivalently, for generalized binomial coefficients,
\begin{equation}\label{eq:ext-chu-vandermonde}
 \binom{x+y}{n}=\sum_{k=0}^n\binom{x}{k}\binom{y}{n-k}.
\end{equation}
\end{proposition}
\begin{proof}
The generalized binomial theorem, proved below, gives
\[
 (1-z)^{-x}(1-z)^{-y}=(1-z)^{-(x+y)}.
\]
The coefficient of \(z^n\) on the left is
\(\sum_{k=0}^n(x)_k(y)_{n-k}/(k!(n-k)!)\), whereas on the right it is
\((x+y)_n/n!\).  Multiplication by \(n!\) proves
\eqref{eq:ext-vandermonde-rising}.  Replacing \(x,y\) by \(-x,-y\), using
\(\binom{x}{n}=\FallingFactorial{x}{n}/n!=(-1)^n(-x)_n/n!\), gives
\eqref{eq:ext-chu-vandermonde}.
\end{proof}

\section{The generalized binomial theorem}

\begin{theorem}[Formal generalized binomial theorem]
Let \(R\) be a commutative \(\RationalNumbers\)-algebra and let \(\alpha\in R\).  In
\(R[[z]]\),
\begin{align}
 (1+z)^\alpha
 &=\sum_{n\ge0}\binom{\alpha}{n}z^n
 =\sum_{n\ge0}\frac{\FallingFactorial{\alpha}{n}}{n!}z^n,\label{eq:ext-binomial-plus}\\
 (1-z)^{-\alpha}
 &=\sum_{n\ge0}\frac{(\alpha)_n}{n!}z^n.\label{eq:ext-binomial-minus}
\end{align}
Analytically, the same series converge for \(|z|<1\), with any consistent
choice of the logarithm used to define the powers.
\end{theorem}
\begin{proof}
Let
\(F(z)=\sum_{n\ge0}\FallingFactorial{\alpha}{n}z^n/n!\).  Since
\(\FallingFactorial{\alpha}{n+1}=(\alpha-n)\FallingFactorial{\alpha}{n}\), coefficient
comparison gives
\[
 (1+z)F'(z)=\alpha F(z),\qquad F(0)=1.
\]
The formal differential equation has a unique solution: once the constant term
is fixed, its coefficient of degree \(n+1\) is determined by coefficients of
lower degree.  The formal exponential
\(\exp(\alpha\log(1+z))\) satisfies the same equation and initial condition,
so it equals \(F\).  This proves \eqref{eq:ext-binomial-plus}; replacing
\(z\) by \(-z\) and \(\alpha\) by \(-\alpha\) proves
\eqref{eq:ext-binomial-minus}.  In the analytic setting, the ratio test gives
radius one unless the series terminates, and differentiation identifies the
sum on every simply connected subdomain of the unit disk.
\end{proof}

\section{Binomial inversion and finite differences}
\index{binomial inversion}

\begin{theorem}[Binomial inversion]
\label{thm:ext-binomial-inversion}
For sequences \((a_n)\) and \((b_n)\) in a commutative ring, the two systems
\begin{align}
 b_n&=\sum_{k=0}^n\binom nk a_k,\label{eq:ext-binomial-transform}\\
 a_n&=\sum_{k=0}^n(-1)^{n-k}\binom nk b_k\label{eq:ext-binomial-inverse}
\end{align}
are equivalent.
\end{theorem}
\begin{proof}
Assume \eqref{eq:ext-binomial-transform} and substitute it into the right side
of \eqref{eq:ext-binomial-inverse}.  Interchanging the two finite sums gives
\begin{align*}
 \sum_{k=0}^n(-1)^{n-k}\binom nk b_k
 &=\sum_{j=0}^n a_j\sum_{k=j}^n
   (-1)^{n-k}\binom nk\binom kj\\
 &=\sum_{j=0}^n a_j\binom nj
   \sum_{r=0}^{n-j}(-1)^{n-j-r}\binom{n-j}{r}.
\end{align*}
The inner sum is \((1-1)^{n-j}\), hence is one for \(j=n\) and zero otherwise.
The result is \(a_n\).  The same calculation in the opposite order proves the
converse, because the inverse matrix is its own inverse after the alternating
signs are inserted.
\end{proof}

\begin{corollary}[Generating-function forms]
Let \(A(z)=\sum_{n\ge0}a_nz^n\) and \(B(z)=\sum_{n\ge0}b_nz^n\).  Then
\eqref{eq:ext-binomial-transform} is equivalent to
\begin{equation}\label{eq:ext-binomial-ogf}
 B(z)=\frac1{1-z}A\!\left(\frac{z}{1-z}\right),
 \qquad
 A(z)=\frac1{1+z}B\!\left(\frac{z}{1+z}\right).
\end{equation}
For exponential generating functions
\(\ExponentialGeneratingFunctionOf{A}(z)=\sum a_nz^n/n!\) and
\(\ExponentialGeneratingFunctionOf{B}(z)=\sum b_nz^n/n!\), the equivalent relation is
\begin{equation}\label{eq:ext-binomial-egf}
 \ExponentialGeneratingFunctionOf{B}(z)=\EulerE^z\ExponentialGeneratingFunctionOf{A}(z),\qquad
 \ExponentialGeneratingFunctionOf{A}(z)=\EulerE^{-z}\ExponentialGeneratingFunctionOf{B}(z).
\end{equation}
\end{corollary}
\begin{proof}
For ordinary generating functions, sum \eqref{eq:ext-binomial-transform} over
\(n\) and use
\(\sum_{n\ge k}\binom nkz^n=z^k/(1-z)^{k+1}\).  Solving the fractional-linear
substitution for its inverse gives the second formula.  For exponential
generating functions,
\[
 \sum_{n\ge0}\sum_{k=0}^n\binom nk a_k\frac{z^n}{n!}
 =\left(\sum_{k\ge0}a_k\frac{z^k}{k!}\right)
  \left(\sum_{j\ge0}\frac{z^j}{j!}\right),
\]
which is \eqref{eq:ext-binomial-egf}.
\end{proof}

\begin{proposition}[Newton series and binomial inversion]
For every polynomial \(p\) of degree at most \(d\),
\begin{equation}\label{eq:ext-newton-binomial}
 p(x)=\sum_{k=0}^d\binom{x}{k}\Delta^kp(0),
 \qquad
 \Delta^kp(0)=\sum_{j=0}^k(-1)^{k-j}\binom kjp(j).
\end{equation}
\end{proposition}
\begin{proof}
The first statement is the Newton expansion already proved in
\cref{thm:newton-expansion}.  Evaluating it at integers \(x=n\) gives
\(p(n)=\sum_{k=0}^n\binom nk\Delta^kp(0)\).  Applying
\cref{thm:ext-binomial-inversion} yields the second identity.  Conversely, the
second identity and binomial inversion recover the first at all nonnegative
integers; two polynomials of degree at most \(d\) agreeing at \(d+1\) points are
equal.
\end{proof}

\section{Factorial differential and difference rules}

\begin{proposition}
For \(n\ge1\),
\begin{align}
 \Delta\FallingFactorial{x}{n}&=n\FallingFactorial{x}{n-1},\label{eq:ext-delta-falling}\\
 \nabla\RisingFactorial{x}{n}&=n\RisingFactorial{x}{n-1},\label{eq:ext-nabla-rising}\\
 \frac{\Differential}{\Differential x}(x)_n
 &=(x)_n\sum_{j=0}^{n-1}\frac1{x+j}
 =(x)_n\bigl(\psi(x+n)-\psi(x)\bigr),\label{eq:ext-pochhammer-derivative}
\end{align}
where \(\nabla f(x)=f(x)-f(x-1)\), and the last identity holds away from poles.
\end{proposition}
\begin{proof}
Factor the two neighboring products:
\[
 \FallingFactorial{x+1}{n}-\FallingFactorial{x}{n}
 =\FallingFactorial{x}{n-1}\bigl((x+1)-(x-n+1)\bigr)
 =n\FallingFactorial{x}{n-1}.
\]
The rising-factorial statement is the reflected identity.  Logarithmic
differentiation of \((x)_n=\prod_{j=0}^{n-1}(x+j)\) gives the finite sum.
Logarithmic differentiation of \eqref{eq:ext-pochhammer-gamma} gives the
digamma difference.
\end{proof}

\begin{remark}
The formulas in this chapter explain why factorial bases are the natural
coordinates for finite-difference operators, just as monomials are natural for
differentiation.  Stirling numbers are precisely the transition matrices
between these two coordinate systems.
\end{remark}

\chapter{Bernoulli, Euler, Genocchi, and Cauchy polynomial families}
\label{ch:ext-bernoulli}
\index{Bernoulli polynomials}
\index{Euler polynomials}
\index{Genocchi numbers}

To avoid collision with Bell numbers, Bernoulli numbers remain \(\beta_n\), and
we write \(\mathsf B_n(x)\) for Bernoulli polynomials, \(\mathsf E_n(x)\) for
Euler polynomials, \(\mathsf \GenocchiNumber{n}\) for Genocchi numbers, and
\(\mathsf C_n(x)\) for the Bernoulli--Cauchy polynomials of the second kind, following the classical difference-calculus tradition \cite{Norlund1924}.

\section{Bernoulli polynomials as an Appell sequence}

\begin{definition}
The Bernoulli polynomials are defined by
\begin{equation}\label{eq:ext-bernoulli-poly-egf}
 \frac{t \EulerE^{xt}}{\EulerE^t-1}
 =\sum_{n\ge0}\mathsf B_n(x)\frac{t^n}{n!}.
\end{equation}
Thus \(\mathsf B_n(0)=\beta_n\).
\end{definition}

\begin{theorem}[Appell, translation, difference, and reflection identities]
For \(n\ge1\),
\begin{align}
 \mathsf B_n'(x)&=n\mathsf B_{n-1}(x),\label{eq:ext-bernoulli-derivative}\\
 \mathsf B_n(x+y)&=\sum_{k=0}^n\binom nk\mathsf B_k(x)y^{n-k},\label{eq:ext-bernoulli-translation}\\
 \mathsf B_n(x+1)-\mathsf B_n(x)&=nx^{n-1},\label{eq:ext-bernoulli-difference}\\
 \mathsf B_n(1-x)&=(-1)^n\mathsf B_n(x).\label{eq:ext-bernoulli-reflection}
\end{align}
Moreover,
\begin{equation}\label{eq:ext-bernoulli-explicit}
 \mathsf B_n(x)=\sum_{k=0}^n\binom nk\beta_k x^{n-k}.
\end{equation}
\end{theorem}
\begin{proof}
Differentiate \eqref{eq:ext-bernoulli-poly-egf} with respect to \(x\); the
result is \(t\) times the same generating function, and coefficient comparison
gives \eqref{eq:ext-bernoulli-derivative}.  Multiplying the generating function
by \(\EulerE^{yt}\) proves \eqref{eq:ext-bernoulli-translation}; taking \(x=0\)
gives \eqref{eq:ext-bernoulli-explicit}.  For the finite difference,
\[
 \frac{t \EulerE^{(x+1)t}}{\EulerE^t-1}-\frac{t \EulerE^{xt}}{\EulerE^t-1}=t \EulerE^{xt},
\]
whose coefficient of \(t^n/n!\) is \(nx^{n-1}\).  Finally,
\[
 \frac{t \EulerE^{(1-x)t}}{\EulerE^t-1}
 =\frac{(-t)\EulerE^{x(-t)}}{\EulerE^{-t}-1};
\]
replace \(t\) by \(-t\) in the defining series to obtain
\eqref{eq:ext-bernoulli-reflection}.
\end{proof}

\begin{corollary}[Faulhaber's formula]
For integers \(N\ge1\) and \(p\ge0\),
\begin{equation}\label{eq:ext-faulhaber}
 \sum_{m=0}^{N-1}m^p
 =\frac{\mathsf B_{p+1}(N)-\mathsf B_{p+1}(0)}{p+1}
 =\frac1{p+1}\sum_{k=0}^{p}\binom{p+1}{k}\beta_kN^{p+1-k}.
\end{equation}
\end{corollary}
\begin{proof}
Set \(n=p+1\) in \eqref{eq:ext-bernoulli-difference}, substitute
\(x=m\), and sum from \(m=0\) to \(N-1\).  The left side telescopes.
Formula \eqref{eq:ext-bernoulli-explicit} gives the second equality; the
\(k=p+1\) term cancels in the difference.
\end{proof}

\section{Euler polynomials and alternating sums}

\begin{definition}
The Euler polynomials are defined by
\begin{equation}\label{eq:ext-euler-poly-egf}
 \frac{2\EulerE^{xt}}{\EulerE^t+1}
 =\sum_{n\ge0}\mathsf E_n(x)\frac{t^n}{n!}.
\end{equation}
\end{definition}

\begin{theorem}
For \(n\ge1\),
\begin{align}
 \mathsf E_n'(x)&=n\mathsf E_{n-1}(x),\label{eq:ext-euler-derivative}\\
 \mathsf E_n(x+1)+\mathsf E_n(x)&=2x^n,\label{eq:ext-euler-difference}\\
 \mathsf E_n(1-x)&=(-1)^n\mathsf E_n(x),\label{eq:ext-euler-reflection}\\
 \mathsf E_n(x+y)&=\sum_{k=0}^n\binom nk\mathsf E_k(x)y^{n-k}.
 \label{eq:ext-euler-translation}
\end{align}
Consequently, for \(N\ge1\),
\begin{equation}\label{eq:ext-alternating-power-sum}
 \sum_{m=0}^{N-1}(-1)^m m^n
 =\frac{\mathsf E_n(0)-(-1)^N\mathsf E_n(N)}2.
\end{equation}
\end{theorem}
\begin{proof}
The derivative and translation identities follow exactly as for Bernoulli
polynomials.  Multiplying \eqref{eq:ext-euler-poly-egf} by \(\EulerE^t+1\) gives
\[
 \sum_{n\ge0}\bigl(\mathsf E_n(x+1)+\mathsf E_n(x)\bigr)\frac{t^n}{n!}
 =2\EulerE^{xt},
\]
which proves \eqref{eq:ext-euler-difference}.  Replacing \(t\) by \(-t\)
proves reflection.  Multiply \eqref{eq:ext-euler-difference} at \(x=m\) by
\((-1)^m\) and sum.  Adjacent terms cancel, leaving
\(\mathsf E_n(0)-(-1)^N\mathsf E_n(N)\), and division by two gives
\eqref{eq:ext-alternating-power-sum}.
\end{proof}

\section{Genocchi numbers}

\begin{definition}
The Genocchi numbers \(\mathsf \GenocchiNumber{n}\) are defined by
\begin{equation}\label{eq:ext-genocchi-egf}
 \frac{2t}{\EulerE^t+1}=\sum_{n\ge0}\mathsf \GenocchiNumber{n}\frac{t^n}{n!}.
\end{equation}
\end{definition}

\begin{proposition}[Reduction to Bernoulli numbers]
For every \(n\ge0\),
\begin{equation}\label{eq:ext-genocchi-bernoulli}
 \mathsf \GenocchiNumber{n}=2(1-2^n)\beta_n.
\end{equation}
In particular, \(\mathsf \GenocchiNumber{0}=0\), \(\mathsf \GenocchiNumber{1}=1\), and
\(\mathsf \GenocchiNumber{2m+1}=0\) for \(m\ge1\).
\end{proposition}
\begin{proof}
The algebraic identity
\[
 \frac{2t}{\EulerE^t+1}=\frac{2t}{\EulerE^t-1}-\frac{4t}{\EulerE^{2t}-1}
\]
expresses the left side as
\[
 2\sum_{n\ge0}\beta_n\frac{t^n}{n!}
 -2\sum_{n\ge0}\beta_n\frac{(2t)^n}{n!}.
\]
Coefficient comparison proves \eqref{eq:ext-genocchi-bernoulli}.  The stated
vanishing follows from \(\beta_{2m+1}=0\) for \(m\ge1\), which itself follows
from Bernoulli reflection at \(x=0,1\) together with
\eqref{eq:ext-bernoulli-difference}.
\end{proof}

\section{Bernoulli--Cauchy polynomials of the second kind}
\index{Cauchy polynomials}

\begin{definition}
Define \(\mathsf C_n(x)\) by
\begin{equation}\label{eq:ext-cauchy-poly-egf}
 \frac{t(1+t)^x}{\log(1+t)}
 =\sum_{n\ge0}\mathsf C_n(x)\frac{t^n}{n!}.
\end{equation}
These are also called Bernoulli polynomials of the second kind.  We use a
distinct letter to prevent confusion with \(\mathsf B_n(x)\).
\end{definition}

\begin{theorem}[Stirling expansion, integral representation, and reflection]
For \(n\ge1\),
\begin{align}
 \mathsf C_n(x)&=\mathsf C_n(0)
 +\sum_{k=1}^n\frac{n}{k}\SignedStirlingFirstKind{n-1}{k-1}x^k,
 \label{eq:ext-cauchy-stirling}\\
 \mathsf C_n(x)&=\int_0^1\FallingFactorial{x+u}{n}\,\Differential u.
 \label{eq:ext-cauchy-integral}
\end{align}
Moreover,
\begin{equation}\label{eq:ext-cauchy-reflection}
 \mathsf C_n(n-2-x)=(-1)^n\mathsf C_n(x).
\end{equation}
Equivalently, the graph is even or odd about \((n-2)/2\) according as \(n\)
is even or odd.
\end{theorem}
\begin{proof}
Differentiate \eqref{eq:ext-cauchy-poly-egf} with respect to \(x\):
\[
 \sum_{n\ge0}\mathsf C_n'(x)\frac{t^n}{n!}
 =t(1+t)^x.
\]
Using
\((1+t)^x=\sum_{m\ge0}\FallingFactorial{x}{m}t^m/m!\), comparison gives
\begin{equation}\label{eq:ext-cauchy-derivative}
 \mathsf C_n'(x)=n\FallingFactorial{x}{n-1}.
\end{equation}
Expand \(\FallingFactorial{x}{n-1}=\sum_{j=0}^{n-1}\SignedStirlingFirstKind{n-1}{j}x^j\) and integrate from
zero to \(x\).  This proves \eqref{eq:ext-cauchy-stirling} after putting
\(k=j+1\).

For the integral representation, observe directly that
\[
 \int_0^1(1+t)^{x+u}\Differential u
 =(1+t)^x\frac{(1+t)-1}{\log(1+t)}
 =\frac{t(1+t)^x}{\log(1+t)}.
\]
Expanding \((1+t)^{x+u}\) by falling factorials and comparing coefficients
proves \eqref{eq:ext-cauchy-integral}.  Finally, use the elementary reflection
\(\FallingFactorial{n-1-y}{n}=(-1)^n\FallingFactorial{y}{n}\):
\begin{align*}
 \mathsf C_n(n-2-x)
 &=\int_0^1\FallingFactorial{n-2-x+u}{n}\Differential u\\
 &=\int_0^1\FallingFactorial{n-1-(x+1-u)}{n}\Differential u\\
 &=(-1)^n\int_0^1\FallingFactorial{x+1-u}{n}\Differential u
 =(-1)^n\mathsf C_n(x),
\end{align*}
where the last equality uses the substitution \(v=1-u\).  This proves
\eqref{eq:ext-cauchy-reflection} without an undetermined integration constant.
\end{proof}

\section{Complementary Bell numbers}
\index{complementary Bell numbers}

\begin{definition}
The complementary Bell numbers are the signed partition sums
\begin{equation}\label{eq:ext-complementary-bell-def}
 \ComplementaryBellNumber{n}:=\sum_{k=0}^n(-1)^k\StirlingSecondKind nk.
\end{equation}
\end{definition}

\begin{theorem}[Generating function and signed Dobi\'{n}ski formula]
One has
\begin{align}
 \sum_{n\ge0}\ComplementaryBellNumber{n}\frac{z^n}{n!}
 &=\exp(1-\EulerE^z),\label{eq:ext-complementary-bell-egf}\\
 \ComplementaryBellNumber{n}
 &=\EulerE\sum_{m\ge0}\frac{(-1)^m m^n}{m!}.\label{eq:ext-complementary-dobinski}
\end{align}
The second series converges absolutely for each fixed \(n\).
\end{theorem}
\begin{proof}
Insert the Stirling EGF \eqref{eq:second-egf} into
\eqref{eq:ext-complementary-bell-def} and interchange the locally finite
formal sums:
\[
 \sum_{n\ge0}\ComplementaryBellNumber{n}\frac{z^n}{n!}
 =\sum_{k\ge0}\frac{(-1)^k(\EulerE^z-1)^k}{k!}
 =\EulerE^{-(\EulerE^z-1)}.
\]
This is \eqref{eq:ext-complementary-bell-egf}.  Expanding the outer exponential
in the equivalent form \(e\,\EulerE^{-\EulerE^z}\) gives
\[
 \EulerE\sum_{m\ge0}\frac{(-1)^m}{m!}\EulerE^{mz}.
\]
The coefficient of \(z^n/n!\) is \eqref{eq:ext-complementary-dobinski}.
Absolute convergence follows from the ratio test applied to \(m^n/m!\).
\end{proof}

\section{Modified Bernoulli numbers from a logarithmic hyperbolic series}

\begin{proposition}
Define \(b_{2m}\) by
\begin{equation}\label{eq:ext-modified-bernoulli-egf}
 \frac12\log\!\left(\frac{\EulerE^{x/2}-\EulerE^{-x/2}}{x/2}\right)
 =\frac12\log2+\sum_{m\ge1}b_{2m}x^{2m}.
\end{equation}
Then
\begin{equation}\label{eq:ext-modified-bernoulli-closed}
 b_{2m}=\frac{\beta_{2m}}{4m(2m)!},\qquad b_{2m+1}=0.
\end{equation}
In particular,
\[
 b_2=\frac1{48},\qquad b_4=-\frac1{5760},\qquad
 b_6=\frac1{362880}.
\]
\end{proposition}
\begin{proof}
After removing the constant \(\tfrac12\log2\), differentiate the left side:
\[
 \frac14\coth(x/2)-\frac1{2x}.
\]
From the Bernoulli generating function one obtains
\[
 \coth(x/2)=\frac2x+2\sum_{m\ge1}\beta_{2m}
 \frac{x^{2m-1}}{(2m)!}.
\]
Substitution cancels the pole and leaves
\(\tfrac12\sum_{m\ge1}\beta_{2m}x^{2m-1}/(2m)!\).
Termwise integration, with zero constant after the explicit
\(\tfrac12\log2\), gives \eqref{eq:ext-modified-bernoulli-closed}.  Odd powers
are absent because the logarithm after removal of the constant is an even
function.
\end{proof}

\chapter{Sheffer sequences and exponential Riordan arrays}
\label{ch:ext-sheffer}
\index{Sheffer sequence}
\index{Riordan array!exponential}

Polynomial sequences of binomial type already arose from complete Bell
polynomials in \cref{thm:binomial-type-bell}.  The natural closure of that
class under multiplication by an invertible exponential series is the class of
Sheffer sequences.  Exponential Riordan arrays package the same operations as
lower-triangular matrix multiplication \cite{Roman1984}.

\section{Characterization of sequences of binomial type}

\begin{theorem}[Generating-function characterization]
\label{thm:ext-binomial-type-characterization}
Let \(p_n(x)\) be polynomials over a field of characteristic zero, with
\(p_0(x)=1\) and \(\deg p_n\le n\).  The following are equivalent.
\begin{enumerate}
\item For all \(x,y\),
\begin{equation}\label{eq:ext-binomial-type-law}
 p_n(x+y)=\sum_{k=0}^n\binom nkp_k(x)p_{n-k}(y).
\end{equation}
\item There is a unique series
\(\beta(z)=\beta_1z+\beta_2z^2/2!+\cdots\) such that
\begin{equation}\label{eq:ext-binomial-type-general-egf}
 \sum_{n\ge0}p_n(x)\frac{z^n}{n!}=\exp(x\beta(z)).
\end{equation}
\end{enumerate}
Moreover, \(\deg p_n=n\) for every \(n\) if and only if
\(\beta_1=\beta'(0)\ne0\); in that case \(\beta\) is a delta series in the
usual umbral-calculus terminology.
\end{theorem}
\begin{proof}
Assume \eqref{eq:ext-binomial-type-general-egf}.  Replacing \(x\) by \(x+y\)
and factoring the exponential yields a product of two EGFs; the EGF product
rule \eqref{eq:egf-product} gives \eqref{eq:ext-binomial-type-law}.

Conversely, put
\(P(x,z)=\sum_{n\ge0}p_n(x)z^n/n!\).  The binomial identity is exactly
\begin{equation}\label{eq:ext-cauchy-exponential-equation}
 P(x+y,z)=P(x,z)P(y,z).
\end{equation}
Setting \(x=y=0\) gives \(P(0,z)=P(0,z)^2\).  Because \(P(0,z)\) has
constant term one, it is invertible, and hence \(P(0,z)=1\).  The formal
logarithm is therefore defined.  Let \(L(x,z)=\log P(x,z)\).  Equation
\eqref{eq:ext-cauchy-exponential-equation} becomes
\(L(x+y,z)=L(x,z)+L(y,z)\).  Write
\(L(x,z)=\sum_{m\ge0}\ell_m(x)z^m\).  Each \(\ell_m(x)\) is a polynomial
satisfying the additive Cauchy equation; comparing coefficients in
\(\ell_m(x+y)\) shows that every term of degree other than one vanishes.
Hence \(\ell_m(x)=c_mx\), so \(L(x,z)=x\beta(z)\).  Since the
constant term of \(P(x,z)\) in \(z\) is \(p_0(x)=1\), one has
\(\beta(0)=0\).  Uniqueness follows from
\(\beta(z)=\partial_xP(0,z)\).  If \(\beta_1\ne0\), the contribution
\((x\beta_1z)^n/n!\) shows that the coefficient of \(x^n\) in \(p_n(x)\) is
\(\beta_1^n\), so \(\deg p_n=n\).  Conversely, exact degree for \(p_1\)
already forces \(\beta_1\ne0\).
\end{proof}

\section{Sheffer sequences}

\begin{definition}
A polynomial sequence \((s_n(x))_{n\ge0}\) is Sheffer for the pair
\((A,\beta)\) if
\begin{equation}\label{eq:ext-sheffer-egf}
 \sum_{n\ge0}s_n(x)\frac{z^n}{n!}
 =A(z)\EulerE^{x\beta(z)},
 \qquad A(0)\ne0,\quad \beta(0)=0,\quad \beta'(0)\ne0.
\end{equation}
The associated sequence \(p_n(x)\) is defined by
\(\sum p_n(x)z^n/n!=\EulerE^{x\beta(z)}\).
\end{definition}

\begin{theorem}[Sheffer identity and convolution]
Write \(A(z)=\sum_{j\ge0}a_jz^j/j!\).  Then
\begin{align}
 s_n(x+y)&=\sum_{j=0}^n\binom nj s_j(x)p_{n-j}(y),
 \label{eq:ext-sheffer-identity}\\
 s_n(x)&=\sum_{j=0}^n\binom nj a_jp_{n-j}(x).
 \label{eq:ext-sheffer-convolution}
\end{align}
Conversely, either identity together with the EGF of \(p_n\) reconstructs
\eqref{eq:ext-sheffer-egf}.
\end{theorem}
\begin{proof}
For the first identity,
\[
 A(z)\EulerE^{(x+y)\beta(z)}
 =\bigl(A(z)\EulerE^{x\beta(z)}\bigr)\EulerE^{y\beta(z)}.
\]
Apply the EGF product rule.  The second identity follows from
\(A(z)\EulerE^{x\beta(z)}=A(z)\sum p_n(x)z^n/n!\).  Conversely, summing either
convolution against \(z^n/n!\) recovers the corresponding product of EGFs.
\end{proof}

\begin{theorem}[Delta operator and lowering law]
Let \(\bar\beta\) be the compositional inverse of \(\beta\), and define the
shift-invariant differential operator
\begin{equation}\label{eq:ext-delta-operator}
 Q=\bar\beta(D_x).
\end{equation}
Then the associated sequence satisfies
\begin{equation}\label{eq:ext-delta-lowering}
 Qp_n(x)=np_{n-1}(x),
\end{equation}
and a Sheffer sequence for \((A,\beta)\) satisfies the same lowering law
\(Qs_n=ns_{n-1}\).
\end{theorem}
\begin{proof}
For any formal series \(q\),
\(q(D_x)\EulerE^{x\beta(z)}=q(\beta(z))\EulerE^{x\beta(z)}\).  With
\(q=\bar\beta\), this becomes
\[
 Qe^{x\beta(z)}=z \EulerE^{x\beta(z)}.
\]
Comparing the coefficient of \(z^n/n!\) gives
\eqref{eq:ext-delta-lowering}.  Since \(A(z)\) is independent of \(x\), the
same calculation applies to \(A(z)\EulerE^{x\beta(z)}\).
\end{proof}

\begin{example}[Appell families]
When \(\beta(z)=z\), one has \(Q=D_x\), and the Sheffer sequence is an Appell
sequence.  Bernoulli polynomials correspond to
\(A(z)=z/(\EulerE^z-1)\); Euler polynomials correspond to
\(A(z)=2/(\EulerE^z+1)\).  Their derivative laws
\eqref{eq:ext-bernoulli-derivative} and \eqref{eq:ext-euler-derivative} are
therefore instances of the general lowering law.
\end{example}

\begin{example}[Bernoulli--Cauchy family]
Equation \eqref{eq:ext-cauchy-poly-egf} has the Sheffer form
\[
 A(t)\EulerE^{x\beta(t)},\qquad
 A(t)=\frac{t}{\log(1+t)},\quad \beta(t)=\log(1+t).
\]
Its delta operator is \(Q=\EulerE^{D_x}-1=\Delta\), and hence
\begin{equation}\label{eq:ext-cauchy-delta}
 \mathsf C_n(x+1)-\mathsf C_n(x)=n\mathsf C_{n-1}(x).
\end{equation}
This identity follows directly by multiplying the generating function by
\((1+t)-1=t\).
\end{example}

\section{Abel polynomials from Lagrange inversion}
\index{Abel polynomials}

\begin{theorem}[Abel's binomial-type sequence]
Fix a scalar \(a\) and put
\begin{equation}\label{eq:ext-abel-polynomial}
 p_0(x)=1,\qquad p_n(x)=x(x-an)^{n-1}\quad(n\ge1).
\end{equation}
Then \((p_n)\) is of binomial type and therefore obeys
\begin{equation}\label{eq:ext-abel-identity}
 (x+y)(x+y-an)^{n-1}
 =\sum_{k=0}^n\binom nk
 x(x-ak)^{k-1}y\bigl(y-a(n-k)\bigr)^{n-k-1},
\end{equation}
with the endpoint factors interpreted through \(p_0=1\).
\end{theorem}
\begin{proof}
Let \(T(z)\) be the unique series satisfying
\(T=z \EulerE^{-aT}\).  Lagrange--B\"urmann gives, for \(n\ge1\),
\begin{align*}
 [z^n]\EulerE^{xT(z)}
 &=\frac1n[u^{n-1}]x \EulerE^{xu}\EulerE^{-anu}\\
 &=\frac{x(x-an)^{n-1}}{n!}.
\end{align*}
Thus
\[
 \EulerE^{xT(z)}=\sum_{n\ge0}p_n(x)\frac{z^n}{n!}.
\]
This is of the form \(\EulerE^{x\beta(z)}\), so
\cref{thm:ext-binomial-type-characterization} proves the binomial law.
Writing that law explicitly gives \eqref{eq:ext-abel-identity}.
\end{proof}

\section{Exponential Riordan arrays}

\begin{definition}
Let \(g(0)\ne0\) and let \(f(z)=f_1z+O(z^2)\) with \(f_1\ne0\).  The
exponential Riordan array \((g,f)\) is the lower-triangular matrix
\begin{equation}\label{eq:ext-riordan-entry}
 R_{n,k}=\frac{n!}{k!}[z^n]g(z)f(z)^k.
\end{equation}
\end{definition}

\begin{theorem}[Action, group law, and inverse]
If \(a=(a_k)\) has EGF \(A(z)=\sum a_kz^k/k!\), then \(b=Ra\) has EGF
\begin{equation}\label{eq:ext-riordan-action}
 B(z)=g(z)A(f(z)).
\end{equation}
Consequently,
\begin{align}
 (g,f)(h,\ell)&=\bigl(g\,(h\circ f),\ \ell\circ f\bigr),
 \label{eq:ext-riordan-product}\\
 (g,f)^{-1}&=\left(\frac1{g\circ\bar f},\bar f\right),
 \label{eq:ext-riordan-inverse}
\end{align}
where \(\bar f=f^{\langle-1\rangle}\).
\end{theorem}
\begin{proof}
By \eqref{eq:ext-riordan-entry},
\begin{align*}
 \sum_{n\ge0}\left(\sum_{k=0}^nR_{n,k}a_k\right)\frac{z^n}{n!}
 &=\sum_{k\ge0}\frac{a_k}{k!}g(z)f(z)^k\\
 &=g(z)A(f(z)),
\end{align*}
which proves the action.  Apply first \((h,\ell)\) and then \((g,f)\):
\[
 A(z)\longmapsto h(z)A(\ell(z))
 \longmapsto g(z)h(f(z))A(\ell(f(z))).
\]
This proves \eqref{eq:ext-riordan-product}.  Substituting
\(\bigl(1/(g\circ\bar f),\bar f\bigr)\) on either side of \((g,f)\) reduces
the second component to \(z\) and the first to one, proving
\eqref{eq:ext-riordan-inverse}.
\end{proof}

\begin{corollary}[Stirling and Bell arrays]
The second-kind and signed first-kind Stirling matrices are the inverse
exponential Riordan arrays
\begin{equation}\label{eq:ext-stirling-riordan}
 (1,\EulerE^z-1),\qquad (1,\log(1+z)).
\end{equation}
The Bell-polynomial matrix associated with
\(F(z)=\sum_{j\ge1}x_jz^j/j!\) is \((1,F)\), since
\begin{equation}\label{eq:ext-bell-riordan}
 \frac{n!}{k!}[z^n]F(z)^k
 =\ExponentialPartialBellPolynomial nk(x_1,\ldots,x_{n-k+1}).
\end{equation}
\end{corollary}
\begin{proof}
The first statement is precisely the column EGF formulas
\eqref{eq:signed-first-egf} and \eqref{eq:second-egf}, together with the fact
that \(\EulerE^z-1\) and \(\log(1+z)\) are compositional inverses.  The second is the
defining Bell-polynomial EGF \eqref{eq:partial-bell-egf}.
\end{proof}

\begin{remark}
Fa\`a di Bruno composition is multiplication in the exponential Riordan group,
while Lagrange inversion computes the compositional inverse needed by the group
inverse.  Thus three topics that are often presented separately are different
faces of one triangular-matrix calculus.
\end{remark}

\chapter{Higher products, multivariate composition, and Good inversion}
\label{ch:ext-multivariate}
\index{Fa\`a di Bruno formula!Fr\'echet form}
\index{Lagrange inversion!multivariate}

\section{Products, reciprocals, and quotients}

\begin{theorem}[Multinomial Leibniz rule]
If \(f_1,\ldots,f_q\) possess \(n\) derivatives, then
\begin{equation}\label{eq:ext-multinomial-leibniz}
 D^n\prod_{r=1}^q f_r(x)
 =\sum_{j_1+\cdots+j_q=n}
 \binom{n}{j_1,\ldots,j_q}
 \prod_{r=1}^q f_r^{(j_r)}(x).
\end{equation}
\end{theorem}
\begin{proof}
Induct on the number \(q\) of factors.  The case \(q=1\) is immediate.  If the
formula holds for \(q-1\) factors, apply the ordinary two-factor Leibniz rule
to \(F=\prod_{r=1}^{q-1}f_r\) and \(f_q\):
\[
 D^n(Ff_q)=\sum_{j_q=0}^n\binom n{j_q}
 D^{n-j_q}F\,f_q^{(j_q)}.
\]
Insert the induction hypothesis in \(D^{n-j_q}F\).  The coefficient of the
term indexed by \(j_1+\cdots+j_q=n\) is
\[
 \binom n{j_q}\binom{n-j_q}{j_1,\ldots,j_{q-1}}
 =\frac{n!}{j_1!\cdots j_q!},
\]
which is the multinomial coefficient in
\eqref{eq:ext-multinomial-leibniz}.  This proof requires only the stated
finite differentiability.
\end{proof}

\begin{theorem}[Higher reciprocal and quotient rules]
Let \(g(x)\ne0\).  For \(r\ge0\),
\begin{equation}\label{eq:ext-reciprocal-derivative}
 D^r\frac1{g(x)}
 =\sum_{k=0}^r(-1)^kk!\,
 \frac{\ExponentialPartialBellPolynomial rk(g'(x),g''(x),\ldots,g^{(r-k+1)}(x))}{g(x)^{k+1}},
\end{equation}
where \(\ExponentialPartialBellPolynomial00=1\).  Hence
\begin{equation}\label{eq:ext-higher-quotient}
 D^n\frac{f}{g}
 =\sum_{r=0}^n\binom nr f^{(n-r)}
 \sum_{k=0}^r(-1)^kk!\,
 \frac{\ExponentialPartialBellPolynomial rk(g',g'',\ldots)}{g^{k+1}}.
\end{equation}
\end{theorem}
\begin{proof}
Apply the Bell-polynomial form of Fa\`a di Bruno,
\eqref{eq:faa-bell}, to \(h(u)=u^{-1}\) and \(u=g(x)\).  Since
\(h^{(k)}(u)=(-1)^kk!u^{-k-1}\), coefficient substitution gives
\eqref{eq:ext-reciprocal-derivative}.  The ordinary two-factor Leibniz rule
applied to \(f\cdot g^{-1}\), followed by the reciprocal formula, gives
\eqref{eq:ext-higher-quotient}.
\end{proof}

\section{Mixed partial derivatives}

For a smooth scalar field \(y=g(x_1,\ldots,x_n)\), write
\(\partial_B y=\prod_{j\in B}\partial_{x_j}y\), with the derivatives applied
jointly rather than multiplied; explicitly,
\(\partial_B y=\partial^{|B|}y/\prod_{j\in B}\partial x_j\).

\begin{theorem}[Set-partition formula for mixed partials]
\label{thm:ext-faa-mixed}
For \(y=g(x_1,\ldots,x_n)\),
\begin{equation}\label{eq:ext-faa-mixed}
 \frac{\partial^n}{\partial x_1\cdots\partial x_n}f(y)
 =\sum_{\pi\in\SetPartitionOperator([n])}f^{(|\pi|)}(y)
   \prod_{B\in\pi}\partial_B y.
\end{equation}
\end{theorem}
\begin{proof}
Induct on \(n\).  The statement for \(n=1\) is the chain rule.  Suppose it
holds for \(n\) and differentiate with respect to \(x_{n+1}\).  In each term,
the derivative either strikes \(f^{(|\pi|)}(y)\), producing
\(f^{(|\pi|+1)}(y)\partial_{n+1}y\); combinatorially this adjoins the singleton
block \(\{n+1\}\).  Or it strikes one factor \(\partial_B y\), replacing it by
\(\partial_{B\cup\{n+1\}}y\); combinatorially this inserts \(n+1\) into the
chosen existing block \(B\).  Every partition of \([n+1]\) is obtained exactly
once by one of these alternatives, proving the induction step.
\end{proof}

For example, the five partitions of \([3]\) give
\begin{align}
 \partial_{123}f(y)={}&f'(y)\partial_{123}y
 +f''(y)\bigl((\partial_1y)(\partial_{23}y)
              +(\partial_2y)(\partial_{13}y)
              +(\partial_3y)(\partial_{12}y)\bigr)\notag\\
 &+f^{(3)}(y)(\partial_1y)(\partial_2y)(\partial_3y).
 \label{eq:ext-faa-mixed-three}
\end{align}

\section{The Fr\'echet Fa\`a di Bruno formula}

\begin{theorem}[Fr\'echet form]
\label{thm:ext-faa-frechet}
Let \(X,Y,Z\) be Banach spaces, let \(g:X\to Y\), \(f:Y\to Z\), and assume
that all Fr\'echet derivatives through order \(n\) exist and are continuous near
\(x\).  Then
\begin{equation}\label{eq:ext-faa-frechet}
 D^n(f\circ g)(x)[h_1,\ldots,h_n]
 =\sum_{\pi\in\SetPartitionOperator([n])}
 D^{|\pi|}f(g(x))
 \left[
   D^{|B|}g(x)[h_j:j\in B]
 \right]_{B\in\pi}.
\end{equation}
The order in which the blocks are supplied is immaterial because higher
Fr\'echet derivatives are symmetric multilinear maps.
\end{theorem}
\begin{proof}
The proof is the same partition induction as in
\cref{thm:ext-faa-mixed}, but we spell out the differential step.  Assume the
formula for \(n\) and differentiate its right side at \(x\) in direction
\(h_{n+1}\).  Differentiating the outer map
\(D^{|\pi|}f(g(x))\) gives
\[
 D^{|\pi|+1}f(g(x))
 \bigl[Dg(x)h_{n+1},\,D^{|B|}g(x)[h_j:j\in B]_{B\in\pi}\bigr],
\]
which corresponds to adjoining the singleton \(\{n+1\}\).  Differentiating
the argument attached to a block \(B\) gives
\[
 D^{|B|+1}g(x)[h_j:j\in B, h_{n+1}],
\]
which corresponds to inserting \(n+1\) into \(B\).  The product rule for a
multilinear map says that these are all terms.  As before, every partition of
\([n+1]\) has a unique predecessor obtained either by deleting its singleton
\(\{n+1\}\) or by deleting \(n+1\) from the unique larger block containing it.
Thus every term appears exactly once.
\end{proof}

\section{Compositions of several functions}

For three sufficiently differentiable scalar functions, two applications of
Fa\`a di Bruno give the explicit nested formula
\begin{align}
 D^n f(g(h(x)))
 &=\sum_{k=0}^n f^{(k)}(g(h(x)))
 \ExponentialPartialBellPolynomial nk\bigl((g\circ h)',(g\circ h)'',\ldots\bigr),
 \label{eq:ext-three-composition-outer}\\
 (g\circ h)^{(j)}
 &=\sum_{\ell=0}^j g^{(\ell)}(h(x))
 \ExponentialPartialBellPolynomial j\ell(h',h'',\ldots).
 \label{eq:ext-three-composition-inner}
\end{align}
Substitution of the second line into the first is a finite closed formula.  Its
combinatorial meaning is a hierarchy of partitions: inner blocks record
collections of derivatives acting on the same occurrence of \(h\), and outer
blocks collect those inner blocks according to occurrences of \(g\).  Flattening
this hierarchy is the partition-theoretic shadow of associativity of formal
series composition.

\section{Good's multivariate Lagrange inversion formula}
\index{Good's formula}

The determinant form below is due to Good; combinatorial and further formal
proofs were developed by Gessel and others \cite{Good1965,Gessel1987}.

Let \(\boldsymbol w=(w_1,\ldots,w_d)\) be the unique vector of formal series
satisfying
\begin{equation}\label{eq:ext-good-system}
 w_i=t_i\phi_i(\boldsymbol w),\qquad \phi_i(\boldsymbol0)\ne0,
 \quad 1\le i\le d.
\end{equation}
For multi-indices \(\boldsymbol n\in\NaturalNumbers^d\), write
\(\boldsymbol t^{\boldsymbol n}=\prod_i t_i^{n_i}\) and
\(\boldsymbol\phi^{\boldsymbol n}=\prod_i\phi_i^{n_i}\).

\begin{lemma}[Multivariate residue substitution]
\label{lem:ext-multivariate-residue-substitution}
Suppose \(\boldsymbol t=\boldsymbol T(\boldsymbol x)\) is a formal change of
variables of the form
\(T_i=x_i u_i(\boldsymbol x)\) with each \(u_i(\boldsymbol0)\) invertible.
Then
\begin{equation}\label{eq:ext-multivariate-residue-substitution}
 \ResidueOperator_{\boldsymbol t} A(\boldsymbol t)\,
   \Differential t_1\cdots\Differential t_d
 =\ResidueOperator_{\boldsymbol x} A(\boldsymbol T(\boldsymbol x))
   \det\!\left(\frac{\partial T_i}{\partial x_j}\right)
   \Differential x_1\cdots\Differential x_d.
\end{equation}
\end{lemma}
\begin{proof}
By linearity it is enough to treat a Laurent monomial
\(A(\boldsymbol t)=\prod_i t_i^{m_i}\).  Regard the integrand as the top
differential form
\[
 \Omega=\prod_i t_i^{m_i}\,\Differential t_1\wedge\cdots\wedge\Differential t_d.
\]
If \(m_i\ne-1\) for some \(i\), then, over characteristic zero,
\[
 \Omega=\Differential\left(
 (-1)^{i-1}\frac{t_i^{m_i+1}}{m_i+1}
 \prod_{j\ne i}t_j^{m_j}
 \bigwedge_{\substack{1\le r\le d\\ r\ne i}}\Differential t_r
 \right).
\]
Pullback commutes with the exterior derivative.  The multivariate residue of
an exact top form is zero: in coordinates it is a sum of derivatives
\(\partial Q_i/\partial x_i\), and no such derivative has an
\(x_i^{-1}\) coefficient.  Thus both sides of
\eqref{eq:ext-multivariate-residue-substitution} vanish whenever some
\(m_i\ne-1\).

It remains to take every \(m_i=-1\).  Since \(T_i=x_i u_i\),
\[
 \bigwedge_{i=1}^d\frac{\Differential T_i}{T_i}
 =\det\left(
 \delta_{ij}+\frac{x_j}{u_i}\frac{\partial u_i}{\partial x_j}
 \right)_{i,j=1}^d
 \bigwedge_{j=1}^d\frac{\Differential x_j}{x_j}.
\]
Every entry after the Kronecker delta has zero constant term, so the determinant
has constant term one and no other term can contribute to the coefficient of
\(x_1^{-1}\cdots x_d^{-1}\).  The transformed residue is therefore one,
which equals the residue of \(\prod_i t_i^{-1}\).  Linearity now proves the
rule for every Laurent series for which the indicated coefficient extractions
are locally finite.
\end{proof}

\begin{theorem}[Lagrange--Good inversion]
\label{thm:ext-good-inversion}
For every formal series \(F\),
\begin{equation}\label{eq:ext-good-inversion}
 [\boldsymbol t^{\boldsymbol n}]F(\boldsymbol w(\boldsymbol t))
 =[\boldsymbol x^{\boldsymbol n}]
 F(\boldsymbol x)\boldsymbol\phi(\boldsymbol x)^{\boldsymbol n}
 \det\!\left(
 \delta_{ij}-\frac{x_j}{\phi_i(\boldsymbol x)}
                  \frac{\partial\phi_i}{\partial x_j}(\boldsymbol x)
 \right)_{i,j=1}^d.
\end{equation}
\end{theorem}
\begin{proof}
Represent the desired coefficient by an iterated residue and make the formal
substitution
\begin{equation}\label{eq:ext-good-change}
 t_i=\frac{x_i}{\phi_i(\boldsymbol x)}.
\end{equation}
The inverse substitution is exactly \(x_i=w_i(\boldsymbol t)\).  The Jacobian
has entries
\[
 \frac{\partial t_i}{\partial x_j}
 =\frac1{\phi_i}\left(
   \delta_{ij}-\frac{x_i}{\phi_i}\frac{\partial\phi_i}{\partial x_j}
 \right).
\]
Therefore
\begin{align*}
 &[\boldsymbol t^{\boldsymbol n}]F(\boldsymbol w(\boldsymbol t))\\
 &\quad=\ResidueOperator_{\boldsymbol t}
 F(\boldsymbol w(\boldsymbol t))
 \prod_i t_i^{-n_i-1}\,\Differential t_1\cdots\Differential t_d\\
 &\quad=\ResidueOperator_{\boldsymbol x}F(\boldsymbol x)
 \prod_i x_i^{-n_i-1}\phi_i(\boldsymbol x)^{n_i}
 \det\!\left(
 \delta_{ij}-\frac{x_i}{\phi_i}
                 \frac{\partial\phi_i}{\partial x_j}
 \right)\Differential\boldsymbol x.
\end{align*}
It remains only to put the determinant in the displayed form.  Let
\(X=\DiagonalOperator(x_1,\ldots,x_d)\),
\(\Phi=\DiagonalOperator(\phi_1,\ldots,\phi_d)\), and let
\(J=(\partial\phi_i/\partial x_j)\).  The determinant just obtained is
\(\det(I-X\Phi^{-1}J)\).  Since \(X\) and \(\Phi\) commute and
\(\det(I-AB)=\det(I-BA)\),
\[
 \det(I-X\Phi^{-1}J)=\det(I-\Phi^{-1}JX),
\]
whose \((i,j)\)-entry is
\(\delta_{ij}-x_j\phi_i^{-1}\partial\phi_i/\partial x_j\).
The remaining residue is exactly the coefficient on the right of
\eqref{eq:ext-good-inversion}.
\end{proof}

\begin{corollary}[Reduction to one variable]
For \(d=1\), Good's formula reads
\begin{equation}\label{eq:ext-good-one-variable}
 [t^n]F(w(t))=[x^n]F(x)\phi(x)^n
 \left(1-\frac{x\phi'(x)}{\phi(x)}\right).
\end{equation}
Formal integration by parts transforms this into the usual
Lagrange--B\"urmann formula \eqref{eq:lagrange-burmann}.
\end{corollary}
\begin{proof}
The residue of the derivative of
\(F(x)\phi(x)^nx^{-n}\) is zero.  Expanding it gives
\[
 0=\ResidueOperator F'(x)\phi(x)^nx^{-n}\Differential x
   +n\ResidueOperator F(x)\phi(x)^{n-1}\phi'(x)x^{-n}\Differential x
   -n\ResidueOperator F(x)\phi(x)^nx^{-n-1}\Differential x.
\]
Consequently the \emph{sum} of the two residue terms represented by the
right-hand side of \eqref{eq:ext-good-one-variable} is
\[
 \ResidueOperator F(x)\phi(x)^nx^{-n-1}\Differential x
 -\ResidueOperator F(x)\phi(x)^{n-1}\phi'(x)x^{-n}\Differential x
 =\frac1n\ResidueOperator F'(x)\phi(x)^nx^{-n}\Differential x.
\]
The last member is precisely the residue form of
\eqref{eq:lagrange-burmann}.
\end{proof}

\chapter{Catalan refinements, polygon dissections, and associahedra}
\label{ch:ext-associahedra}
\index{Catalan numbers}
\index{Narayana numbers}
\index{associahedron}

The base chapters derive Catalan and Fuss--Catalan coefficients from Lagrange
inversion and explain the signed face expansion underlying compositional
inversion.  Here we add two refinements supplied by the supporting dossiers:
Narayana peak enumeration and a complete count of faces of every dimension of
the associahedron \cite{Stasheff1963,Loday2004}.

\section{Narayana numbers from a bivariate Lagrange equation}

Let \(C(z,u)\) count Dyck paths, where \(z\) marks semilength and \(u\) marks
peaks; the empty path contributes one.  Every nonempty Dyck path has the unique
first-return decomposition
\[
 U\,A\,D\,B.
\]
If \(A\) is empty, the initial \(UD\) is a peak; otherwise it is not.  Hence
\begin{equation}\label{eq:ext-narayana-functional}
 C(z,u)=1+zC(z,u)\bigl(C(z,u)+u-1\bigr).
\end{equation}

\begin{theorem}[Narayana formula]
For \(1\le k\le n\), the number of Dyck paths of semilength \(n\) with exactly
\(k\) peaks is
\begin{equation}\label{eq:ext-narayana-number}
 N(n,k)=\frac1n\binom nk\binom n{k-1}.
\end{equation}
Consequently,
\begin{equation}\label{eq:ext-narayana-sum-catalan}
 \sum_{k=1}^nN(n,k)=\frac1{n+1}\binom{2n}{n}.
\end{equation}
\end{theorem}
\begin{proof}
Put \(Y=C-1\).  Equation \eqref{eq:ext-narayana-functional} becomes
\[
 Y=z(1+Y)(u+Y).
\]
Lagrange inversion gives
\begin{align*}
 [z^nu^k]Y
 &=\frac1n[t^{n-1}u^k](1+t)^n(u+t)^n\\
 &=\frac1n\binom nk[t^{k-1}](1+t)^n
 =\frac1n\binom nk\binom n{k-1},
\end{align*}
proving \eqref{eq:ext-narayana-number}.  Summing over \(k\) is equivalent to
setting \(u=1\), where the functional equation reduces to
\(C=1+zC^2\).  Formula \eqref{eq:catalan-lagrange} then proves
\eqref{eq:ext-narayana-sum-catalan}.  Alternatively, Vandermonde convolution
sums the closed form directly.
\end{proof}

\begin{corollary}[Symmetry]
\begin{equation}\label{eq:ext-narayana-symmetry}
 N(n,k)=N(n,n+1-k).
\end{equation}
\end{corollary}
\begin{proof}
Apply \(\binom nk=\binom n{n-k}\) to both binomial factors in
\eqref{eq:ext-narayana-number}.
\end{proof}

\section{Polygon dissections and Schr\"oder trees}

Fix a convex polygon with \(N\ge3\) sides and distinguish one boundary edge as
the root.  A dissection using \(j\) noncrossing diagonals has \(j+1\) regions.
Its rooted planar dual is a plane rooted tree: each region is an internal node,
each non-root boundary edge is a leaf, and the cyclic order around regions
orders the children.  Every internal node has at least two children.  Conversely,
gluing polygons according to such a tree reconstructs the dissection.  Thus
dissections of an \(N\)-gon with \(j\) diagonals correspond to plane rooted
trees with
\begin{equation}\label{eq:ext-dissection-tree-parameters}
 L=N-1\quad\hbox{leaves},\qquad r=j+1\quad\hbox{internal nodes},
\end{equation}
all internal outdegrees at least two.

Let \(T(z,u)\) count these trees, with \(z\) marking leaves and \(u\) internal
nodes.  A tree is either a leaf or an internal node carrying an ordered sequence
of at least two subtrees, so
\begin{equation}\label{eq:ext-schroeder-tree-functional}
 T=z+u\frac{T^2}{1-T}.
\end{equation}

\begin{theorem}[Number of polygon dissections]
\label{thm:ext-polygon-dissections}
The number \(D(N,j)\) of dissections of a convex \(N\)-gon by exactly \(j\)
noncrossing diagonals is
\begin{equation}\label{eq:ext-polygon-dissection-number}
 D(N,j)=\frac1{N-1}\binom{N-3}{j}\binom{N+j-1}{j+1},
 \qquad 0\le j\le N-3.
\end{equation}
\end{theorem}
\begin{proof}
Solve \eqref{eq:ext-schroeder-tree-functional} for the Lagrange form:
\[
 T=z\frac{1-T}{1-(1+u)T}
   =z\left(1-\frac{uT}{1-T}\right)^{-1}.
\]
For \(L\ge1\), Lagrange inversion yields
\begin{align*}
 [z^Lu^r]T
 &=\frac1L[t^{L-1}u^r]
   \left(1-\frac{ut}{1-t}\right)^{-L}\\
 &=\frac1L\binom{L+r-1}{r}
   [t^{L-1-r}](1-t)^{-r}\\
 &=\frac1L\binom{L+r-1}{r}\binom{L-2}{r-1}.
\end{align*}
Substitute \(L=N-1\) and \(r=j+1\) from
\eqref{eq:ext-dissection-tree-parameters}; this is exactly
\eqref{eq:ext-polygon-dissection-number}.  The rooted edge causes no additional
factor because every polygon in the count is rooted in the same prescribed way.
\end{proof}

At the two extremes, \(D(N,0)=1\), while
\[
 D(N,N-3)=\frac1{N-1}\binom{2N-4}{N-2}
\]
is the Catalan number counting triangulations of the \(N\)-gon.

\section{The complete face vector of the associahedron}

We now use the common convention that \(\AssociahedronByDimension{n-2}\) is the \((n-2)\)-dimensional
associahedron whose vertices are parenthesizations of \(n\) letters, equivalently
triangulations of an \((n+1)\)-gon.  A set of \(j\) pairwise noncrossing
diagonals determines a face of codimension \(j\): its vertices are the
triangulations containing those diagonals.  Hence a \(d\)-dimensional face
corresponds to a dissection with
\(j=n-2-d\) diagonals.

\begin{theorem}[Face numbers]
\label{thm:ext-associahedron-face-numbers}
For \(0\le d\le n-2\), the number of \(d\)-faces of \(\AssociahedronByDimension{n-2}\) is
\begin{equation}\label{eq:ext-associahedron-face-number}
 f_d(\AssociahedronByDimension{n-2})=\frac1n\binom{n-2}{d}
                \binom{2n-d-2}{n-1}.
\end{equation}
In particular,
\begin{align}
 f_0(\AssociahedronByDimension{n-2})&=\frac1n\binom{2n-2}{n-1}=\CatalanNumber{n-1},
 \label{eq:ext-associahedron-vertices}\\
 f_{n-3}(\AssociahedronByDimension{n-2})&=\frac{(n+1)(n-2)}2,
 \label{eq:ext-associahedron-facets}\\
 f_{n-2}(\AssociahedronByDimension{n-2})&=1.
\end{align}
\end{theorem}
\begin{proof}
Apply \cref{thm:ext-polygon-dissections} with \(N=n+1\) and
\(j=n-2-d\):
\begin{align*}
 f_d(\AssociahedronByDimension{n-2})
 &=\frac1n\binom{n-2}{n-2-d}
   \binom{2n-d-2}{n-1-d}\\
 &=\frac1n\binom{n-2}{d}\binom{2n-d-2}{n-1},
\end{align*}
where the last equality uses symmetry of the second binomial coefficient.
This proves \eqref{eq:ext-associahedron-face-number}.  The three special cases
follow by direct simplification.  For facets, the result is also the number of
diagonals in an \((n+1)\)-gon.
\end{proof}

\begin{example}
For \(\AssociahedronByDimension{2}\), the formula gives \((f_0,f_1,f_2)=(5,5,1)\), the face vector of a
pentagon.  For \(\AssociahedronByDimension{3}\), it gives \((14,21,9,1)\): fourteen vertices,
twenty-one edges, nine polygonal facets, and the polytope itself.
\end{example}

\section{The Narayana \texorpdfstring{$h$}{h}-vector}

For a simple \(D\)-polytope \(P\), define its \(h\)-polynomial by
\begin{equation}\label{eq:ext-simple-h-polynomial}
 h_P(u)=\sum_{d=0}^D f_d(P)(u-1)^d.
\end{equation}

\begin{theorem}[Associahedral \(h\)-polynomial]
\label{thm:ext-associahedron-h}
For \(\AssociahedronByDimension{n-2}\),
\begin{equation}\label{eq:ext-associahedron-h}
 h_{\AssociahedronByDimension{n-2}}(u)=\sum_{r=0}^{n-2}N(n-1,r+1)u^r
 =\sum_{r=0}^{n-2}\frac1{n-1}
   \binom{n-1}{r}\binom{n-1}{r+1}u^r.
\end{equation}
\end{theorem}
\begin{proof}
The coefficient of \(u^r\) in
\eqref{eq:ext-simple-h-polynomial}, using
\eqref{eq:ext-associahedron-face-number}, is
\begin{align*}
 &\sum_{d=r}^{n-2}(-1)^{d-r}\binom dr
 \frac1n\binom{n-2}{d}\binom{2n-d-2}{n-1}\\
 &\quad=\frac1n\binom{n-2}{r}
 \sum_{j=0}^{n-2-r}(-1)^j\binom{n-2-r}{j}
 \binom{2n-r-j-2}{n-1}.
\end{align*}
The alternating Vandermonde identity
\begin{equation}\label{eq:ext-alternating-vandermonde}
 \sum_{j=0}^M(-1)^j\binom Mj\binom{R-j}{S}
 =\binom{R-M}{S-M}
\end{equation}
follows by extracting \([x^S]\) from
\((1+x)^R(1-(1+x)^{-1})^M=x^M(1+x)^{R-M}\).  Apply it with
\(M=n-2-r\), \(R=2n-r-2\), and \(S=n-1\).  The coefficient becomes
\[
 \frac1n\binom{n-2}{r}\binom n{r+1}
 =\frac1{n-1}\binom{n-1}{r}\binom{n-1}{r+1},
\]
which is \(N(n-1,r+1)\) by
\eqref{eq:ext-narayana-number}.
\end{proof}

\begin{corollary}[Dehn--Sommerville symmetry in this case]
The coefficients of \(h_{\AssociahedronByDimension{n-2}}(u)\) are palindromic.
\end{corollary}
\begin{proof}
This is the Narayana symmetry \eqref{eq:ext-narayana-symmetry} with
\(n\) replaced by \(n-1\).
\end{proof}

\section{Tamari adjacency and the geometric meaning of associativity}

Two vertices of \(\AssociahedronByDimension{n-2}\) are adjacent exactly when their triangulations differ by
a single diagonal flip, or equivalently when their parenthesizations differ by
one local use of \((ab)c\leftrightarrow a(bc)\).  Orienting each edge by the
forward rotation \((ab)c\to a(bc)\) produces the Hasse diagram of the Tamari
lattice.  The pentagonal \(\AssociahedronByDimension{2}\) is therefore not merely a picture of five
parenthesizations: its boundary records the coherence relation asserting that
all reassociation paths between four factors agree.  Higher-dimensional faces
encode simultaneous compatible reassociations, while
\eqref{eq:ext-associahedron-face-number} counts them exactly.

\chapter{Gamma, digamma, polygamma, and Barnes \texorpdfstring{$G$}{G}}
\label{ch:ext-gamma-barnes}
\index{gamma function}
\index{digamma function}
\index{polygamma function}
\index{Barnes G-function@Barnes \(G\)-function}

The first-kind Stirling chapters already derive convergent series for
\(\log\Gamma\), \(\psi\), and the Stieltjes constants.  We complement that
coefficient viewpoint with the classical products and with the next function in
the gamma hierarchy, Barnes's \(G\)-function \cite{NISTDLMF,WhittakerWatson1927}.

\section{Euler's integral and product for the gamma function}

For \(\Re z>0\), define
\begin{equation}\label{eq:ext-gamma-integral}
 \Gamma(z)=\int_0^\infty t^{z-1}\EulerE^{-t}\,\Differential t.
\end{equation}
Integration by parts gives
\begin{equation}\label{eq:ext-gamma-recurrence}
 \Gamma(z+1)=z\Gamma(z),\qquad \Gamma(1)=1.
\end{equation}
The boundary term vanishes because \(t^ze^{-t}\to0\) at infinity and
\(t^z\to0\) at zero when \(\Re z>0\).

\begin{theorem}[Euler limit and Weierstrass product]
For \(\Re z>0\),
\begin{equation}\label{eq:ext-euler-gamma-limit}
 \Gamma(z)=\lim_{N\to\infty}
 \frac{N!\,N^z}{z(z+1)\cdots(z+N)}.
\end{equation}
Consequently, for all \(z\in\ComplexNumbers\),
\begin{equation}\label{eq:ext-weierstrass-gamma}
 \frac1{\Gamma(z+1)}
 =\EulerE^{\gamma z}\prod_{k=1}^\infty
 \left(1+\frac zk\right)\EulerE^{-z/k},
\end{equation}
where \(\gamma=\lim_{N\to\infty}(H_N-\log N)\).  The product converges
locally uniformly.
\end{theorem}
\begin{proof}
The beta integral and \(N\) integrations by parts give
\[
 \int_0^1t^{z-1}(1-t)^N\Differential t
 =\frac{N!}{z(z+1)\cdots(z+N)}.
\]
Multiply by \(N^z\) and substitute \(u=Nt\):
\[
 \frac{N!N^z}{z(z+1)\cdots(z+N)}
 =\int_0^N u^{z-1}\left(1-\frac uN\right)^N\Differential u.
\]
For fixed \(u\), the factor in parentheses tends to \(\EulerE^{-u}\); moreover
\((1-u/N)^N\le \EulerE^{-u}\) on \(0\le u\le N\).  Dominated convergence on a
large fixed interval, followed by an exponentially small tail estimate, proves
\eqref{eq:ext-euler-gamma-limit}.

Apply the limit to \(\Gamma(z+1)\) and use the recurrence to write
\[
 \frac1{\Gamma(z+1)}
 =\lim_{N\to\infty}N^{-z}\prod_{k=1}^N\left(1+\frac zk\right).
\]
Since \(\log N=H_N-\gamma+o(1)\), this becomes
\[
 \lim_{N\to\infty}\EulerE^{\gamma z+o(1)}
 \prod_{k=1}^N\left(1+\frac zk\right)\EulerE^{-z/k},
\]
which is \eqref{eq:ext-weierstrass-gamma}.  For every compact set, the estimate
\(\log(1+z/k)-z/k=O(k^{-2})\) is uniform once the finitely many small values
of \(k\) are removed.  Hence the product converges locally uniformly and
defines an entire function.  It agrees with \(1/\Gamma(z+1)\) on the half-plane
where Euler's limit was derived, and therefore everywhere by analytic
continuation.
\end{proof}

\section{Digamma and polygamma series}

Define
\begin{equation}\label{eq:ext-digamma-def}
 \psi(z)=\frac{\Gamma'(z)}{\Gamma(z)},\qquad
 \psi^{(m)}(z)=\frac{\Differential^m}{\Differential z^m}\psi(z).
\end{equation}

\begin{theorem}[Series and Hurwitz-zeta representation]
For \(z\notin\{-1,-2,\ldots\}\),
\begin{equation}\label{eq:ext-digamma-series}
 \psi(z+1)=-\gamma+\sum_{k=1}^\infty\frac{z}{k(k+z)}.
\end{equation}
For \(m\ge1\) and \(z\notin\{0,-1,-2,\ldots\}\),
\begin{equation}\label{eq:ext-polygamma-hurwitz}
 \psi^{(m)}(z)=(-1)^{m+1}m!
 \sum_{k=0}^\infty\frac1{(z+k)^{m+1}}
 =(-1)^{m+1}m!\,\HurwitzZeta{m+1}{z}.
\end{equation}
Also,
\begin{equation}\label{eq:ext-digamma-recurrence}
 \psi(z+1)=\psi(z)+\frac1z,
 \qquad
 \psi^{(m)}(z+1)=\psi^{(m)}(z)+\frac{(-1)^m m!}{z^{m+1}}.
\end{equation}
\end{theorem}
\begin{proof}
Take logarithms in \eqref{eq:ext-weierstrass-gamma} and differentiate.  Local
uniform convergence permits termwise differentiation and gives
\[
 -\psi(z+1)=\gamma+\sum_{k\ge1}
 \left(\frac1{k+z}-\frac1k\right),
\]
which is \eqref{eq:ext-digamma-series}.  Replacing \(z\) by \(z-1\) and
differentiating \(m\) times gives the absolutely and locally uniformly
convergent series \eqref{eq:ext-polygamma-hurwitz}.  Finally, logarithmic
differentiation of \(\Gamma(z+1)=z\Gamma(z)\) proves the first recurrence;
further differentiation proves the second.
\end{proof}

\begin{corollary}[Pochhammer logarithmic derivatives]
For \(n\ge1\),
\begin{align}
 \frac{\Differential}{\Differential x}\log(x)_n
 &=\psi(x+n)-\psi(x)=\sum_{j=0}^{n-1}\frac1{x+j},
 \label{eq:ext-pochhammer-log-derivative}\\
 \frac{\Differential^m}{\Differential x^m}\log(x)_n
 &=\psi^{(m-1)}(x+n)-\psi^{(m-1)}(x)
 =(-1)^{m-1}(m-1)!\sum_{j=0}^{n-1}\frac1{(x+j)^m}
 \label{eq:ext-pochhammer-higher-log-derivative}
\end{align}
for \(m\ge1\).
\end{corollary}
\begin{proof}
Differentiate the gamma quotient \eqref{eq:ext-pochhammer-gamma} and use the
recurrences in \eqref{eq:ext-digamma-recurrence}; telescoping gives the finite
sums.
\end{proof}

\section{Barnes's canonical product}

\begin{definition}
Barnes's \(G\)-function is defined by the locally uniformly convergent product
\begin{equation}\label{eq:ext-barnes-product}
 \BarnesGFunction(z+1)=(2\pi)^{z/2}
 \exp\!\left[-\frac{z+(1+\gamma)z^2}{2}\right]
 \prod_{k=1}^\infty
 \left(1+\frac zk\right)^k
 \exp\!\left(\frac{z^2}{2k}-z\right).
\end{equation}
\end{definition}

\begin{proposition}[Convergence and functional equation]
The product \eqref{eq:ext-barnes-product} converges locally uniformly and defines
an entire function satisfying
\begin{equation}\label{eq:ext-barnes-recurrence}
 \BarnesGFunction(z+1)=\Gamma(z)\BarnesGFunction(z),\qquad \BarnesGFunction(1)=1.
\end{equation}
Consequently, for integers \(n\ge2\),
\begin{equation}\label{eq:ext-barnes-integers}
 \BarnesGFunction(n)=1!\,2!\cdots(n-2)!
 =\prod_{k=1}^{n-2}k^{n-1-k}.
\end{equation}
\end{proposition}
\begin{proof}
For \(z\) in a compact set,
\[
 k\log(1+z/k)+\frac{z^2}{2k}-z=O(k^{-2})
\]
uniformly, so the logarithm of the product converges locally uniformly.
To verify the recurrence, let \(G_N\) denote the product truncated at \(k=N\).
The weighted factors telescope:
\[
 \prod_{k=1}^N
 \left(\frac{1+(z+1)/k}{1+z/k}\right)^k
 =\frac{(N+z+1)^N}{\prod_{k=1}^N(z+k)}.
\]
A direct division of the remaining exponential and \((2\pi)^{z/2}\) factors
gives
\begin{align*}
 \frac{G_N(z+2)}{G_N(z+1)}
 ={}&\sqrt{2\pi}\,
 \frac{(N+z+1)^N}{\prod_{k=1}^N(z+k)}\\
 &\times\exp\!\left(
 -(1+\gamma)z-\frac{2+\gamma}{2}
 +\frac{2z+1}{2}H_N-N\right).
\end{align*}
Use Stirling's formula
\(N!\sim\sqrt{2\pi N}(N/e)^N\),
\(H_N=\log N+\gamma+O(N^{-1})\), and
\((1+(z+1)/N)^N\to \EulerE^{z+1}\).  The displayed expression differs by a factor
tending to one from Euler's limit
\[
 \Gamma(z+1)=\lim_{N\to\infty}
 \frac{N^zN!}{\prod_{k=1}^N(z+k)}.
\]
Therefore \(\BarnesGFunction(z+2)/\BarnesGFunction(z+1)=\Gamma(z+1)\), equivalent to
\eqref{eq:ext-barnes-recurrence}.  Setting \(z=0\) in the product gives
\(\BarnesGFunction(1)=1\).  Iterating the recurrence at positive integers proves the first
product in \eqref{eq:ext-barnes-integers}; collecting equal factors proves the
second.
\end{proof}

\begin{proposition}[Alexeiewsky's integral]
\label{prop:ext-alexeiewsky}
For real \(x\ge0\),
\begin{equation}\label{eq:ext-alexeiewsky}
 \int_0^x\log\Gamma(1+t)\Differential t
 =\frac{x}{2}\log(2\pi)-\frac{x(x+1)}2
  +x\log\Gamma(x+1)-\log \BarnesGFunction(x+1).
\end{equation}
The identity extends holomorphically to any simply connected domain avoiding
the poles and zeros after compatible logarithm branches are chosen.
\end{proposition}
\begin{proof}
Termwise differentiation of the locally uniformly convergent Barnes product,
followed by the digamma series \eqref{eq:ext-digamma-series}, gives
\begin{align*}
 \frac{\Differential}{\Differential x}\log \BarnesGFunction(x+1)
 &=\frac12\log(2\pi)-\frac12-(1+\gamma)x
   +\sum_{k\ge1}\frac{x^2}{k(k+x)}\\
 &=\frac12\log(2\pi)-\frac12-x+x\psi(x+1).
\end{align*}
Differentiate the right-hand side of \eqref{eq:ext-alexeiewsky}.  The terms
other than \(\log\Gamma(x+1)\) cancel by the displayed logarithmic derivative,
so its derivative is \(\log\Gamma(x+1)\), exactly the derivative of the left
side.  Both sides vanish at \(x=0\), proving the identity.
\end{proof}

\section{Euler--Maclaurin asymptotics for Barnes \texorpdfstring{$G$}{G}}

\begin{theorem}[Complete algebraic asymptotic expansion]
\label{thm:ext-barnes-asymptotic}
For every fixed \(p\ge1\), as \(x\to+\infty\),
\begin{align}
 \log \BarnesGFunction(x+1)={}&\frac{x^2}{2}\log x-\frac{3x^2}{4}
 +\frac{x}{2}\log(2\pi)-\frac1{12}\log x+\RiemannZeta'(-1)
 \notag\\
 &+\sum_{r=1}^{p-1}
 \frac{\beta_{2r+2}}{4r(r+1)x^{2r}}
 +O(x^{-2p}).
 \label{eq:ext-barnes-asymptotic}
\end{align}
The proof first determines every coefficient and the constant on positive
integers, then uses Alexeiewsky's integral to establish the same expansion for
all real \(x\to+\infty\).
\end{theorem}
\begin{proof}
We give the real integer derivation, which determines every displayed
coefficient and the constant.  From \eqref{eq:ext-barnes-integers},
\begin{equation}\label{eq:ext-barnes-sum-reduction}
 \log \BarnesGFunction(n+1)=n\log\Gamma(n+1)-\sum_{k=1}^{n}k\log k.
\end{equation}
Apply Euler--Maclaurin to \(f(x)=x\log x\).  Since
\[
 f'(x)=\log x+1,
 \qquad
 f^{(m)}(x)=(-1)^m(m-2)!x^{1-m}\quad(m\ge2),
\]
one obtains, to arbitrary order,
\begin{align}
 \sum_{k=1}^{n}k\log k
 ={}&\left(\frac{n^2}{2}+\frac n2+\frac1{12}\right)\log n
 -\frac{n^2}{4}+\frac1{12}-\RiemannZeta'(-1)
 \notag\\
 &-\sum_{r=1}^{p-1}
 \frac{\beta_{2r+2}}{2r(2r+1)(2r+2)n^{2r}}
 +O(n^{-2p}).
 \label{eq:ext-hyperfactorial-asymptotic}
\end{align}
Here the constant is identified by applying Euler--Maclaurin to
\(\RiemannZeta(s)=\sum_{k\ge1}k^{-s}\), differentiating the continued formula at
\(s=-1\), and comparing finite parts.  Explicitly, the first three tail terms
at \(s=-1\) are
\[
 \frac{n^2}{2}\log n-\frac{n^2}{4}
 +\frac n2\log n+\frac1{12}(\log n+1),
\]
while the differentiated Bernoulli tails give exactly the inverse even powers
in \eqref{eq:ext-hyperfactorial-asymptotic}; moving
\(-\sum_{k\le n}k\log k\) to the other side leaves \(\RiemannZeta'(-1)\).

The ordinary Stirling expansion is
\begin{equation}\label{eq:ext-stirling-loggamma-use}
 \log\Gamma(n+1)=\left(n+\frac12\right)\log n-n
 +\frac12\log(2\pi)
 +\sum_{r=1}^{p}\frac{\beta_{2r}}{2r(2r-1)n^{2r-1}}
 +O(n^{-2p-1}).
\end{equation}
Substitute \eqref{eq:ext-hyperfactorial-asymptotic} and
\eqref{eq:ext-stirling-loggamma-use} into
\eqref{eq:ext-barnes-sum-reduction}.  The polynomial and logarithmic terms
collect to the first line of \eqref{eq:ext-barnes-asymptotic}.  At each inverse
even power, the contribution from \(n\log\Gamma(n+1)\) and that from the
hyperfactorial sum combine as
\[
 \frac{\beta_{2r+2}}{(2r+2)(2r+1)}
 +\frac{\beta_{2r+2}}{2r(2r+1)(2r+2)}
 =\frac{\beta_{2r+2}}{4r(r+1)},
\]
which proves every coefficient in the second line.  The Euler--Maclaurin
remainder estimates give \(O(n^{-2p})\).

It remains to justify the asserted expansion between the integers.  Solve
Alexeiewsky's identity \eqref{eq:ext-alexeiewsky} for \(\log \BarnesGFunction(x+1)\), split
the integral at \(1\), and insert the Stirling expansion for
\(\log\Gamma(1+t)\) uniformly for \(t\ge1\).  If that expansion is truncated
after \(\beta_{2p}\), its remainder is \(O(t^{-2p-1})\); multiplication by
\(x\) at the endpoint contributes \(O(x^{-2p})\), while integration of the
remainder contributes an \(x\)-independent constant plus \(O(x^{-2p})\).
The elementary terms give
\[
 \frac{x^2}{2}\log x-\frac{3x^2}{4}
 +\frac{x}{2}\log(2\pi)-\frac1{12}\log x,
\]
and, for \(r\ge1\), the term \(\beta_{2r+2}\) contributes
\(\beta_{2r+2}/(4r(r+1)x^{2r})\).  Thus
\eqref{eq:ext-barnes-asymptotic} holds for all real \(x\to+\infty\) with one
constant independent of \(x\).  The integer calculation above identifies that
constant as \(\RiemannZeta'(-1)\), completing the proof.
\end{proof}

\begin{remark}
The constant \(A=\exp(1/12-\RiemannZeta'(-1))\) is the Glaisher--Kinkelin constant.
Equation \eqref{eq:ext-hyperfactorial-asymptotic} is equivalently the asymptotic
law for the hyperfactorial \(\prod_{k=1}^n k^k\).  Thus the constant appearing
in Barnes \(G\) is forced by the same finite-part calculation that continues
the zeta function to \(-1\).
\end{remark}

\chapter{Euler--Maclaurin summation and Darboux methods}
\label{ch:ext-em-darboux}
\index{Euler--Maclaurin formula}
\index{Darboux formula}

Bernoulli polynomials have two complementary analytic roles.  Periodizing them
turns sums into integrals with controlled remainders; using an auxiliary
polynomial in repeated integration by parts turns endpoint derivatives into an
exact Taylor-type identity.  Both mechanisms are developed here from first
principles; the normalization agrees with the modern DLMF convention
\cite{NISTDLMF}.

\section{Periodic Bernoulli functions}

For \(m\ge1\), write
\begin{equation}\label{eq:ext-periodic-bernoulli}
 \widetilde{\mathsf B}_m(x)=\mathsf B_m(\{x\}),
 \qquad \FractionalPart{x}=x-\Floor{x}.
\end{equation}
For \(m\ge2\) this is continuous at the integers because
\(\mathsf B_m(1)=\mathsf B_m(0)\).  Away from integers,
\begin{equation}\label{eq:ext-periodic-bernoulli-derivative}
 \frac{\Differential}{\Differential x}\frac{\widetilde{\mathsf B}_m(x)}{m!}
 =\frac{\widetilde{\mathsf B}_{m-1}(x)}{(m-1)!}.
\end{equation}

\begin{proposition}[Fourier series]
For \(m\ge2\),
\begin{equation}\label{eq:ext-bernoulli-fourier-complex}
 \widetilde{\mathsf B}_m(x)
 =-\frac{m!}{(2\pi \ImaginaryUnit)^m}
 \sum_{k\in\IntegerNumbers\setminus\{0\}}\frac{\EulerE^{2\pi ikx}}{k^m},
\end{equation}
with pointwise convergence everywhere and absolute convergence for \(m\ge2\)
except that the \(m=2\) derivative series is only conditionally convergent.  In
particular,
\begin{equation}\label{eq:ext-bernoulli-fourier-even}
 \widetilde{\mathsf B}_{2p}(x)
 =(-1)^{p+1}\frac{2(2p)!}{(2\pi)^{2p}}
 \sum_{k=1}^\infty\frac{\cos(2\pi kx)}{k^{2p}},
\end{equation}
and therefore
\begin{equation}\label{eq:ext-periodic-bernoulli-bound}
 \left|\widetilde{\mathsf B}_{2p}(x)\right|
 \le\frac{2(2p)!\RiemannZeta(2p)}{(2\pi)^{2p}}.
\end{equation}
\end{proposition}
\begin{proof}
For nonzero \(k\), let
\(c_k^{(m)}=\int_0^1\mathsf B_m(x)\EulerE^{-2\pi ikx}\Differential x\).  Integration by
parts, using \(\mathsf B_m(1)=\mathsf B_m(0)\) for \(m\ne1\), gives
\[
 c_k^{(m)}=\frac{m}{2\pi ik}c_k^{(m-1)}.
\]
At the last step, direct integration of \(\mathsf B_1(x)=x-1/2\) gives
\(c_k^{(1)}=-1/(2\pi ik)\).  Hence
\(c_k^{(m)}=-m!/(2\pi ik)^m\).  The mean coefficient is zero because
\(\int_0^1\mathsf B_m(x)\Differential x=0\) for \(m\ge1\), by
\eqref{eq:ext-bernoulli-derivative}.  This proves
\eqref{eq:ext-bernoulli-fourier-complex}; pairing \(k\) with \(-k\) gives the
even form.  Taking absolute values and summing proves
\eqref{eq:ext-periodic-bernoulli-bound}.
\end{proof}

\section{Euler--Maclaurin with an explicit remainder}

\begin{theorem}[Euler--Maclaurin summation]
\label{thm:ext-euler-maclaurin}
Let \(a<b\) be integers and \(f\in C^{2p}([a,b])\).  Then
\begin{align}
 \sum_{k=a}^{b}f(k)
 ={}&\int_a^b f(x)\Differential x+\frac{f(a)+f(b)}2
 +\sum_{r=1}^{p}\frac{\beta_{2r}}{(2r)!}
   \bigl(f^{(2r-1)}(b)-f^{(2r-1)}(a)\bigr)
 \notag\\
 &-\frac1{(2p)!}\int_a^b
   \widetilde{\mathsf B}_{2p}(x)f^{(2p)}(x)\Differential x.
 \label{eq:ext-euler-maclaurin}
\end{align}
The remainder satisfies
\begin{equation}\label{eq:ext-euler-maclaurin-bound}
 |R_p|\le
 \frac{2\RiemannZeta(2p)}{(2\pi)^{2p}}
 \int_a^b|f^{(2p)}(x)|\Differential x.
\end{equation}
\end{theorem}
\begin{proof}
On each interval \((k,k+1)\),
\(\widetilde{\mathsf B}_1(x)=x-k-1/2\).  Integration by parts gives
\[
 \int_k^{k+1}\widetilde{\mathsf B}_1(x)f'(x)\Differential x
 =\frac{f(k)+f(k+1)}2-\int_k^{k+1}f(x)\Differential x.
\]
Summing from \(k=a\) to \(b-1\) yields the exact starting identity
\begin{equation}\label{eq:ext-em-start}
 \sum_{k=a}^{b}f(k)
 =\int_a^b f(x)\Differential x+\frac{f(a)+f(b)}2
 +\int_a^b\widetilde{\mathsf B}_1(x)f'(x)\Differential x.
\end{equation}
Now use \eqref{eq:ext-periodic-bernoulli-derivative} repeatedly.  The first
integration gives
\[
 \int_a^b\widetilde{\mathsf B}_1f'
 =\frac{\beta_2}{2!}\bigl(f'(b)-f'(a)\bigr)
 -\frac1{2!}\int_a^b\widetilde{\mathsf B}_2f''.
\]
The next integration introduces \(\widetilde{\mathsf B}_3\), but its boundary
term vanishes because \(\beta_3=0\); the following integration produces the
\(\beta_4\) endpoint term.  Continuing in pairs gives exactly
\eqref{eq:ext-euler-maclaurin}.  The odd Bernoulli numbers beyond
\(\beta_1\) vanish by reflection, so no omitted endpoint terms remain.
Finally apply \eqref{eq:ext-periodic-bernoulli-bound} inside the remainder
integral to obtain \eqref{eq:ext-euler-maclaurin-bound}.
\end{proof}

\begin{corollary}[Infinite-tail form]
Suppose \(f^{(j)}(x)\to0\) as \(x\to\infty\) for
\(0\le j\le2p-1\), and the remainder integral converges.  Then
\begin{align}
 \sum_{k=a}^{\infty}f(k)
 ={}&\int_a^\infty f(x)\Differential x+\frac{f(a)}2
 -\sum_{r=1}^{p}\frac{\beta_{2r}}{(2r)!}f^{(2r-1)}(a)
 \notag\\
 &-\frac1{(2p)!}\int_a^\infty
 \widetilde{\mathsf B}_{2p}(x)f^{(2p)}(x)\Differential x.
 \label{eq:ext-euler-maclaurin-tail}
\end{align}
\end{corollary}
\begin{proof}
Apply \cref{thm:ext-euler-maclaurin} on \([a,b]\) and let the integer
\(b\to\infty\).  The endpoint hypotheses remove all terms at \(b\); the
assumed convergence of the improper remainder integral gives the limit of its
truncations.
\end{proof}

\begin{example}[Harmonic numbers]
Apply the finite formula \eqref{eq:ext-euler-maclaurin} to \(f(x)=1/x\) on
\([1,n]\), and absorb all lower-end terms and the limiting remainder into the
constant \(\gamma=\lim_{n\to\infty}(H_n-\log n)\).  Since
\(f^{(m)}(x)=(-1)^m m!x^{-m-1}\), one obtains, for fixed \(p\),
\begin{equation}\label{eq:ext-harmonic-asymptotic}
 H_n=\log n+\gamma+\frac1{2n}
 -\sum_{r=1}^{p-1}\frac{\beta_{2r}}{2r\,n^{2r}}
 +O(n^{-2p}).
\end{equation}
Indeed, the tail of the remainder integral is
\(O(\int_n^\infty x^{-2p-1}\Differential x)=O(n^{-2p})\) by
\eqref{eq:ext-periodic-bernoulli-bound}, which gives the displayed error
order.
\end{example}

\section{Darboux's exact integration-by-parts identity}

\begin{theorem}[Darboux formula]
\label{thm:ext-darboux-formula}
Let \(f\) be analytic on a neighborhood of the line segment joining \(a\) to
\(z\), and let \(\phi\) be a polynomial of degree \(n\).  Then
\begin{align}
 \phi^{(n)}(0)\bigl(f(z)-f(a)\bigr)
 ={}&\sum_{m=1}^n(-1)^{m-1}(z-a)^m
 \Bigl(\phi^{(n-m)}(1)f^{(m)}(z)
       -\phi^{(n-m)}(0)f^{(m)}(a)\Bigr)
 \notag\\
 &+(-1)^n(z-a)^{n+1}
 \int_0^1\phi(t)f^{(n+1)}(a+t(z-a))\Differential t.
 \label{eq:ext-darboux-formula}
\end{align}
\end{theorem}
\begin{proof}
Put \(d=z-a\) and define
\[
 S(t)=\sum_{m=1}^n(-1)^m d^m\phi^{(n-m)}(t)
 f^{(m)}(a+td).
\]
Differentiation produces two sums.  After shifting the index in the terms where
the derivative hits \(f^{(m)}\), all intermediate terms cancel, leaving
\begin{equation}\label{eq:ext-darboux-telescope}
 S'(t)=-d\phi^{(n)}(t)f'(a+td)
 +(-1)^nd^{n+1}\phi(t)f^{(n+1)}(a+td).
\end{equation}
Because \(\phi^{(n)}\) is constant, integration from zero to one turns the
first term into
\(-\phi^{(n)}(0)(f(z)-f(a))\).  Move it to the other side and substitute the
endpoint values of \(S\); the result is exactly
\eqref{eq:ext-darboux-formula}.
\end{proof}

\begin{corollary}[Taylor formula with integral remainder]
Taking \(\phi(t)=(t-1)^n\) gives
\begin{equation}\label{eq:ext-taylor-integral-remainder}
 f(z)=\sum_{m=0}^n\frac{f^{(m)}(a)}{m!}(z-a)^m
 +\frac{(z-a)^{n+1}}{n!}\int_0^1(1-t)^n
 f^{(n+1)}(a+t(z-a))\Differential t.
\end{equation}
\end{corollary}
\begin{proof}
Here \(\phi^{(n)}=n!\), all derivatives of order below \(n\) vanish at
\(t=1\), and
\(\phi^{(n-m)}(0)=n!(-1)^m/m!\).  Substitute these values into
\eqref{eq:ext-darboux-formula} and divide by \(n!\).
\end{proof}

\section{Darboux transfer for algebraic singularities}

\begin{theorem}[Elementary Darboux transfer]
\label{thm:ext-darboux-transfer}
Let \(0<|\rho|<R\), let
\(\alpha_j\notin\{0,-1,-2,\ldots\}\), and suppose
\begin{equation}\label{eq:ext-darboux-decomposition}
 F(z)=\sum_{j=1}^Jc_j(1-z/\rho)^{-\alpha_j}+H(z),
\end{equation}
where \(H\) is analytic in \(|z|<R\).  Then, for every
\(|\rho|<R_0<R\),
\begin{equation}\label{eq:ext-darboux-transfer}
 [z^n]F(z)=\rho^{-n}\sum_{j=1}^Jc_j
 \frac{\Gamma(n+\alpha_j)}{\Gamma(\alpha_j)\Gamma(n+1)}
 +O(R_0^{-n}).
\end{equation}
Moreover, for fixed \(\alpha\),
\begin{align}
 \frac{\Gamma(n+\alpha)}{\Gamma(n+1)}
 =n^{\alpha-1}\biggl(&1+\frac{\alpha(\alpha-1)}{2n}
 +\frac{\alpha(\alpha-1)(\alpha-2)(3\alpha-1)}{24n^2}
 +O(n^{-3})\biggr).
 \label{eq:ext-gamma-ratio-asymptotic}
\end{align}
\end{theorem}
\begin{proof}
The generalized binomial theorem gives the exact coefficient
\[
 [z^n](1-z/\rho)^{-\alpha}
 =\rho^{-n}\frac{(\alpha)_n}{n!}
 =\rho^{-n}\frac{\Gamma(n+\alpha)}{\Gamma(\alpha)\Gamma(n+1)}.
\]
Cauchy's estimate on the circle \(|z|=R_0\) gives
\([z^n]H=O(R_0^{-n})\), proving
\eqref{eq:ext-darboux-transfer}.  Finally insert Stirling's expansion for the
two gamma factors, subtract their logarithms, expand through \(n^{-2}\), and
exponentiate.  The logarithmic coefficients are
\(\alpha(\alpha-1)/2\) at order \(n^{-1}\) and
\(-\alpha(\alpha-1)(2\alpha-1)/12\) at order \(n^{-2}\); adding half the
square of the first coefficient after exponentiation gives the second
coefficient displayed in \eqref{eq:ext-gamma-ratio-asymptotic}.
\end{proof}

\begin{example}[Catalan asymptotics from the square-root singularity]
The Catalan generating function can be written
\begin{equation}\label{eq:ext-catalan-local}
 C(z)=\frac{1-\sqrt{1-4z}}{2z}
 =\frac2{1+\sqrt{1-4z}}
 =2-2X+2X^2-2X^3+\cdots,
 \qquad X=\sqrt{1-4z}.
\end{equation}
The integer powers of \(X^2=1-4z\) contribute only finitely many coefficients;
the successive half-integer powers predict the descending powers of \(n\) in
the asymptotic expansion.  For a completely explicit remainder proof, start
from the exact coefficient formula
\(\CatalanNumber{n}=\binom{2n}{n}/(n+1)\).  Stirling's expansion gives
\[
 \binom{2n}{n}=\frac{4^n}{\sqrt{\pi n}}
 \left(1-\frac1{8n}+\frac1{128n^2}+O(n^{-3})\right),
\]
while
\[
 \frac1{n+1}=\frac1n\left(1-\frac1n+\frac1{n^2}+O(n^{-3})\right).
\]
Multiplication yields
\begin{equation}\label{eq:ext-catalan-asymptotic}
 \CatalanNumber{n}=\frac{4^n}{\sqrt\pi\,n^{3/2}}
 \left(1-\frac9{8n}+O(n^{-2})\right)
\end{equation}
and, one order more precisely,
\begin{equation}\label{eq:ext-catalan-asymptotic-two}
 \CatalanNumber{n}=\frac{4^n}{\sqrt\pi\,n^{3/2}}
 \left(1-\frac9{8n}+\frac{145}{128n^2}+O(n^{-3})\right).
\end{equation}
Expanding the half-integer singular terms with
\eqref{eq:ext-gamma-ratio-asymptotic} gives the same coefficients; the exact
binomial calculation supplies the required control of the unshown local
terms.
\end{example}

\chapter{M\"obius inversion, explicit reversions, and coefficient algorithms}
\label{ch:ext-inversion-algorithms}
\index{M\"obius inversion!partition lattice}
\index{series reversion!algorithms}

This final extension chapter collects material that was distributed among the
five drafts: the incidence-algebra explanation of cumulants, inverse
trigonometric coefficient evaluations, Kepler's equation, and triangular or
Newton-style algorithms for all elementary formal-series operations.

\section{M\"obius inversion on the partition lattice}

Let \(\Pi_n\) denote the lattice of set partitions of \([n]\), ordered by
refinement, with bottom \(\PartitionLatticeMinimum\) and top \(\PartitionLatticeMaximum\).

\begin{theorem}[M\"obius function to the top]
\label{thm:ext-partition-mobius}
For \(n\ge1\) and \(\pi\in\Pi_n\),
\begin{equation}\label{eq:ext-partition-mobius}
 \mu_{\Pi_n}(\pi,\PartitionLatticeMaximum)
 =(-1)^{|\pi|-1}(|\pi|-1)!.
\end{equation}
More generally,
\begin{equation}\label{eq:ext-partition-mobius-general}
 \mu_{\Pi_n}(\PartitionLatticeMinimum,\pi)
 =\prod_{B\in\pi}(-1)^{|B|-1}(|B|-1)!.
\end{equation}
\end{theorem}
\begin{proof}
The interval \([\pi,\PartitionLatticeMaximum]\) is canonically isomorphic to
\(\Pi_{|\pi|}\): coarsening \(\pi\) is the same as partitioning its set of
blocks.  It therefore suffices to compute
\(u_m=\mu_{\Pi_m}(\PartitionLatticeMinimum,\PartitionLatticeMaximum)\).  The interval
\([\PartitionLatticeMinimum,\sigma]\) is the direct product
\(\prod_{B\in\sigma}\Pi_{|B|}\), and M\"obius functions multiply on direct
products; hence
\begin{equation}\label{eq:ext-mobius-block-product}
 \mu(\PartitionLatticeMinimum,\sigma)=\prod_{B\in\sigma}u_{|B|}.
\end{equation}
The defining M\"obius relation gives
\[
 \sum_{\sigma\in\Pi_m}\mu(\PartitionLatticeMinimum,\sigma)
 =\begin{cases}1,&m=0,1,\\0,&m\ge2.\end{cases}
\]
By the exponential formula, the EGF of the left side is
\[
 \exp\!\left(\sum_{m\ge1}u_m\frac{z^m}{m!}\right).
\]
It must therefore equal \(1+z\).  Taking the formal logarithm gives
\[
 \sum_{m\ge1}u_m\frac{z^m}{m!}=\log(1+z)
 =\sum_{m\ge1}(-1)^{m-1}\frac{z^m}{m},
\]
so \(u_m=(-1)^{m-1}(m-1)!\).  The interval isomorphism proves
\eqref{eq:ext-partition-mobius}, and substituting the values of \(u_{|B|}\)
into \eqref{eq:ext-mobius-block-product} proves
\eqref{eq:ext-partition-mobius-general}.
\end{proof}

\begin{corollary}[Moment--cumulant inversion]
If
\begin{equation}\label{eq:ext-moment-partition}
 m_n=\sum_{\pi\in\Pi_n}\prod_{B\in\pi}\kappa_{|B|},
\end{equation}
then
\begin{equation}\label{eq:ext-cumulant-partition}
 \kappa_n=\sum_{\pi\in\Pi_n}
 (-1)^{|\pi|-1}(|\pi|-1)!\prod_{B\in\pi}m_{|B|}.
\end{equation}
\end{corollary}
\begin{proof}
For each partition \(\sigma\), define
\(M(\sigma)=\prod_{B\in\sigma}m_{|B|}\) and
\(K(\sigma)=\prod_{B\in\sigma}\kappa_{|B|}\).  Applying
\eqref{eq:ext-moment-partition} independently inside every block of \(\sigma\)
gives
\[
 M(\sigma)=\sum_{\pi\le\sigma}K(\pi).
\]
M\"obius inversion on \(\Pi_n\) yields
\(K(\PartitionLatticeMaximum)=\sum_{\pi\le\PartitionLatticeMaximum}
 \mu(\pi,\PartitionLatticeMaximum)M(\pi)\).  Since
\(K(\PartitionLatticeMaximum)=\kappa_n\), substitution of
\eqref{eq:ext-partition-mobius} proves
\eqref{eq:ext-cumulant-partition}.
\end{proof}

\section{Inverse trigonometric coefficient identities}

\begin{theorem}[Inverse sine]
For \(|y|<1\),
\begin{equation}\label{eq:ext-arcsin-series}
 \arcsin y=\sum_{m\ge0}\frac1{2m+1}\binom{2m}{m}
 \frac{y^{2m+1}}{4^m}.
\end{equation}
Equivalently, Lagrange inversion gives the coefficient identity
\begin{equation}\label{eq:ext-sine-power-coefficient}
 [x^{2m}]\left(\frac{x}{\sin x}\right)^{2m+1}
 =\frac1{4^m}\binom{2m}{m}.
\end{equation}
\end{theorem}
\begin{proof}
The generalized binomial theorem gives
\[
 \frac{\Differential}{\Differential y}\arcsin y=(1-y^2)^{-1/2}
 =\sum_{m\ge0}\binom{2m}{m}\frac{y^{2m}}{4^m}.
\]
Integrate term by term from zero to \(y\) to obtain
\eqref{eq:ext-arcsin-series}.  On the other hand, the analytic inverse-function
form \eqref{eq:analytic-lagrange-coeff}, applied to \(f(x)=\sin x\), gives
\[
 [y^{2m+1}]\arcsin y
 =\frac1{2m+1}[x^{2m}]
 \left(\frac{x}{\sin x}\right)^{2m+1}.
\]
Comparing this with the coefficient in \eqref{eq:ext-arcsin-series} proves
\eqref{eq:ext-sine-power-coefficient}.
\end{proof}

\begin{theorem}[Inverse tangent]
For \(|y|<1\),
\begin{equation}\label{eq:ext-arctan-series}
 \arctan y=\sum_{m\ge0}\frac{(-1)^m}{2m+1}y^{2m+1},
\end{equation}
and
\begin{equation}\label{eq:ext-tangent-power-coefficient}
 [x^{2m}]\left(\frac{x}{\tan x}\right)^{2m+1}=(-1)^m.
\end{equation}
\end{theorem}
\begin{proof}
Integrate the geometric series
\((1+y^2)^{-1}=\sum_{m\ge0}(-1)^my^{2m}\) to obtain
\eqref{eq:ext-arctan-series}.  Lagrange inversion for \(f(x)=\tan x\) gives
\[
 [y^{2m+1}]\arctan y
 =\frac1{2m+1}[x^{2m}]
 \left(\frac{x}{\tan x}\right)^{2m+1}.
\]
Coefficient comparison proves \eqref{eq:ext-tangent-power-coefficient}.
\end{proof}

\section{Kepler's equation}
\index{Kepler equation}
\index{Bessel functions}

For eccentricity \(0\le e<1\), Kepler's equation is
\begin{equation}\label{eq:ext-kepler}
 E=M+e\sin E.
\end{equation}
Since \(\Differential(E-e\sin E)/\Differential E=1-e\cos E>0\), the map from eccentric anomaly
\(E\) to mean anomaly \(M\) is globally invertible on the real line, with a
\(2\pi\)-periodic correction \(E(M)-M\).

\begin{theorem}[Reversion and Fourier--Bessel forms]
Formally in \(e\),
\begin{equation}\label{eq:ext-kepler-reversion}
 E=M+\sum_{k\ge1}\frac{\EulerE^k}{k!}
 \left(\frac{\Differential}{\Differential M}\right)^{k-1}\sin^kM.
\end{equation}
For \(0\le e<1\), the convergent Fourier representation is
\begin{equation}\label{eq:ext-kepler-bessel}
 E=M+2\sum_{n\ge1}\frac{J_n(ne)}{n}\sin(nM).
\end{equation}
\end{theorem}
\begin{proof}
Equation \eqref{eq:ext-kepler} is the shifted implicit equation
\eqref{eq:shifted-reversion-equation} with \(x=M\), \(y=e\), and
\(f(v)=\sin v\).  Formula \eqref{eq:shifted-lagrange-v} gives
\eqref{eq:ext-kepler-reversion}.

For the Fourier form, write
\(E-M=\sum_{n\ge1}c_n\sin(nM)\); oddness removes cosine terms.  Changing
variables from \(M\) to \(E\) gives
\begin{align*}
 c_n
 &=\frac1\pi\int_0^{2\pi}(E-M)\sin(nM)\Differential M\\
 &=\frac e\pi\int_0^{2\pi}\sin E\,
   \sin\bigl(n(E-e\sin E)\bigr)(1-e\cos E)\Differential E.
\end{align*}
Since
\[
 \frac{\Differential}{\Differential E}\cos\bigl(n(E-e\sin E)\bigr)
 =-n(1-e\cos E)\sin\bigl(n(E-e\sin E)\bigr),
\]
integration by parts, with vanishing boundary term, gives
\[
 c_n=\frac e{\pi n}\int_0^{2\pi}
 \cos E\cos(nE-ne\sin E)\Differential E.
\]
Using
\[
 J_m(t)=\frac1{2\pi}\int_0^{2\pi}\cos(mE-t\sin E)\Differential E
\]
and the product-to-sum formula, the integral divided by \(\pi\) is
\(J_{n-1}(ne)+J_{n+1}(ne)\).  The Bessel recurrence
\[
 J_{n-1}(t)+J_{n+1}(t)=\frac{2n}{t}J_n(t)
\]
therefore gives \(c_n=2J_n(ne)/n\) when \(e>0\); the case \(e=0\) follows
directly, or by continuity.  This proves the stated value of every Fourier
coefficient.  For fixed
\(e<1\), the map \(E\mapsto E-e\sin E\) has nonzero derivative on the real
axis.  The analytic inverse-function theorem, applied uniformly on the compact
quotient \(E\bmod 2\pi\), extends its inverse to a strip
\(|\Im M|<\eta_e\).  Hence the periodic analytic function \(E(M)-M\) has
Fourier coefficients \(O(\EulerE^{-\eta_e n})\) by shifting the coefficient contour
to the boundary of any smaller strip.  Its Fourier series therefore converges
absolutely and uniformly, proving \eqref{eq:ext-kepler-bessel} rather than
merely identifying its formal coefficients.
\end{proof}

\section{Triangular coefficient recurrences}

Let \(A(z)=\sum_{n\ge0}a_nz^n\),
\(B(z)=\sum_{n\ge0}b_nz^n\), and
\(C(z)=\sum_{n\ge0}c_nz^n\).  The following recurrences determine the result
through degree \(N\) using only lower coefficients.

\begin{proposition}[Product, reciprocal, and quotient]
If \(C=AB\), then
\begin{equation}\label{eq:ext-alg-product}
 c_n=\sum_{j=0}^na_jb_{n-j}.
\end{equation}
If \(a_0\ne0\) and \(C=1/A\), then
\begin{equation}\label{eq:ext-alg-reciprocal}
 c_0=a_0^{-1},\qquad
 c_n=-a_0^{-1}\sum_{j=1}^na_jc_{n-j}.
\end{equation}
If \(C=B/A\), then
\begin{equation}\label{eq:ext-alg-quotient}
 c_n=a_0^{-1}\left(b_n-\sum_{j=1}^na_jc_{n-j}\right).
\end{equation}
\end{proposition}
\begin{proof}
Take the coefficient of \(z^n\) in \(C=AB\), \(AC=1\), and \(AC=B\),
respectively, and isolate the term containing \(c_n\).
\end{proof}

\begin{proposition}[Logarithm, exponential, and arbitrary powers]
Assume \(a_0=1\).  If
\(L=\log A=\sum_{n\ge1}\ell_nz^n\), then
\begin{equation}\label{eq:ext-alg-log}
 \ell_n=a_n-\frac1n\sum_{j=1}^{n-1}j\ell_j a_{n-j}.
\end{equation}
Conversely, if \(A=\EulerE^L\), then
\begin{equation}\label{eq:ext-alg-exp}
 a_0=1,\qquad a_n=\frac1n\sum_{j=1}^n j\ell_j a_{n-j}.
\end{equation}
For \(C=A^\alpha\),
\begin{equation}\label{eq:ext-alg-power}
 c_0=1,
 \qquad
 c_n=\frac1n\sum_{j=1}^n\bigl((\alpha+1)j-n\bigr)a_jc_{n-j}.
\end{equation}
\end{proposition}
\begin{proof}
The identity \(A'=AL'\) gives
\(na_n=\sum_{j=1}^n j\ell_j a_{n-j}\).  Isolating \(\ell_n\) proves
\eqref{eq:ext-alg-log}; isolating \(a_n\) proves
\eqref{eq:ext-alg-exp}.  For powers, logarithmic differentiation gives
\(AC'=\alpha A'C\).  The coefficient of \(z^{n-1}\) is
\[
 \sum_{j=0}^{n-1}(n-j)a_jc_{n-j}
 =\alpha\sum_{j=1}^n j a_jc_{n-j}.
\]
Isolating \(nc_n\) yields \eqref{eq:ext-alg-power}.
\end{proof}

\begin{proposition}[Composition and reversion]
Suppose \(B(0)=0\).  Then
\begin{equation}\label{eq:ext-alg-composition}
 [z^n]A(B(z))
 =\sum_{k=0}^na_k\OrdinaryPartialBellPolynomial nk(b_1,\ldots,b_{n-k+1}).
\end{equation}
If \(A(0)=0\), \(a_1\ne0\), and
\(B=A^{\langle-1\rangle}\), then
\begin{align}
 b_1&=a_1^{-1},\label{eq:ext-alg-reversion-first}\\
 b_n&=-a_1^{-1}\sum_{k=2}^na_k
 \OrdinaryPartialBellPolynomial nk(b_1,\ldots,b_{n-k+1}),\qquad n\ge2.
 \label{eq:ext-alg-reversion-triangular}
\end{align}
\end{proposition}
\begin{proof}
Equation \eqref{eq:ext-alg-composition} is the ordinary composition theorem
\eqref{eq:ordinary-composition-bell}.  Set \(A(B(z))=z\).  The coefficient of
\(z\) gives \(a_1b_1=1\).  For \(n\ge2\), the coefficient vanishes and its
\(k=1\) term is \(a_1b_n\); every term with \(k\ge2\) depends only on
\(b_1,\ldots,b_{n-1}\).  Isolating \(b_n\) proves
\eqref{eq:ext-alg-reversion-triangular}.
\end{proof}

\section{Newton iteration and precision doubling}

Triangular recurrences are transparent and excellent for symbolic work.  With
fast polynomial multiplication, Newton iteration is asymptotically faster at
very high truncation orders.

\begin{theorem}[Reciprocal Newton iteration]
Suppose \(A(0)\) is invertible and \(C_m\) satisfies
\(AC_m\equiv1\pmod{z^m}\).  Define
\begin{equation}\label{eq:ext-newton-reciprocal}
 C_{2m}=C_m(2-AC_m)\pmod{z^{2m}}.
\end{equation}
Then \(AC_{2m}\equiv1\pmod{z^{2m}}\).
\end{theorem}
\begin{proof}
Write \(AC_m=1-E\) with \(E=O(z^m)\).  Then
\[
 AC_m(2-AC_m)=(1-E)(1+E)=1-E^2,
\]
and \(E^2=O(z^{2m})\).
\end{proof}

\begin{theorem}[Newton reversion]
Let \(A(0)=0\), \(A'(0)\) invertible, and suppose
\(A(B_m(z))\equiv z\pmod{z^m}\).  Put
\begin{equation}\label{eq:ext-newton-reversion}
 B_{2m}=B_m-\frac{A(B_m)-z}{A'(B_m)}\pmod{z^{2m}}.
\end{equation}
Then \(A(B_{2m})\equiv z\pmod{z^{2m}}\).
\end{theorem}
\begin{proof}
Let \(E=A(B_m)-z=O(z^m)\) and
\(\Delta=-E/A'(B_m)=O(z^m)\).  Formal Taylor expansion gives
\[
 A(B_m+\Delta)=A(B_m)+A'(B_m)\Delta+O(\Delta^2)
 =z+O(z^{2m}).
\]
The reciprocal of \(A'(B_m)\) exists because its constant term is
\(A'(0)\).  This proves precision doubling.
\end{proof}

\begin{remark}[Complexity]
With schoolbook multiplication, filling all triangular recurrences through
order \(N\) costs between \(O(N^2)\) and \(O(N^3)\), depending on whether Bell
polynomials are reused.  Newton iteration reduces reciprocal, logarithm,
exponential, and reversion to a logarithmic number of truncated multiplications
and compositions.  The algebraic identities above are independent of any
particular software package or coefficient representation.
\end{remark}


\part{Archive Concordance and Reference Record}

\chapter{Concordance with the source archive}
\label{ch:concordance}

The supplied package contains five substantial LaTeX drafts in
\nolinkurl{articles/}, together with encyclopedic reference dossiers in
\nolinkurl{Wikipedia/} and \nolinkurl{MathWorld/}
\cite{WikipediaDossier,MathWorldDossier}.  The drafts are not five successive
chapters: they are parallel treatments with a large common core and different
strengths.  This synthesis therefore follows mathematical dependence rather
than file order.  Repeated definitions and proofs have been deduplicated;
distinct examples, algorithms, geometric interpretations, and special-function
extensions have been retained and reconciled.  In the tables below,
``normalized'' means algebraically equivalent after applying the notation and
boundary conventions fixed in the opening chapters; ``corrected'' means that an
index, hypothesis, Jacobian, convergence domain, or undefined placeholder has
been repaired and the repaired statement proved.

\section{The five draft manuscripts}
\renewcommand{\arraystretch}{1.15}
\begin{longtable}{>{\raggedright\arraybackslash}p{0.28\textwidth}
                  >{\raggedright\arraybackslash}p{0.29\textwidth}
                  >{\raggedright\arraybackslash}p{0.31\textwidth}}
\toprule
\textbf{Draft} & \textbf{Principal strengths retained} &
\textbf{Role in the unified article}\\
\midrule
\endfirsthead
\toprule
\textbf{Draft} & \textbf{Principal strengths retained} &
\textbf{Role in the unified article}\\
\midrule
\endhead
\midrule
\multicolumn{3}{r}{\small Continued on the next page}\\
\endfoot
\bottomrule
\endlastfoot

\nolinkurl{Combinatorial_Formulae_and_Inversion_Theorems.tex}
& Foundational proofs for factorial bases, Stirling inversion, Bell-polynomial
  composition, Fa\`a di Bruno, and Lagrange inversion; careful distinction
  between formal and analytic arguments.
& Supplies proof details and normalization checks throughout
  Chapters~\ref{ch:formal}--\ref{ch:ext-multivariate}.\\[1ex]

\nolinkurl{Formulae_Compendium.tex}
& Compact cross-family formula tables, explicit reversion examples, higher
  product and quotient rules, and useful coefficient recurrences.
& Its tables have been converted into derived theorems and algorithms rather
  than repeated as an independent catalogue; see especially
  Chapters~\ref{ch:ext-pochhammer}, \ref{ch:ext-multivariate}, and
  \ref{ch:ext-inversion-algorithms}.\\[1ex]

\nolinkurl{Formulae_Monograph_Package/Formulae_Monograph.tex}
& Broad special-function coverage, restricted partition variants, asymptotic
  methods, and extended worked examples.
& Provides much of the breadth in Chapters~\ref{ch:stirling-asymptotics},
  \ref{ch:stirling-variants}, \ref{ch:ext-gamma-barnes}, and
  \ref{ch:ext-em-darboux}.\\[1ex]

\nolinkurl{Formulae_Monograph_Source_and_PDF/Formulae_Monograph.tex}
& Geometric and polyhedral interpretations, computational observations, and a
  second independent audit of conventions and boundary cases.
& Its nonduplicate material is incorporated in the Eulerian, Catalan, and
  associahedral chapters, while duplicated exposition has been collapsed into
  the common proofs.\\[1ex]

\nolinkurl{combinatorial_coefficient_calculus.tex}
& The most coherent global organization: formal series, the Stirling--Bell--
  Eulerian chain, composition calculus, inversion, and an archive concordance.
& Serves as the structural backbone.  Every inherited result was nevertheless
  checked against the other four drafts, and the new factorial, umbral,
  multivariate, special-function, and algorithmic parts were integrated into
  its dependency order.\\
\end{longtable}
\renewcommand{\arraystretch}{1}

\section{The supporting topic dossiers}
The Wikipedia snapshots supply detailed topic inventories and historical or
combinatorial interpretations; the MathWorld dossier adds Bernoulli, Euler,
Genocchi, Cauchy, Pochhammer, gamma/Barnes, Darboux, Euler--Maclaurin,
multivariate inversion, and associahedral material.  These have not been copied
as isolated encyclopedia entries.  They have been recast in the common
coefficient language and supplied with proofs in
Chapters~\ref{ch:ext-pochhammer}--\ref{ch:ext-inversion-algorithms}.  The table
that follows records the absorption of the principal Wikipedia topic files.

\renewcommand{\arraystretch}{1.15}
\begin{longtable}{>{\raggedright\arraybackslash}p{0.18\textwidth}
                  >{\raggedright\arraybackslash}p{0.37\textwidth}
                  >{\raggedright\arraybackslash}p{0.33\textwidth}}
\toprule
\textbf{Source file} & \textbf{Topics represented} & \textbf{Location and status}\\
\midrule
\endfirsthead
\toprule
\textbf{Source file} & \textbf{Topics represented} & \textbf{Location and status}\\
\midrule
\endhead
\midrule
\multicolumn{3}{r}{\small Continued on the next page}\\
\endfoot
\bottomrule
\endlastfoot


\nolinkurl{Associahedron.txt / Associahedron.pdf}
& Stasheff polytopes, equivalent realizations, polygon dissections and
  triangulations, the Tamari order, face enumeration, and appearances in
  homotopy theory and moduli spaces.
& Chapters 22 and~\ref{ch:ext-associahedra}.  The polygon-dissection model is
  used to derive the complete face vector and the Narayana $h$-vector, with
  dimensions and index shifts stated explicitly.\\[1ex]

\nolinkurl{Catalan_number.txt / Catalan_number.pdf}
& The Catalan recurrence and closed form; binary trees, Dyck paths,
  parenthesizations, triangulations, and noncrossing structures; generating
  functions, convolution, and asymptotic growth.
& Chapters 22 and~\ref{ch:ext-associahedra}, together with the Darboux example
  in Chapter~\ref{ch:ext-em-darboux}.  Lagrange inversion proves the closed
  form, the dissection chapter organizes the bijections, and Stirling--Darboux
  analysis supplies a controlled full asymptotic expansion.\\[1ex]

\nolinkurl{Bell_number.txt}
& Set partitions and equivalence relations; squarefree factorizations; rhyme
  schemes and restricted-growth words; the card-shuffle and pattern-avoidance
  interpretations; the Bell--Aitken triangle; binomial and Stirling sums;
  Spivey-type identities; the exponential generating function and its
  differential equation; Dobiński's formula and Poisson moments; Touchard
  congruences and modular periods; contour representation; inequalities,
  log-convexity, saddle-point growth, and Bell-prime computation.
& Chapters 9--11.  The many uses of the letter $B$ are separated.  Finite lists
  of computed primes and periods are treated as reproducible computations, while
  symbolic claims are proved.  The source's unnamed higher saddle coefficients
  are replaced by a defined recursive saddle expansion.\\[1ex]

\nolinkurl{Bell_polynomials.txt}
& Exponential partial and complete Bell polynomials; ordinary Bell polynomials;
  weighted-partition interpretation and low-degree tables; generating functions;
  pointing and convolution recurrences; derivatives and scaling; Stirling, Lah,
  Touchard, and symmetric-function specializations; exponential--logarithmic
  inversion; Hessenberg determinants; moments and cumulants; cycle indices;
  Hermite polynomials; binomial-type sequences; Faà di Bruno, inverse-series, and
  Laplace applications.
& Chapters 12--15 and 19--24.  All conventions are reconciled through one
  coefficient formula.  The quadratic differential recurrence is proved by
  translating every derivative back to the master generating function.\\[1ex]

\nolinkurl{Eulerian_number.txt}
& Descents and ascents; recurrence, symmetry, explicit formula, and exponential
  generating function; Eulerian polynomials; Worpitzky's identity and power sums;
  negative polylogarithms; conversion to the Stirling basis; the Irwin--Hall
  integral; alternating sums and Bernoulli numbers; cube, hypersimplex, and
  permutohedron interpretations; type-B and second-order Eulerian numbers.
& Chapters 16--18.  The type-B descent statistic is indexed at positions
  $0,1,\ldots,n-1$ with $\pi_0=0$; this corrects the one-step shift in the archived
  formula.  The second-order interpolation formula is proved from its differential
  recurrence and generating function.\\[1ex]

\nolinkurl{Fa#U00e0_di_Bruno's_formula.txt}
& Set-partition, multiplicity, and Bell-polynomial forms of the higher chain rule;
  worked low-order derivatives; the exponential specialization; mixed partial
  derivatives and multivariate composition.
& Chapter 19.  All forms follow from the same labelled-partition count, and the
  multivariate formula is proved by distributing derivative labels among inner
  components.  The filename's encoded accent is rendered throughout as
  ``Faà di Bruno.''\\[1ex]

\nolinkurl{Lagrange_inversion_theorem.txt}
& Formal inversion by residues; the Lagrange--Bürmann coefficient formula;
  analytic local inversion; coefficient expressions through Bell polynomials;
  algebraic examples, Lambert $W$, Catalan and Fuss--Catalan series, and the
  rooted-tree/associahedral interpretation.
& Chapters 21--22.  Formal and analytic hypotheses are stated separately.
  Coefficient extraction, residues, and contour integrals are shown to be three
  realizations of the same identity.\\[1ex]

\nolinkurl{Lagrange_reversion_theorem.txt}
& The shifted two-parameter implicit equation; reversion coefficients and their
  derivatives; the delta-function derivation.
& Chapter 23.  The delta argument is decoded as a change-of-variables statement:
  its Jacobian and orientation/absolute-value issue are made explicit rather than
  hidden in symbolic delta notation.\\[1ex]

\nolinkurl{Stirling_number.txt}
& Unified first- and second-kind notation; falling and rising factorial bases;
  inverse triangular matrices; Lah numbers; Stirling transforms; finite
  differences; differential-operator normal ordering; symmetric-function
  descriptions; extensions to negative indices.
& Chapters 2--6.  Signed and unsigned first-kind numbers are never conflated, and
  the boundary convention makes every matrix identity valid without exceptional
  edge cases.\\[1ex]

\nolinkurl{Stirling_numbers_of_the_first_kind.txt}
& Permutation cycles and recurrences; tables and fixed-$k$ formulas; harmonic and
  near-diagonal identities; finite sums and convolutions; congruences; ordinary
  and exponential generating functions; explicit double sums; real-rootedness
  and asymptotics; zeta and Hurwitz-zeta sums; convergent expansions for
  $\log\Gamma$, the digamma function, and Stieltjes constants; generalized
  factorial coefficients.
& Chapters 2--8.  Near-diagonal partition formulas are derived by Lagrange
  inversion.  The special-function series are supplied with their actual
  convergence domains and justified by Binet/Hermite integral transforms.\\[1ex]

\nolinkurl{Stirling_numbers_of_the_second_kind.txt}
& Set partitions, surjections, and recurrences; Bell-number relations; explicit
  inclusion--exclusion formula; parity; rational and exponential generating
  functions; finite identities, bounds, saddle asymptotics, row maxima and
  log-concavity; Poisson and fixed-point moments; rhyme schemes; $r$-Stirling,
  associated, and reduced variants.
& Chapters 2--11.  The two backward recurrences presented tentatively in the
  source are proved exactly by coefficient extraction.  Strict log-concavity is
  obtained from real-rooted Touchard polynomials, and reduced numbers are proved
  through a chromatic-polynomial argument.\\
\end{longtable}
\renewcommand{\arraystretch}{1}

\subsection{The MathWorld dossier}
The following grouped inventory records every MathWorld topic file in the
archive.  A title followed by an accompanying \nolinkurl{.nb} notebook is
listed only once; the notebook is retained in the archive as provenance and as
an independently inspectable symbolic source, while all identities used here
are derived in the text.

\renewcommand{\arraystretch}{1.15}
\begin{longtable}{>{\raggedright\arraybackslash}p{0.22\textwidth}
                  >{\raggedright\arraybackslash}p{0.37\textwidth}
                  >{\raggedright\arraybackslash}p{0.30\textwidth}}
\toprule
\textbf{Topic group} & \textbf{MathWorld files represented} &
\textbf{Absorption into the monograph}\\
\midrule
\endfirsthead
\toprule
\textbf{Topic group} & \textbf{MathWorld files represented} &
\textbf{Absorption into the monograph}\\
\midrule
\endhead
\midrule
\multicolumn{3}{r}{\small Continued on the next page}\\
\endfoot
\bottomrule
\endlastfoot

Factorial and classical polynomial calculus
& \emph{Pochhammer Symbol}; \emph{Bernoulli Number}; \emph{Bernoulli
  Polynomial}; \emph{Bernoulli Polynomial of the Second Kind}; \emph{Euler
  Polynomial}; \emph{Genocchi Number}; and \emph{Modified Bernoulli Number}.
& Chapters~\ref{ch:ext-pochhammer} and~\ref{ch:ext-bernoulli}.  Their generating
  functions, difference equations, reflection laws, multiplication formulas,
  and conversions to factorial bases are proved in one umbral framework.\\[1ex]

Bell and Dobi\'nski material
& \emph{Bell Number}; \emph{Bell Polynomial}; \emph{Dobi\'nski's Formula}; and
  \emph{Complementary Bell Number}.
& Chapters 9--15 and~\ref{ch:ext-bernoulli}.  Set-partition derivations,
  exponential generating functions, moment representations, Bell-polynomial
  recurrences, and the complementary sequence are connected without conflating
  their several conventional uses of $B_n$.\\[1ex]

Stirling and Catalan families
& \emph{Stirling Number of the First Kind}; \emph{Stirling Number of the Second
  Kind}; and \emph{Catalan Number}.
& Chapters 2--8, 21--22, and~\ref{ch:ext-associahedra}.  The Stirling triangles
  are treated as inverse changes of factorial basis; Catalan numbers are then
  obtained by reversion and refined by Narayana and polygon-dissection counts.\\[1ex]

Gamma hierarchy
& \emph{Digamma Function}; \emph{Polygamma Function}; and \emph{Barnes
  G-Function}.
& Chapter~\ref{ch:ext-gamma-barnes}.  Weierstrass products and logarithmic
  differentiation yield the gamma hierarchy; Alexeiewsky's integral and
  Euler--Maclaurin summation yield the Barnes functional equation, integer
  values, and asymptotic expansion.\\[1ex]

Summation and asymptotic transfer
& \emph{Euler--Maclaurin Integration Formulas} and \emph{Darboux's Formula}.
& Chapter~\ref{ch:ext-em-darboux}.  Periodic Bernoulli functions give an exact
  remainder, and local singular expansions are transferred to coefficients
  with explicit error control; harmonic and Catalan expansions serve as worked
  tests.\\
\end{longtable}
\renewcommand{\arraystretch}{1}


\section{Formula-status ledger}
The displayed identities audited across the five drafts and supporting topic files
fall into four editorial classes.
\begin{enumerate}
\item \emph{Exact identities} are reproduced, usually in a notation that prevents
      collisions, and proved by a bijection, coefficient extraction, finite
      differences, matrix inversion, or a generating-function calculation.
\item \emph{Equivalent normal forms} are consolidated.  For example, the
      multiplicity sum, Bell-polynomial identity, and set-partition form of Faà di
      Bruno's formula are one theorem viewed in three coordinate systems.
\item \emph{Corrected statements} retain the source's intended content but repair
      an index shift, a missing range or analyticity hypothesis, or a suppressed
      Jacobian.  The type-B Eulerian and delta-function reversion formulas are the
      main examples.
\item \emph{Schematic or computational statements} are replaced by something
      checkable.  Undefined asymptotic placeholders become recursively defined
      saddle coefficients; finite tables and modular searches come with exact
      recurrences or certificate procedures.
\end{enumerate}
The resulting distinction is important: a theorem is not silently weakened into
numerical evidence, and a numerical observation is not silently promoted to a
theorem.

\section{Implementation-neutral algorithms}
The transient software lists in encyclopedic source material are best replaced by
algorithms whose correctness follows from the proved recurrences.  To compute all
entries through row $N$, initialize the boundary values and fill either Stirling
triangle by its two-term recurrence; this uses $O(N^2)$ arithmetic operations and
$O(N)$ memory if only the preceding row is retained.  Bell numbers are then row
sums of the second-kind triangle, or may be produced by the binomial recurrence.
Eulerian numbers admit the same triangular complexity.  Partial Bell polynomials
can be filled by the pointing recurrence, with exact polynomial or symbolic-series
arithmetic as appropriate.  For inverse series, Lagrange--Bürmann gives each target
coefficient independently from a truncated power of the defining series; Newton
iteration is asymptotically faster for very high orders, while the Lagrange form is
usually more transparent for symbolic coefficients.  Every one of these algorithms
is valid over any coefficient ring in which the displayed divisions make sense.

\section{Dependency map}
The logical flow of the monograph may be summarized as
\[
\begin{gathered}
 \text{factorial bases and finite differences}
 \longrightarrow \text{Stirling transforms}\\
 \longrightarrow \text{Bell partition calculus}
 \longrightarrow \text{composition and differentiation}.
\end{gathered}
\]
with Eulerian numbers providing the parallel change of coordinates between powers
and shifted binomial coefficients.  Formal residues then turn implicit substitution
into Lagrange inversion.  This dependency map explains why identities that appear
unrelated in the separate source files repeatedly reduce to the same three
operations: change of basis, labelled decomposition, and coefficient extraction.

\backmatter

\chapter*{Selected bibliography and provenance}
\addcontentsline{toc}{chapter}{Selected bibliography and provenance}

\begin{thebibliography}{99}

\bibitem{Bell1934}
E.~T. Bell,
``Exponential polynomials,''
\emph{Annals of Mathematics}, second series, vol.~35 (1934), pp.~258--277.

\bibitem{Broder1984}
A.~Z. Broder,
``The $r$-Stirling numbers,''
\emph{Discrete Mathematics}, vol.~49 (1984), pp.~241--259.

\bibitem{Charalambides2002}
C.~A. Charalambides,
\emph{Enumerative Combinatorics},
Chapman \& Hall/CRC, 2002.

\bibitem{Claesson2001}
A. Claesson,
``Generalized pattern avoidance,''
\emph{European Journal of Combinatorics}, vol.~22 (2001), pp.~961--971.

\bibitem{Comtet1974}
L. Comtet,
\emph{Advanced Combinatorics: The Art of Finite and Infinite Expansions},
translated by J.~W. Nienhuys, D. Reidel, 1974.

\bibitem{FlajoletSedgewick2009}
P. Flajolet and R. Sedgewick,
\emph{Analytic Combinatorics},
Cambridge University Press, 2009.

\bibitem{GKP1994}
R.~L. Graham, D.~E. Knuth, and O. Patashnik,
\emph{Concrete Mathematics}, second edition,
Addison--Wesley, 1994.

\bibitem{Harper1967}
L.~H. Harper,
``Stirling behavior is asymptotically normal,''
\emph{Annals of Mathematical Statistics}, vol.~38 (1967), pp.~410--414.

\bibitem{Henrici1974}
P. Henrici,
\emph{Applied and Computational Complex Analysis}, vol.~1,
Wiley, 1974.

\bibitem{HsuShiue1998}
L.~C. Hsu and P.~J.-S. Shiue,
``A unified approach to generalized Stirling numbers,''
\emph{Advances in Applied Mathematics}, vol.~20 (1998), pp.~366--384.

\bibitem{Johnson2002}
W.~P. Johnson,
``The curious history of Faà di Bruno's formula,''
\emph{American Mathematical Monthly}, vol.~109 (2002), pp.~217--234.

\bibitem{MoserWyman1955}
L. Moser and M. Wyman,
``An asymptotic formula for the Bell numbers,''
\emph{Transactions of the Royal Society of Canada}, section III, vol.~49
(1955), pp.~49--54.

\bibitem{NISTDLMF}
F.~W.~J. Olver, A.~B. Olde Daalhuis, D.~W. Lozier, B.~I. Schneider,
R.~F. Boisvert, C.~W. Clark, B.~R. Miller, B.~V. Saunders,
H.~S. Cohl, and M.~A. McClain, editors,
\emph{NIST Digital Library of Mathematical Functions},
National Institute of Standards and Technology.

\bibitem{Riordan1958}
J. Riordan,
\emph{An Introduction to Combinatorial Analysis},
Wiley, 1958.

\bibitem{Roman1984}
S. Roman,
\emph{The Umbral Calculus},
Academic Press, 1984.

\bibitem{Spivey2008}
M.~Z. Spivey,
``A generalized recurrence for Bell numbers,''
\emph{Journal of Integer Sequences}, vol.~11 (2008), Article 08.2.5.

\bibitem{Stanley2012}
R.~P. Stanley,
\emph{Enumerative Combinatorics}, vol.~1, second edition,
Cambridge University Press, 2012.

\bibitem{Temme1993}
N.~M. Temme,
``Asymptotic estimates of Stirling numbers,''
\emph{Studies in Applied Mathematics}, vol.~89 (1993), pp.~233--243.

\bibitem{Wilf2005}
H.~S. Wilf,
\emph{generatingfunctionology}, third edition,
A K Peters, 2005.

\bibitem{Good1965}
I.~J. Good,
``The generalization of Lagrange's expansion and the enumeration of trees,''
\emph{Proceedings of the Cambridge Philosophical Society}, vol.~61 (1965),
pp.~499--517.

\bibitem{Gessel1987}
I.~M. Gessel,
``A combinatorial proof of the multivariable Lagrange inversion formula,''
\emph{Journal of Combinatorial Theory, Series A}, vol.~45 (1987), pp.~178--195.

\bibitem{Loday2004}
J.-L. Loday,
``Realization of the Stasheff polytope,''
\emph{Archiv der Mathematik}, vol.~83 (2004), pp.~267--278.

\bibitem{Norlund1924}
N.~E. N\o rlund,
\emph{Vorlesungen \"uber Differenzenrechnung},
Springer, Berlin, 1924.

\bibitem{Stasheff1963}
J.~D. Stasheff,
``Homotopy associativity of H-spaces I and II,''
\emph{Transactions of the American Mathematical Society}, vol.~108 (1963),
pp.~275--312.

\bibitem{WhittakerWatson1927}
E.~T. Whittaker and G.~N. Watson,
\emph{A Course of Modern Analysis}, fourth edition,
Cambridge University Press, 1927.

\bibitem{WikipediaDossier}
The \nolinkurl{Wikipedia/} directory in \emph{Formulae-2.zip}, containing
reference snapshots on Bell numbers and polynomials, Eulerian numbers,
Fa\`a di Bruno's formula, Lagrange inversion and reversion, and Stirling
numbers of both kinds.

\bibitem{MathWorldDossier}
The \nolinkurl{MathWorld/} directory in \emph{Formulae-2.zip}, containing
MathWorld PDF and notebook materials on factorial polynomial families,
special functions, inversion, asymptotic methods, and associahedra.

\bibitem{SourceArchive}
\emph{Formulae-2.zip}, the source package supplied for this project: five
substantial LaTeX drafts in \nolinkurl{articles/}, reference snapshots in
\nolinkurl{Wikipedia/}, and MathWorld PDF/notebook materials in
\nolinkurl{MathWorld/}.

\end{thebibliography}

\printindex
\end{document}
